%File: anonymous-submission-latex-2026.tex
\documentclass[letterpaper]{article} % DO NOT CHANGE THIS
\usepackage[submission]{aaai2026}  % DO NOT CHANGE THIS
\usepackage{times}  % DO NOT CHANGE THIS
\usepackage{helvet}  % DO NOT CHANGE THIS
\usepackage{courier}  % DO NOT CHANGE THIS
\usepackage[hyphens]{url}  % DO NOT CHANGE THIS
\usepackage{graphicx} % DO NOT CHANGE THIS
\urlstyle{rm} % DO NOT CHANGE THIS
\def\UrlFont{\rm}  % DO NOT CHANGE THIS
\usepackage{natbib}  % DO NOT CHANGE THIS AND DO NOT ADD ANY OPTIONS TO IT
\usepackage{caption} % DO NOT CHANGE THIS AND DO NOT ADD ANY OPTIONS TO IT
\frenchspacing  % DO NOT CHANGE THIS
\setlength{\pdfpagewidth}{8.5in} % DO NOT CHANGE THIS
\setlength{\pdfpageheight}{11in} % DO NOT CHANGE THIS
%
% These are recommended to typeset algorithms but not required. See the subsection on algorithms. Remove them if you don't have algorithms in your paper.
\usepackage{algorithm}
\usepackage{algorithmic}

\usepackage{amsmath} 
\usepackage{amsthm}
\newtheorem{definition}{Definition}  
\newtheorem{theorem}{Theorem}
\newtheorem{lemma}{Lemma}
\usepackage{booktabs}
\usepackage{multirow}

\usepackage{makecell}

%
% These are are recommended to typeset listings but not required. See the subsubsection on listing. Remove this block if you don't have listings in your paper.
\usepackage{newfloat}
\usepackage{listings}
\DeclareCaptionStyle{ruled}{labelfont=normalfont,labelsep=colon,strut=off} % DO NOT CHANGE THIS
\lstset{%
	basicstyle={\footnotesize\ttfamily},% footnotesize acceptable for monospace
	numbers=left,numberstyle=\footnotesize,xleftmargin=2em,% show line numbers, remove this entire line if you don't want the numbers.
	aboveskip=0pt,belowskip=0pt,%
	showstringspaces=false,tabsize=2,breaklines=true}
\floatstyle{ruled}
\newfloat{listing}{tb}{lst}{}
\floatname{listing}{Listing}
%
% Keep the \pdfinfo as shown here. There's no need
% for you to add the /Title and /Author tags.
\pdfinfo{
/TemplateVersion (2026.1)
}

\usepackage{amsfonts}
% \usepackage[pagebackref=true,breaklinks=true,letterpaper=true,colorlinks,bookmarks=false]{hyperref}

% \newcommand{\LYY}[1]{{\color{magenta}{\bf LYY:} #1}}


% DISALLOWED PACKAGES
% \usepackage{authblk} -- This package is specifically forbidden
% \usepackage{balance} -- This package is specifically forbidden
% \usepackage{color (if used in text)
% \usepackage{CJK} -- This package is specifically forbidden
% \usepackage{float} -- This package is specifically forbidden
% \usepackage{flushend} -- This package is specifically forbidden
% \usepackage{fontenc} -- This package is specifically forbidden
% \usepackage{fullpage} -- This package is specifically forbidden
% \usepackage{geometry} -- This package is specifically forbidden
% \usepackage{grffile} -- This package is specifically forbidden
% \usepackage{hyperref} -- This package is specifically forbidden
% \usepackage{navigator} -- This package is specifically forbidden
% (or any other package that embeds links such as navigator or hyperref)
% \indentfirst} -- This package is specifically forbidden
% \layout} -- This package is specifically forbidden
% \multicol} -- This package is specifically forbidden
% \nameref} -- This package is specifically forbidden
% \usepackage{savetrees} -- This package is specifically forbidden
% \usepackage{setspace} -- This package is specifically forbidden
% \usepackage{stfloats} -- This package is specifically forbidden
% \usepackage{tabu} -- This package is specifically forbidden
% \usepackage{titlesec} -- This package is specifically forbidden
% \usepackage{tocbibind} -- This package is specifically forbidden
% \usepackage{ulem} -- This package is specifically forbidden
% \usepackage{wrapfig} -- This package is specifically forbidden
% DISALLOWED COMMANDS
% \nocopyright -- Your paper will not be published if you use this command
% \addtolength -- This command may not be used
% \balance -- This command may not be used
% \baselinestretch -- Your paper will not be published if you use this command
% \clearpage -- No page breaks of any kind may be used for the final version of your paper
% \columnsep -- This command may not be used
% \newpage -- No page breaks of any kind may be used for the final version of your paper
% \pagebreak -- No page breaks of any kind may be used for the final version of your paperr
% \pagestyle -- This command may not be used
% \tiny -- This is not an acceptable font size.
% \vspace{- -- No negative value may be used in proximity of a caption, figure, table, section, section, subsection, or reference
% \vskip{- -- No negative value may be used to alter spacing above or below a caption, figure, table, section, section, subsection, or reference

\setcounter{secnumdepth}{0} %May be changed to 1 or 2 if section numbers are desired.

% The file aaai2026.sty is the style file for AAAI Press
% proceedings, working notes, and technical reports.
%

% Title

% Your title must be in mixed case, not sentence case.
% That means all verbs (including short verbs like be, is, using,and go),
% nouns, adverbs, adjectives should be capitalized, including both words in hyphenated terms, while
% articles, conjunctions, and prepositions are lower case unless they
% directly follow a colon or long dash
\title{Flatness Matters for LoRA-Based Continual Learning: A PAC-Bayesian Perspective}

% LYY: Harnessing Flat Minima for Continual LoRA: A Theoretically-Grounded Framework


\author{
    %Authors
    % All authors must be in the same font size and format.
    Written by AAAI Press Staff\textsuperscript{\rm 1}\thanks{With help from the AAAI Publications Committee.}\\
    AAAI Style Contributions by Pater Patel Schneider,
    Sunil Issar,\\
    J. Scott Penberthy,
    George Ferguson,
    Hans Guesgen,
    Francisco Cruz\equalcontrib,
    Marc Pujol-Gonzalez\equalcontrib
}
\affiliations{
    %Afiliations
    \textsuperscript{\rm 1}Association for the Advancement of Artificial Intelligence\\
    % If you have multiple authors and multiple affiliations
    % use superscripts in text and roman font to identify them.
    % For example,

    % Sunil Issar\textsuperscript{\rm 2},
    % J. Scott Penberthy\textsuperscript{\rm 3},
    % George Ferguson\textsuperscript{\rm 4},
    % Hans Guesgen\textsuperscript{\rm 5}
    % Note that the comma should be placed after the superscript

    1101 Pennsylvania Ave, NW Suite 300\\
    Washington, DC 20004 USA\\
    % email address must be in roman text type, not monospace or sans serif
    proceedings-questions@aaai.org
%
% See more examples next
}

%Example, Single Author, ->> remove \iffalse,\fi and place them surrounding AAAI title to use it
\iffalse
\title{My Publication Title --- Single Author}
\author {
    Author Name
}
\affiliations{
    Affiliation\\
    Affiliation Line 2\\
    name@example.com
}
\fi

\iffalse
%Example, Multiple Authors, ->> remove \iffalse,\fi and place them surrounding AAAI title to use it
\title{My Publication Title --- Multiple Authors}
\author {
    % Authors
    First Author Name\textsuperscript{\rm 1},
    Second Author Name\textsuperscript{\rm 2},
    Third Author Name\textsuperscript{\rm 1}
}
\affiliations {
    % Affiliations
    \textsuperscript{\rm 1}Affiliation 1\\
    \textsuperscript{\rm 2}Affiliation 2\\
    firstAuthor@affiliation1.com, secondAuthor@affilation2.com, thirdAuthor@affiliation1.com
}
\fi


% REMOVE THIS: bibentry
% This is only needed to show inline citations in the guidelines document. You should not need it and can safely delete it.
% \usepackage{bibentry}
% END REMOVE bibentry

\begin{document}

\maketitle

\begin{abstract}
Continual learning with parameter-efficient adaptation faces fundamental challenges in balancing knowledge retention and generalization across sequential tasks. This work investigates the role of loss landscape flatness in shaping generalization within the Low-Rank Adaptation (LoRA) framework for continual learning. We revisit PAC-Bayesian generalization theory and establish new bounds that explicitly incorporate flatness metrics into continual adaptation with LoRA. Theoretical analysis reveals how flat minima contribute to mitigating forgetting and improving robustness through structured posterior distributions over task sequences.  Based on the proposed theoretical framework, we conduct experiments to assess the effects of flatness-aware perturbations, including Sharpness-Aware Minimization (SAM), Gaussian noise, and Fisher-guided noise, under both local (LoRA-only) and global (full-model) scopes. Experimental results across diverse continual learning benchmarks demonstrate that flatness in parameter-efficient adaptation plays a critical role in enhancing generalization and stability, offering new insights into the interplay between optimization geometry and continual learning dynamics.



% Our analysis reveals how flat minima contribute to mitigating forgetting and improving robustness through structured posterior distributions over task sequences. Building on this theoretical foundation, we conduct experiments to examine the effects of flatness-aware perturbations—such as SAM, Gaussian, and Fisher-guided noise—under both local (LoRA-only) and global (full model) scopes. Results across diverse continual learning benchmarks demonstrate that flatness within parameter-efficient adaptation plays a critical role in enhancing generalization and stability, offering new insights into the interaction between optimization geometry and continual learning dynamics. 


% We derive task-dependent bounds that highlight the role of expected sharpness in shaping robustness and mitigating forgetting. 

 


% We systematically examine the impact of perturbation type (SAM, Gaussian, Fisher-guided) and perturbation scope (LoRA-only vs. full-model) on two representative adaptation strategies: SeqLoRA and IncLoRA. Our results reveal that promoting flatness consistently improves forward transfer and reduces forgetting.






% Building on this theoretical foundation, we perform comprehensive experiments across multiple continual learning benchmarks to assess the effectiveness of flatness-promoting strategies. Specifically, we investigate Sharpness-Aware Minimization (SAM), Gaussian noise, and Fisher-guided perturbations under both local (LoRA-only) and global (full-model) scopes. 



% Extensive experiments across diverse benchmarks validate our analysis, showing that promoting flatness—via SAM, Gaussian, or Fisher-guided perturbations—consistently enhances generalization and stability in LoRA-based continual learning. Our findings provide both theoretical and empirical insights into the interplay between optimization geometry and parameter-efficient adaptation.




% This work investigates the role of loss landscape flatness in shaping generalization within the Low-Rank Adaptation (LoRA) framework for continual learning. We revisit PAC-Bayesian generalization theory and establish new bounds that explicitly incorporate flatness metrics into continual adaptation with LoRA. Our analysis reveals how flat minima contribute to mitigating forgetting and improving robustness through structured posterior distributions over task sequences. Building on this theoretical foundation, we conduct systematic experiments to examine the effects of flatness-aware perturbations—such as SAM, Gaussian, and Fisher-guided noise—under both local (LoRA-only) and global (full model) scopes. Results across diverse continual learning benchmarks demonstrate that flatness within parameter-efficient adaptation plays a critical role in enhancing generalization and stability, offering new insights into the interaction between optimization geometry and continual learning dynamics.

% Continual learning with large-scale models requires balancing efficient adaptation, knowledge retention, and robust generalization across tasks. In this work, we study Incremental Low-Rank Adaptation (IncLoRA) for continual learning, with a focus on how the flatness of the loss landscape affects model generalization. Using a PAC-Bayesian framework, we derive a new generalization bound for IncLoRA, revealing a theoretical link between loss flatness and the way knowledge is structured across tasks. Based on this analysis, we introduce Fisher Information guided Flat-Minima Orthogonal Incremental LoRA (FFOI), a method that combines global flatness regularization, orthogonality constraints between task-specific adapters, and knowledge alignment guided by Fisher information. Extensive experiments on multiple continual learning benchmarks confirm the effectiveness of FFOI and provide empirical support for the theoretical findings on flatness, orthogonality, and knowledge retention.
% Motivated by these challenges, we first establish a novel PAC-Bayesian generalization bound for Multi-LoRA in continual learning, which highlights the joint roles of flat loss landscape , and independence of different LoRa parameters and transfer of knowledge between different tasks. Building on this theoretical insight, we propose Fisher Information guided Flat-Minima Orthogonal Incremental LoRA (FFOI), a novel method that unifies sharpness-aware empirical risk minimization with Fisher information-based flatness regularization and orthogonality constraints among LoRA adapters.
\end{abstract}



\section{Introduction}



The widespread deployment of large language models (LLMs) in real-world applications has created a growing demand for efficient and flexible adaptation to diverse and continuously emerging downstream tasks. In real‐world settings, tasks emerge incrementally, making continual learning the natural framework for model updates. However, fully fine‐tuning an LLM for each new task is often impractical due to enormous computational and storage costs. Parameter‐efficient techniques, such as Low‐Rank Adaptation (LoRA), address this challenge by introducing small, task‐specific modules while freezing the bulk of the model’s weights. However, even within this efficient multi LoRA paradigm, the balance between mitigating forgetting and ensuring generalization to future tasks still remains an open and challenging problem.


Although LoRA dramatically reduces overhead, it still faces the core difficulty of balancing knowledge retention (avoiding catastrophic forgetting) with the ability to generalize to future tasks.



However, the computational and storage demands of full-model retraining for each task are often prohibitive. Parameter-efficient fine-tuning methods, such as Low-Rank Adaptation (LoRA), have emerged as practical solutions, allowing models to accommodate each new task by learning compact, task-specific modules while keeping the majority of model parameters fixed. However, even within this efficient multi LoRA paradigm, the balance between mitigating forgetting and ensuring generalization to future tasks still remains an open and challenging problem.


The connection between flat minima and generalization has long been a central theme in deep learning research. Both early and recent studies have shown that models converging to flat regions of the loss surface are less sensitive to parameter perturbations, thereby exhibiting enhanced robustness and generalization~\shortcite{hinton1993keeping,NIPS1994_01882513,FlatMinima1997,chaudhariEntropySGDBiasingGradient2019,bahri-etal-2022-sharpness,zhouUnderstandingConvergenceGeneralization2024,zhangExploringFlatMinima2024a,yangRevisitingFlatnessAwareOptimization2025}. Building on these insights, methods such as Stochastic Weight Averaging (SWA)~\shortcite{izmailovAveragingWeightsLeads2019} and Sharpness-Aware Minimization (SAM)~\shortcite{foretSharpnessAwareMinimizationEfficiently2021} have explicitly promote flatness, achieving notable gains—especially in computer vision tasks. Moreover, emerging evidence suggests that flat minima may also play a vital role in alleviating catastrophic forgetting~\cite{mirzadehUnderstandingRoleTraining2020,JMLRAnEmpiricalInvestigatio}, a core challenge in continual learning scenarios. However, the effects and mechanisms of flat minima remain underexplored in the context of parameter-efficient fine-tuning, particularly within LoRA systems for large language models. In particular, there is a critical lack of systematic investigation into how flat minima influence both generalization and forgetting when adapting LLMs to a sequence of downstream tasks through continual LoRA adaptation.

\begin{figure}[t]
  \centering
  \includegraphics[width=1\linewidth]{images/show_Loss_different-caijian.pdf}
  \caption{Illustration of the three types of risk in hierarchical PAC-Bayesian continual learning: future task generalization risk $L(Q)$, expected generalization risk on observed tasks $L(Q_{1:n})$, and empirical risk on observed tasks $\hat{L}(Q_{1:n})$.}
  \label{fig concept: 3 different definaiton}
\end{figure}

The main contributions of the paper can be summarized as:
\begin{itemize}
    \item 
    \item 
    \item 
\end{itemize}



% \LYY{The story should be: The core contradiction of continual learning lies in plasticity and stability. Through a new theoretical analysis (PAC-Bayes bound), we mathematically revealed for the first time that in the era of large models, in the continual LoRA scenario, solving this contradiction requires optimizing three orthogonal dimensions at the same time: empirical risk, solution flatness, and cross-task knowledge structure. Existing methods only partially solve one or two dimensions, so the effect is limited. Our FFOI framework is the first solution designed to fully optimize this theoretical bound. Its three components precisely correspond to the three requirements of the theory, so it can achieve performance breakthroughs.}

% The main contributions of this paper are: 
% \begin{enumerate}
%     \item \LYY{We provide the first theoretical characterization of the generalization problem of incremental LoRA in continual learning. Our PAC-Bayes bound reveals the intrinsic connection between flatness, orthogonality, and knowledge transfer in overcoming forgetting.}
%     \item \LYY{We design a theoretically-grounded FFOI framework that synergizes three components - flatness optimization, orthogonality constraints, and Fisher-guided alignment.}
%     \item \LYY{We conduct comprehensive experimental validation on multiple challenging benchmarks, which not only demonstrate the superior performance of FFOI but also provide direct support for the intrinsic mechanisms predicted by the theoretical framework.}
% \end{enumerate}

\section{Realted Work}

\subsection{LoRA-Based Continual Learning}

A critical challenge in adapting large pre-trained models to downstream continual learning is to leverage prior knowledge for new tasks while maintaining generalizability to future tasks. Parameter-efficient strategies—including adapter\cite{pmlr-v97-houlsby19a}, prompts\cite{Wang_2022_CVPR}, and low-rank adaptation\shortcite{hu2022lora}—tackle this challenge by supporting incremental learning with minimal parameter updates, thereby reducing both resource consumption and the risk of catastrophic forgetting. Among these, Low-Rank Adaptation (LoRA)\cite{hu2022lora} has become particularly prominent, as it updates only a small subset of parameters while keeping the backbone fixed. LoRA-based approaches, such as SDLoRA\cite{wu2025sdlorascalabledecoupledlowrank}, have demonstrated strong effectiveness in mitigating forgetting and reducing resource usage. 
However, LoRA-based continual learning faces persistent challenges in handling diverse and long task sequences. In partially shared parameter settings, the complexity of task similarity and interference highlights the stability-plasticity trade-off. Recent studies have introduced orthogonalization strategies, such as OLoRA~\cite{NEURIPS2020_70d85f35,wang2023orthogonalsubspacelearninglanguage}, to constrain parameter updates and reduce interference across tasks. Although effective in preserving task-specific representations, these methods may hinder positive transfer when tasks are correlated. Balancing interference reduction and knowledge sharing remains a fundamental problem in LoRA-based continual learning.

% However, LoRA-based continual learning still faces challenges when dealing with diverse or long task sequences. In settings where parameters are not fully shared, task similarity and interference become more complex, further highlighting the stability-plasticity trade-off. To address this, recent work has explored orthogonalization strategies, such as OLoRA\cite{NEURIPS2020_70d85f35,wang2023orthogonalsubspacelearninglanguage}, which constrain parameter updates to reduce interference across tasks. While these methods can alleviate forgetting by isolating task-specific knowledge, they may also limit positive transfer when tasks are related. Thus, balancing interference reduction and knowledge alignment remains a key challenge, motivating further principled advances in LoRA-based continual learning.



% A critical challenge in adapting large pre-trained models to downstream continual learning is to leverage prior knowledge for new tasks while maintaining generalizability to future tasks.







% Continual adaptation from large pre-trained models to downstream tasks has emerged as a central scenario in modern machine learning, where models must incrementally acquire new knowledge under strict efficiency constraints while preserving previously learned capabilities. Parameter-efficient adaptation strategies—including adapter modules, prompt-based methods, and low-rank adaptation—have been widely adopted to facilitate incremental learning without full model retraining. Adapter-based techniques, such as ADA, employ lightweight modules in conjunction with knowledge distillation to efficiently adapt pre-trained transformers, while prompt-based approaches (e.g., L2P) dynamically select task-specific prompts to guide the model’s behavior. Low-Rank Adaptation (LoRA) has emerged as a prominent paradigm, updating only the low-rank matrices while keeping the backbone frozen. These LoRA-based methods, exemplified by SDLoRA and related variants, have shown significant effectiveness in reducing catastrophic forgetting and minimizing resource consumption, establishing a strong foundation for continual learning in large-scale pre-trained models.


% Despite their merits, LoRA-based continual learning methods encounter inherent limitations when faced with heterogeneous or lengthy task sequences. In scenarios where model parameters are not globally shared—such as multi-branch or multi-head architectures—task similarity introduces additional complexity, intensifying the classical stability-plasticity trade-off. To address these challenges, recent studies have proposed orthogonalization-based approaches, including Orthogonal Gradient Descent (OGD) and Orthogonal LoRA (OLoRA), which constrain parameter updates to be orthogonal to those associated with previous tasks. While such strategies can effectively mitigate catastrophic forgetting by isolating task-specific knowledge, they may inadvertently restrict beneficial transfer between related tasks. Achieving an optimal balance between interference reduction and knowledge alignment thus remains an open problem, motivating the pursuit of more principled and theoretically grounded solutions within the LoRA-based continual learning framework.



% However, despite their efficiency and effectiveness, LoRA-based continual learning remains fundamentally challenged by task interference and the stability-plasticity dilemma—particularly when tasks are diverse or the adaptation spans long sequences


\subsection{Flat Minima in Continual Learning}
It has been found that flat local minima leads to better generalization capabilities than sharp minima in the sense that a flat minimizer is more robust when the test loss is shifted due to random perturbations\shortcite{NEURIPS2020_518a38cc,chen2022visiontransformersoutperformresnets,bahri2022sharpnessawareminimizationimproveslanguage}. To promote flat minima in continual learning, three mainstream strategies have been explored. The first focuses on directly optimizing flatness metrics\shortcite{foretSharpnessAwareMinimizationEfficiently2021,zhangGradientNormAware2023,liRevisitingRandomWeight2024}, with Sharpness-Aware Minimization (SAM)~\cite{foretSharpnessAwareMinimizationEfficiently2021} as a representative approach. The second approach leverages mode connectivity, constructing interpolating between multiple solutions to traverse flat paths in parameter space\shortcite{NEURIPS2018_be3087e7,izmailov2019averagingweightsleadswider,mirzadeh2020linearmodeconnectivitymultitask}. 
The third comes down to obtaining well distributed representations, such as by using pre-training\shortcite{JMLR:v24:22-0496,ramasesh2022effect}, wider network architectures\shortcite{pmlr-v162-mirzadeh22a,ramasesh2020anatomycatastrophicforgettinghidden}  and self-supervised learning\shortcite{hu2022doesselfsupervisedpretrainingperform,madaan2022representationalcontinuityunsupervisedcontinual}.
However, the relationship between loss landscape geometry and generalization in LoRA-based continual learning remains insufficiently characterized. In particular, the interaction between flatness-aware optimization and parameter-efficient adaptation lacks systematic analysis from both empirical and theoretical perspectives, limiting current understanding of generalization and forgetting in such models.


% The relationship between loss landscape geometry and generalization in LoRA-based continual learning remains an important and underexplored topic, motivating further investigation into flatness-aware optimization within parameter-efficient frameworks.


\subsection{PAC-Bayesian Theory for Generalization in Continual Learning}

The PAC-Bayesian framework provides a principled approach for quantifying generalization by establishing bounds on the gap between performance on the training data and the underlying data distribution~\cite{mcallesterPACBayesianStochasticModel2003}. This framework has been extended to continual and lifelong learning, where generalization depends not only on empirical risk within individual tasks but also on the consistency of model updates across tasks \cite{pentinaPACBayesianBoundLifelong2014}. Wang et al.~\cite{wangCoSCLCooperationSmall2022} further investigate this trade-off from the perspective of VC dimension, highlighting that flatter solutions in the loss landscape contribute to tighter generalization bounds. However, existing PAC-Bayesian analyses rarely incorporate explicit measures of loss landscape geometry, despite growing empirical evidence that flatter minima improve generalization and mitigate forgetting.  This gap is particularly relevant for parameter-efficient continual learning, where flatness-aware optimization techniques are increasingly explored. However, their impact on PAC-Bayesian generalization guarantees remains underexplored. Incorporating loss landscape geometry into PAC-Bayesian analysis could offer a clearer theoretical understanding of generalization and forgetting in LoRA-based continual learning.










% The PAC-Bayesian framework provides a principled foundation for analyzing generalization in stochastic learning by establishing generalization bounds in terms of both empirical risk and the divergence between posterior and prior distributions over model parameters~\cite{mcallesterPACBayesianStochasticModel2003}.  This framework has been extended to continual and lifelong learning, where generalization depends not only on the empirical risk within individual tasks but also on the consistency of model updates across tasks \cite{pentinaPACBayesLifelongLearning2014}. Wang et al.\cite{wangCoSCLCooperationSmall2022} further explore this trade-off from the perspective of VC dimension, emphasizing that flatter solutions in the loss landscape lead to tighter generalization bounds. Despite these advances, existing PAC-Bayesian analyses rarely incorporate explicit measures of loss landscape geometry into their bounds,despite increasing empirical evidence that flatter minima lead to better generalization and forgetting resistance. This gap is particularly relevant given the growing adoption of flatness-aware optimization techniques (e.g., SAM) in practice. Bridging this gap through a systematic integration of loss landscape geometry into PAC-Bayesian generalization analysis constitutes an important step toward a deeper theoretical understanding of generalization and forgetting in LoRA-based continual learning.












% The PAC-Bayesian theory expresses probably approximately correct (PAC) bounds on the generalization risk \shortcite{mcallesterPACBayesianStochasticModel2003}.  It is a a strong link between frequentist PAC-Bayesian risk bounds and the Bayesian marginal likelihood. Pentina extend  it to the A PAC-Bayesian Bound for Lifelong Learning, provide a unified view on  forgetting-generalization tradeoff.   Liyuan Wang theoretically analyze the generalization errors for learning plasticity and memory stability in continual learning\cite{wangCoSCLCooperationSmall2022} from VC dimension ,point out that flatness of loss landscape is a key factor for upper-bounded. However, existing PAC-Bayesian analyses rarely consider the influence of loss landscape flatness, despite increasing empirical evidence that flatter minima lead to better generalization and forgetting resistance. Bridging this gap remains an important direction for both theory and practice in parameter-efficient continual learning. 



% The PAC-Bayesian framework offers a principled approach for analyzing generalization in continual learning, capturing how model updates and task distribution shifts affect generalization error across task sequences. Recent studies highlight that forgetting and performance degradation are shaped not only by empirical loss, but also by the discrepancy between learned solutions for different tasks. However, existing PAC-Bayesian analyses rarely consider the influence of loss landscape flatness, despite increasing empirical evidence that flatter minima lead to better generalization and forgetting resistance. Bridging this gap remains an important direction for both theory and practice in parameter-efficient continual learning.

\section{Notation and Preliminaries}

\subsection{Continual Learning Problem Setting }

We formalize continual learning (CL) with the concept of the task environment framework~\cite{pentinaPACBayesianBoundLifelong2014}. Let $\mathcal{X}$ and $\mathcal{Y}$ denote the input and output spaces, respectively, and let $\mathcal{F} \subset \{f : \mathcal{X} \rightarrow \mathcal{Y}\}$ denote the hypothesis space. The loss function is $\ell : \mathcal{F} \times \mathcal{X} \times \mathcal{Y} \rightarrow \mathbb{R}_{\geq 0}$.

We assume an underlying environment with a task distribution $\mathcal{T}$, such that each task $t \in \{1, 2, \ldots, n\}$ is independently sampled from $\mathcal{T}$. For each task $t$, the agent receives a training set $S_t = \{(x_{tj}, y_{tj})\}_{j=1}^{m_t}$ and a test set $T_t$, both drawn i.i.d.\ from a task-specific data distribution $\mathcal{D}_t$. The empirical loss for function $f \in \mathcal{F}$ on task $t$ is defined as
\begin{equation}
    \hat{L}_{S_t}(f) = \frac{1}{m_t} \sum_{j=1}^{m_t} \ell(f(x_{tj}), y_{tj}).
\end{equation}

\subsection{Parameter-Efficient CL: LoRA Adaptation}

Modern large-scale pretrained models are often adapted to downstream task using parameter-efficient methods. In the context of continual learning, we distinguish two principal strategies:
\begin{itemize}
    % \item \textbf{Full-parameter Finetuning (SeqFT):} All parameters are updated for each new task.
    \item \textbf{Sequence LoRA (SeqLoRA):} A single  pair of LoRA matrices $(A, B)$ is introduced and trained sequentially across multiple tasks, with backbone weights $W_0$ frozen throughout. After $N$ tasks:
    $$
        \text{SeqLoRA:} \quad W_0 + AB
    $$
    \item \textbf{Incremental LoRA (IncLoRA):} For each new task $t$, an additional LoRA pair $(A_t, B_t)$ is appended. During training on task $t$, only $(A_t, B_t)$ are updated, with all previous LoRA modules and $W_0$ kept fixed.  After $N$ tasks:
    \begin{equation*}
        \text{IncLoRA:} \quad W_0 + \sum_{t=1}^N A_t B_t
    \end{equation*}
    % This incremental adaptation scheme provides greater flexibility for continual learning and serves as the basis for all subsequent theoretical and algorithmic developments in this work.
\end{itemize}


\subsection{Flatness Metrics in Continual Learning}


Flatness of the loss landscape is widely recognized as a important indicator of generalization and robustness in deep learning~\citep{liVisualizingLossLandscape2018, foretSharpnessAwareMinimizationEfficiently2021}. Several quantitative metrics have been proposed, of which two are particularly relevant:

\begin{definition}[$\epsilon$-Sharpness]\shortcite{foretSharpnessAwareMinimizationEfficiently2021}
The adversarial sharpness of a solution $\hat{\theta}$ on dataset $S$ is defined as
\begin{equation}
    \epsilon\text{-Sh}(\hat{\theta}, \rho) :=\max_{\epsilon \in B(\hat\theta, \rho)} \big( \hat{L}_{S}(\hat{\theta} + \epsilon) - \hat{L}_{S}(\hat{\theta}) \big)
\end{equation}
where $B(\hat\theta, \rho)$ denotes a ball of radius $\rho$ centered at $\hat\theta$.
\end{definition}

This metric evaluates the worst-case loss increase under bounded adversarial perturbations and is useful for analyzing the sharpness of local minima. However, its computation typically requires inner maximization, making it less practical for large-scale continual learning scenarios. Instead, we adopt a task-aware expected sharpness metric, which measures the average-case sensitivity of the solution to random Gaussian perturbations:

\begin{definition}[Task-Aware Expected Sharpness]\shortcite{liRevisitingRandomWeight2024}
For task $i$, the expected sharpness of a solution $\theta$ is defined as
\begin{equation}\label{def: Task-Aware Expected Sharpness}
    \mathrm{E\text{-}Sh}_i(\theta, \Sigma_Q) := \mathbb{E}_{\epsilon \sim \mathcal{N}(0,\,\Sigma_Q)} \left[ \hat{L}_{\mathcal{D}_i}(\theta + \epsilon) - \hat{L}_{\mathcal{D}_i}(\theta) \right],
\end{equation}
where $\mathcal{N}(0,\,\Sigma_Q)$ is a Gaussian distribution, typically instantiated as $\mathcal{N}(0,\, \sigma^2 \mathbf{I})$.
\end{definition}

% Unlike $\epsilon${-Sh}, task-aware expected sharpness is straightforward to compute, which models uncertainty over solutions using distributions of the model. For this reason, both 

% Our theoretical analysis and practical training procedures consistently use the expected sharpness metric to quantify and promote flatness. In addition to task-aware expected sharpness, we also employ two complementary metrics to characterize flatness: the maximal eigenvalue of the Hessian matrix \shortcite{NEURIPS2020_518a38cc} and the mean diagonal Fisher information\shortcite{NEURIPS2022_c859b99b}. To further visualize and compare the flatness of solutions under different methods, we supplement these quantitative measures with loss landscape plots\shortcite{NEURIPS2018_a41b3bb3}.
 % and naturally aligns with the PAC-Bayesian framework

\subsection{PAC-Bayesian Generalization Theory}

The PAC-Bayesian framework provides a principled approach to quantifying generalization by considering a distribution over model parameters, rather than a single deterministic solution. The classical generalization bound of McAllester~\citep{mcallesterPACBayesianStochasticModel2003} states that, for any posterior $Q$ and prior $P$ over the hypothesis space, the expected test risk is upper-bounded by the empirical risk plus a complexity term involving the KL divergence between $Q$ and $P$:

%in stochastic model selection 
% Given a prior $P$ and a posterior $Q$ over the hypothesis space $\mathcal{F}$, the empirical risk under $Q$ is defined as $\hat{L}(Q) = \mathbb{E}_{f \sim Q}[\hat{L}_S(f)]$. The classical generalization bound of McAllester~\citep{mcallesterPACBayesianStochasticModel2003} is as follows:

\begin{theorem}[PAC-Bayes Generalization Bound]
For any loss $\ell': \mathcal{F} \times \mathcal{X} \times \mathcal{Y} \to [0,1]$, with probability at least $1-\delta$ over the choice of training set $S$:
\begin{equation}
    \mathbb{E}_{f \sim Q}[L_{\mathcal{D}}^{\ell'}(f)] \leq \mathbb{E}_{f \sim Q}[\hat{L}_S^{\ell'}(f)] + \sqrt{\frac{KL(Q\|P) + \log \frac{m}{\delta}}{2(m-1)}}
\end{equation}
where $m = |S|$ and $KL(Q\|P)$ is the Kullback–Leibler divergence.
\end{theorem}

% This result connects generalization directly to the quality of fit on the training set and the divergence from a chosen prior. However, a key limitation of this classical bound is that it does not explicitly account for the geometry of the loss landscape, which is known to play a critical role in robustness and generalization. Addressing this gap is central to our subsequent theoretical development, where we extend the PAC-Bayesian framework to incorporate flatness-aware terms in the context of continual learning and parameter-efficient adaptation.


\section{PAC-Bayesian Generalization Theory with Flatness for Continual Learning}
\subsection{PAC-Bayesian Bound with Flatness for Continual Learning}




To rigorously analyze continual learning with LoRA adaptation, we use a hierarchical Bayesian framework~\cite{pentinaPACBayesianBoundLifelong2014}: the prior $P$ over model parameters is treated as a random variable drawn from a hyperprior $\mathcal{P}$, which is recursively updated to a hyperposterior $\mathcal{Q}$ as tasks are observed, thus reflecting accumulated knowledge.


We distinguish three core risk notions in continual learning (see Fig.~\ref{fig concept: 3 different definaiton}):
\begin{itemize}
    \item \textbf{Future-Task Generalization Risk:} The expected performance on future, unseen tasks, reflecting the generalization ability of continual learning, but generally intractable in practice since both the distribution over future tasks and the underlying data distribution for each task are not accessible.
    \begin{equation}
        L(\mathcal{Q}) := \mathbb{E}_{t \sim \mathcal{T}}\mathbb{E}_{P \sim \mathcal{Q}}L_{\mathcal{D}_t}\big(Q_t(\mathcal{D}_t, P)\big),
    \end{equation}
    where $Q_t(\mathcal{D}_t, P )$ is the posterior obtained by training the learning algorithm with prior $P$ and data distribution $\mathcal{D}_t$.
    \item \textbf{Expected Generalization Risk on Observed Tasks:}  The expected test performance over all previously observed tasks, which serves as a practical proxy for future-task generalization risk, but still remains an idealized quantity as it relies on the unknown data distributions of each task.
    \begin{equation}
        L(\mathcal{Q}_{1:n}) := \frac{1}{n} \sum_{t=1}^n \mathbb{E}_{P_{t} \sim \mathcal{Q}_{t}} \mathbb{E}_{f \sim Q_t(P_t)} [ L_{\mathcal{D}_t}(f) ],
    \end{equation}
    where  $\mathcal{Q}_{1:n} = (\mathcal{Q}_1, \ldots, \mathcal{Q}_n)$ denotes the set of hyperposteriors specifying the distributions over task-specific priors for all observed tasks, $Q_t(P_t)$ being the posterior obtained by training on the $t$-th task with prior $P_t$.
    \item \textbf{Empirical Risk on Observed Tasks:} The average training loss across all observed tasks, which is directly computable from the available training data and forms the empirical basis for generalization analysis.
    \begin{equation}
        \hat{L}(\mathcal{Q}_{1:n}) := \frac{1}{n} \sum_{t=1}^n \mathbb{E}_{P_t \sim \mathcal{Q}_t(P_t)} \mathbb{E}_{f \sim Q_t}  [ L_{\mathcal{S}_t}(f) ].
    \end{equation}
\end{itemize}

% Our objective is to establish a generalization bound that quantitatively connects these quantities, bridge the gap between the empirical risk $\hat{L}(\mathcal{Q}_{1:n})$ (directly optimizeable) and the generalization risk of the future task $L(\mathcal{Q})$ (ultimate target), thereby allowing for a reliable generalization assessment prediction for future tasks from observed data.


% These risk definitions clarify the central challenge of continual learning: how to quantitatively relate the empirical risk computed on observed data to get better performance on all learned task and the expected generalization risk on future tasks . 

These risk definitions clarify the central challenge of continual learning: to quantitatively connect the empirical risk on observed tasks to both  reduced forgetting across all learned tasks and enhanced generalization to future, unseen tasks. Classical PAC-Bayesian analysis offers generalization guarantees~\shortcite{pentinaPACBayesianBoundLifelong2014} but does not explicitly address two critical aspects in continual learning: the geometry of the loss landscape, and the sequential, task-dependent nature of posterior updates. To address these limitations, we establish a general PAC-Bayesian bound for incremental continual learning that incorporates a task-aware flatness term and allows for arbitrary statistical dependencies across tasks.


% While classical PAC-Bayesian analysis provides a generalization guarantee for lifelong learning~\cite{pentinaPACBayesianBoundLifelong2014}, it neither accounts for the flatness of the loss landscape—a property critical for robustness and mitigating forgetting—nor does it directly address the sequential, task-dependent posteriors of incremental learning. To bridge this gap, we first present a general PAC-Bayesian bound for incremental continual learning that explicitly includes a task-aware flatness term and accommodates arbitrary statistical dependencies among tasks



% The central objective is to bridge the gap from $\hat{L}(\mathcal{Q}_{1:n})$ (optimizable) to $L(\mathcal{Q})$ (ultimate goal) through explicit, quantitative generalization bounds.
% This relationship highlights that the empirical continual risk $\hat{L}(\mathcal{Q}_{1:n})$ provides a data-driven estimate for the average risk on observed tasks, which in turn serves as a proxy for the continual generalization risk $L(\mathcal{Q})$ on future tasks. Our goal is to establish a generalization bound that quantitatively connects these quantities and enables reliable prediction of future task performance from observed data.







% We formally distinguish three types of risk in continual learning, as illustrated in Fig.~\ref{fig concept: 3 different definaiton}.

% \begin{itemize}
%     \item \textbf{Continual generalization risk}: the expected performance on a future, unseen task drawn from the task environment,
%     \begin{equation}
%         L(\mathcal{Q}) := \mathbb{E}_{t \sim \mathcal{T}}\mathbb{E}_{P \sim \mathcal{Q}}L_{\mathcal{D}_t}\big(Q_t(\mathcal{D}_t, P)\big),
%     \end{equation}
%     which quantifies the generalization ability for thr future task, but is typically intractable.
%     \item \textbf{Average generalization risk on observed tasks}:
%     \begin{equation}
%         L(\mathcal{Q}_{1:n}) := \frac{1}{n} \sum_{t=1}^n \mathbb{E}_{P_{1:n} \sim \mathcal{Q}_{1:n}} \mathbb{E}_{f \sim Q_t(P_t)} [ L_{\mathcal{D}_t}(f) ],
%     \end{equation}
%     reflecting expected performance over observed tasks but still depending on unknown data distributions.
%     \item \textbf{Empirical continual risk}:
%     \begin{equation}
%         \hat{L}(\mathcal{Q}_{1:n}) := \frac{1}{n} \sum_{t=1}^n \mathbb{E}_{P \sim \mathcal{Q}} \mathbb{E}_{f \sim Q_t} \frac{1}{m_t} \sum_{j=1}^{m_t} \ell(f(x_{tj}), y_{tj}),
%     \end{equation}
%     which is the practical training quantity and the foundation for generalization analysis.
% \end{itemize}

% This relationship highlights that the empirical continual risk $\hat{L}(\mathcal{Q}_{1:n})$ provides a data-driven estimate for the average risk on observed tasks, which in turn serves as a proxy for the continual generalization risk $L(\mathcal{Q})$ on future tasks. Our goal is to establish a generalization bound that quantitatively connects these quantities and enables reliable prediction of future task performance from observed data.


% is the KL divergence between the hyperposterior and hyperprior, reflecting meta-level (hyperprior) complexity, expanded as
% is the KL divergence between the joint hyperposterior and hyperprior, reflecting the meta level complexity. it can be expanded as 

% in the incremantal learning process. 

% :

\begin{theorem}[PAC-Bayes Bound with flatness for Continual Learning]
\label{thm:general-pacbayes}
Consider $n$ sequential tasks under a hierarchical Bayesian continual learning framework. For any hyperposterior $\mathcal{Q}$ and any family of task-wise posteriors $\{Q_t\}$, with probability at least $1-\delta$,
\begin{align}
& L(\mathcal{Q}_{1:n}) \leq \frac{1}{n} \sum_{t=1}^n \mathbb{E}_{P_{t} \sim \mathcal{Q}_{t}} \Big[ \widehat{L}_{S_t}(\mathbf{W}_t) 
+ {\mathrm{E\text{-}Sh}}_t(\mathbf{W}_t, \Sigma_Q) \Big] \notag\\
& \quad  +\frac{1}{n(\bar{m}_n-1)} \sqrt{
\frac{
KL^{\text{task}}_n + KL^{\text{hyper}}_n + \log(\bar{m}_n/\delta)
}{2(\bar{m}_n-1)}
}
\end{align}
where
\begin{itemize}
\item $\mathbf{W}_t$ denotes the model parameters after training on task $t$: for full finetuning, this refers to all updated model weights; for LoRA, it includes the backbone and all LoRA branches up to task $t$.
\item $KL^{\text{task}}_n = KL(Q_{1:n}||P_{1:n})$ denotes the joint KL divergence over all task posteriors and priors, reflecting parameter-level complexity.
\item $KL^{\text{hyper}}_n$ is the KL divergence between the hyperposterior and hyperprior, reflecting the meta-level (hyperprior) complexity. This term captures the dependency among hyperpriors induced by the hierarchical Bayesian modeling, and can be expanded as
$$
\begin{aligned}
    KL^{\text{hyper}}_n
    &=KL(\mathcal{Q}_1 |\mathcal{P}_1)  \\
    & \quad + \sum_{t=2}^n  \mathbb{E}_{\mathcal{Q}_{1:t-1}} [ KL(\mathcal{Q}_t | \mathcal{P}_t(\cdot|\mathcal{Q}_{1:t-1})) ]
\end{aligned}
$$ 
where each hyperprior $\mathcal{P}_t$ depends on the previous tasks in the incremental learning process.
\item $\bar{m}_n = n / (\sum_{t=1}^n 1/m_t)$ is the harmonic mean of the sample sizes.
\end{itemize}
\end{theorem}

% denotes the joint KL divergence over all tasks, reflect the paramter level complexity.
% $$

The general PAC-Bayes bound established above provides a unified and principled framework for continual learning, integrating empirical risk, flatness of the loss landscape, and hierarchical Bayesian dependencies into a single generalization guarantee.   This result highlights the essential role of flatness for mitigating forgetting and improving robustness, while explicitly modeling the sequential nature of knowledge accumulation across tasks. However, the joint KL divergence term $KL(Q_{1:n}|P_{1:n})$ remains challenging to optimize and interpret in practice, primarily due to complex parameter dependencies and the sequential, history-dependent updates of hyperdistributions required by continual learning.
% primarily due to complex parameter dependencies and the recursive updates of hyperdistribution in continual learning.
% the need to update hyperpriors in a sequential, task-dependent manner throughout continual learning.
% due to complex parameter dependencies and the recursive  hyperdistribution updates inherent in continual learning. 
These challenges naturally motivate structural regularization, such as orthogonality constraints, to decouple task-specific adaptations, simplify the KL term, and enable a more tractable and interpretable formulation. We formalize this idea in the following specialized bound.

% These challenges naturally motivate the need for structural regularization—such as orthogonality constraints—which can decouple task-specific adaptations, simplify the KL term, and enable a more tractable and interpretable formulation. This consideration leads us to the following specialized bound.


% This motivates the introduction of explicit regularization to induce a more tractable and interpretable learning process.



% The general PAC-Bayes bound stated above is widely applicable; however, in continual learning, the joint KL divergence term $KL(Q_{1:n} | P_{1:n})$ is challenging to optimize and interpret, as it requires handling the sampling process from recursively updated hyperdistributions across tasks, especially in large-scale, high-dimensional settings.


% While the above bound provides a theoretically general guarantee, its joint KL term $KL(Q_{1:n}\|P_{1:n})$ is difficult to optimize and interpret in practice, especially in large-scale continual learning, where the lack of structural constraints can lead to entangled parameter dependencies and prohibitive computational costs. This motivates the introduction of explicit regularization to induce a more tractable and interpretable learning process.


\subsection{PAC-Bayes Bound for Orthogonal IncLoRA}



To further enhance tractability and interpretability, we impose a stronger structural assumption—orthogonality among task-specific parameter adaptations. This constraint is not only theoretically motivated by our analysis above but also reflects a common strategy in continual learning practice, as seen in works such as OLoRA and orthogonal gradient projection methods, which use orthogonality to minimize task interference and promote knowledge disentanglement. By explicitly incorporating orthogonality, we obtain a sharper and more operationally useful bound that directly guides the design of scalable continual learning algorithms.


\begin{theorem}[PAC-Bayes Bound for Orthogonal Incremental LoRA with Flatness]
Under the orthogonality constraint, the generalization risk satisfies::
\begin{align}
      &  L(\mathcal{Q}_{1:n}) \leq  \frac{1}{n} \sum_{t=1}^{n}
        \mathbb{E}_{P_{1:n} \sim \mathcal{Q}_{1:n}}\big[
            \widehat{L}_{S_t}(\mathbf{W}_{t}) \notag
 + \text{E-Sh}_t(\mathbf{W}_{t} ,\, \Sigma_Q)
        \big] \notag\\
&  \qquad + \frac{(b-a) }{n(\bar{m}_n - 1)}\sqrt{
            \frac{KL_{n}^{\text{task}} + KL_{n}^{\text{hyper}} +\log(\bar{m}_{n}/\delta)}{2(\bar{m}_{n} - 1)}
        }
\end{align}
where $KL^{\text{task}}_n = \sum_{t=1}^n KL(Q_t \| P_t) $.
\end{theorem}

This specialized result demonstrates that, under orthogonality, the KL regularization becomes fully separable across tasks, while empirical risk and flatness are penalized on a per-task basis. Importantly, this formulation systematically combines three key algorithmic elements: empirical risk minimization with flatness promotion, parameter independence across tasks, and the deeper inter-task dependencies captured by the hierarchical hyperprior–hyperposterior structure. In contrast to previous approaches that typically address these aspects in isolation, our theory provides a comprehensive foundation for continual learning method .



\section{Understanding Flatness in LoRA-Based Continual Learning}



\subsection{Experimental Design}

% \begin{table*}[t]
% \centering
% \begin{tabular}{lcc|cc}
% \toprule
% \multirow{2}{*}{\textbf{Method}} & \multicolumn{2}{c|}{\textbf{Perturbation Scope}} & \multicolumn{2}{c}{\textbf{Noise Type}} \\
% \cmidrule(lr){2-3} \cmidrule(lr){4-5}
% & \textbf{LoRA-Only} & \textbf{Global} & \textbf{Gaussian} & \textbf{Fisher-Guided} \\
% \midrule
% LoRA-SAM       & \checkmark &            &             &             \\
% LoRA-Gaussian  & \checkmark &            & \checkmark  &             \\
% LoRA-Fisher    & \checkmark &            &             & \checkmark  \\
% \midrule
% Full-Gaussian  &            & \checkmark & \checkmark  &             \\
% Full-Fisher    &            & \checkmark & \checkmark   & \checkmark  \\
% \bottomrule
% \end{tabular}
% \caption{Summary of perturbation strategies. "LoRA-Only" restricts perturbations to the current adapters; "Global" applies them to the whole model. "Gaussian" refers to isotropic random noise, while "Fisher-Guided" scales perturbations inversely by Fisher information. "Full-Fisher" combines Fisher-guided noise for LoRA parameters and Gaussian noise for the rest.}
% \label{tab:design_flatness}
% \end{table*}

To empirically assess the role of flatness in LoRA-based continual learning, we build on PAC-Bayesian theory, which links local curvature to generalization and forgetting. In practice, LoRA-based methods update only lightweight adapters while keeping the backbone and prior adapters frozen, creating a mismatch between the scope of parameter updates and the assumptions behind global flatness regularization. This discrepancy raises two central questions:
\begin{enumerate}
    \item Is promoting flatness within the LoRA subspace sufficient to improve robustness and mitigate forgetting, or is flatness over the full model necessary?
    \item How do different perturbation strategies trade off knowledge retention, adaptability, and generalization during continual learning?
\end{enumerate}


% Building on the PAC-Bayesian analysis established in the previous section, we aim to empirically investigate how flatness affects continual learning within LoRA-based adaptation. While the PAC-Bayes bound emphasizes the importance of global flatness for generalization and forgetting mitigation, LoRA-based continual learning only updates newly introduced adapters, leaving the backbone and previous adapters frozen. This discrepancy raises two key questions: 
% \begin{itemize}
%     \item Does promoting flatness solely within the updated LoRA subspace suffice to improve robustness, or is global flatness necessary.
%     \item How do different perturbation strategies affect the relationship between flatness, generalization, and forgetting.
% \end{itemize}

 

% To address these questions, we conduct controlled experiments along two factors: \textbf{perturbation scope} and \textbf{perturbation strategy}.

% The first factor, perturbation scope, examines whether restricting perturbations to the updated LoRA adapters is sufficient for improving robustness, or whether promoting flatness across the entire model—including the frozen backbone and prior adapters—yields additional benefits. Accordingly, we consider two settings: LoRA-only, where perturbations are applied solely to the current LoRA adapters; and Global, where perturbations are applied to the entire model.

% The second factor, perturbation type, examines how different noise strategies influence robustness and knowledge retention. Specifically, we evaluate three commonly used forms of perturbation: (i) Sharpness-Aware Minimization (SAM), which applies adversarial perturbations along sharp directions to encourage flatter solutions; (ii) Gaussian noise, representing isotropic unstructured perturbations; and (iii) Fisher-guided noise, where the perturbation magnitude is scaled inversely by accumulated Fisher information to selectively protect parameters deemed critical by prior tasks.

% The second factor, perturbation type, examines how different strategies for introducing perturbations influence robustness and knowledge retention during continual learning. Specifically, we investigate three representative forms of perturbation: Sharpness-Aware Minimization (SAM), Gaussian noise, and Fisher-guided Gaussian noise.
% Each is detailed below.

% The second factor we examine is the perturbation strategy, which determines how noise is introduced to affect robustness and knowledge retention in continual learning. We study three representative methods: (1) adversarial perturbations via SAM, (2) random weight perturbation via isotropic Gaussian noise, and (3) adaptive random weight perturbation  via Fisher-guided structured noise.

To answer these questions, we design a controlled empirical study along two axes:
\paragraph{Perturbation Scope.} We consider two settings:
\begin{itemize}
    \item \textbf{LoRA-only}: Perturbations are applied solely to the newly introduced LoRA adapters.
    \item \textbf{Global}: Perturbations are applied to the entire model, including frozen backbone and past adapters.
\end{itemize}

\paragraph{Perturbation Strategy.} We evaluate three representative noise types commonly used to encourage flatness:

\begin{itemize}
    \item \textbf{Adversarial Perturbation (SAM)}~\cite{foretSharpnessAwareMinimizationEfficiently2021}: Perturbs weights in the direction of steepest loss ascent to induce flatter minima:
    \begin{equation}
    \epsilon^*(\theta, \rho) := \arg \max_{\|\epsilon\|_2 \leq \rho} \hat{L}_S(\theta + \epsilon) \approx \rho \cdot \frac{g}{\|g\|_2}
    \end{equation}

    \item \textbf{Random Weight Perturbation (Gaussian)}: Applies isotropic Gaussian noise across all dimensions:
    \begin{equation}
    \epsilon \sim \mathcal{N}(0,\ \sigma^2 \cdot I)
    \end{equation}

    \item \textbf{Adaptive Random Perturbation (Fisher-guided)}: Modulates the variance of injected noise based on parameter importance. Following \cite{liRevisitingRandomWeight2024}, we employ Fisher information as a sensitivity proxy, but adopt an exponential decay rather than square-root normalization to better preserve relative contrast:
    \begin{equation}
    \epsilon_{ij} \sim \mathcal{N}\left(0,\ \left[\exp\left(-\alpha \cdot \tilde{\mathcal{F}}_{ij} \right)\right]^2 \cdot \sigma^2 \right)
    \end{equation}
    where $\tilde{\mathcal{F}}_{ij}$ is the normalized Fisher score for parameter $(i,j)$.
\end{itemize}


% The second factor we investigate is the \textbf{perturbation strategy}, which controls how noise is introduced during training to promote robustness and generalization. We consider three representative methods:

% \begin{itemize}
%     \item \textbf{Adversarial perturbation (SAM)}\cite{foretSharpnessAwareMinimizationEfficiently2021} applies perturbations along the sharpest gradient direction, encouraging the optimizer to find flatter minima through worst-case local perturbation::
%     \begin{equation}
%     \epsilon^*(\theta, \rho) := \arg \max_{\|\epsilon\|_2 \leq \rho} \hat{L}_S(\theta + \epsilon) \approx \rho \cdot \frac{g}{\|g\|_2}
%     \end{equation}

%     \item \textbf{Random weight perturbation (Gaussian)} applies isotropic, unstructured noise drawn from a standard normal distribution:
%     \begin{equation}
%     \epsilon \sim \mathcal{N}(\mathbf{0},\ \sigma^2 \cdot \mathbf{I})
%     \end{equation}

%     \item  \textbf{Adaptive random perturbation (Fisher-guided)} adjusts the noise magnitude per parameter based on its estimated importance. Following Li et al. \shortcite{liRevisitingRandomWeight2024}, we use Fisher Information to modulate perturbation strength. However, unlike their square-root normalization, which tends to smooth parameter-wise differences, we adopt an exponential decay:
% \begin{equation}
% \epsilon_{ij} \sim \mathcal{N}\left(0,\ \left[\exp\left(-\alpha \cdot \tilde{\mathcal{F}}_{ij} \right)\right]^2 \cdot \sigma^2 \right)
% \end{equation}
% This formulation better preserves the original contrast in importance, allowing sharper suppression of sensitive parameters while maintaining relative variation across dimensions.
% \end{itemize}





% The second factor we investigate is the \textbf{perturbation strategy}, which controls how noise is introduced during training to enhance robustness and generalization. We examine three representative approaches:

% \begin{itemize}
%     \item \textbf{Adversarial perturbation (SAM)} introduces noise along the direction of sharpest local curvature, encouraging the optimizer to find flatter minima through worst-case local perturbation:
%     \begin{equation}
%     \epsilon^*(\theta, \rho) := \arg \max_{\|\epsilon\|_2 \leq \rho} \hat{L}_S(\theta + \epsilon) \approx \rho \cdot \frac{g}{\|g\|_2}
%     \end{equation}

%     \item \textbf{Random weight perturbation (Gaussian)} applies isotropic, unstructured noise drawn from a standard normal distribution:
%     \begin{equation}
%     \epsilon \sim \mathcal{N}(\mathbf{0},\ \sigma^2 \cdot \mathbf{I})
%     \end{equation}

%     \item \textbf{Adaptive random perturbation (Fisher-guided)} adjusts noise magnitude per parameter based on its estimated importance, as captured by the Fisher information:
%     \begin{equation}
%     \epsilon_{ij} \sim \mathcal{N}\left(0,\ \left[\exp\left(-\alpha \cdot \tilde{\mathcal{F}}_{ij} \right)\right]^2 \cdot \sigma^2 \right)
%     \end{equation}
%     where $\tilde{\mathcal{F}}_{ij}$ is the normalized Fisher score for parameter $(i,j)$.
% \end{itemize}

% While our method is related to the Adaptive Random Weight Perturbation (ARWP) proposed by Li et al. (2024), there is a key difference in how the adaptive scaling is implemented. ARWP applies a square-root-based normalization derived from historical gradient magnitudes, which smooths the differences across parameters and leads to a more uniform noise distribution. In contrast, our method applies an exponential decay function to the Fisher scores, which more sharply suppresses noise on critical parameters while retaining relative variation in noise across dimensions. This preserves the original contrast in parameter sensitivity, allowing noise to concentrate more on less informative directions without washing out the structure entirely. As a result, our Fisher-guided perturbation provides a more expressive form of adaptive regularization that better aligns with continual learning needs.
















% The second factor we investigate is the \textbf{perturbation strategy}, which determines how noise is introduced to affect robustness and knowledge retention during continual learning. We study three representative approaches: (1) \textbf{adversarial perturbations} via Sharpness-Aware Minimization (SAM), which explicitly target sharp directions in the loss landscape; (2) \textbf{random weight perturbation} using isotropic Gaussian noise, which explores parameter-space flatness in an unstructured, stochastic manner; and (3) \textbf{adaptive random weight perturbation} guided by Fisher information, which modulates noise magnitude based on the estimated importance of each parameter. The three strategies above can be formalized as follows:

% \begin{itemize}
%     \item \textbf{Adversarial perturbation (SAM):}
%     \begin{equation}
%     \epsilon^*(\theta, \rho) := \arg \max_{\|\epsilon\|_2 \leq \rho} \hat{L}_S(\theta + \epsilon) \approx \rho \cdot \frac{g}{\|g\|_2}
%     \end{equation}

%     \item \textbf{Random weight perturbation (Gaussian):}
%     \begin{equation}
%     \epsilon \sim \mathcal{N}(\mathbf{0},\ \sigma^2 \cdot \mathbf{I})
%     \end{equation}

%     \item \textbf{Adaptive random weight perturbation (Fisher-guided):}
%     \begin{equation}
%     \epsilon_{ij} \sim \mathcal{N}\left(0,\ \left[\exp\left(-\alpha \cdot \tilde{\mathcal{F}}_{ij} \right)\right]^2 \cdot \sigma^2 \right)
%     \end{equation}
%     where
%     \begin{equation}
%     \tilde{\mathcal{F}}_{ij} = \frac{\log(\mathcal{F}_{ij} + \epsilon) - \min}{\max - \min + \epsilon}
%     \end{equation}
% \end{itemize}










% \paragraph{Sharpness-Aware Minimization (SAM).}
% SAM promotes flat minima by introducing worst-case perturbations along gradient directions, thereby explicitly steering optimization toward flatter regions of the loss landscape.  The worst-case perturbation is
% \begin{equation}
% \begin{aligned}
% \epsilon^*(\hat{\theta}, \rho):=\arg \max _{\|\epsilon\|_2 \leq \rho} \hat{L}_S(\theta+\epsilon) 
% \approx \rho \frac{g}{\|g\|_2}
% \end{aligned}
% \end{equation}
% where \( g = \nabla_{\theta} L(\theta) \) is the gradient. The parameters are then updated on the perturbed point \( \theta' = \theta + \epsilon^* \), encouraging robustness to worst-case local perturbations.

% \paragraph{ Isotropic  Gaussian Noise.}
% In contrast to SAM's adversarial design, Gaussian perturbations provide an unstructured, stochastic approach to exploring flatness. Following prior PAC-Bayesian literature, we adopt isotropic Gaussian noise to estimate expected sharpness:
% \begin{equation}
% \epsilon \sim  \mathcal{N}(\mathbf{0}, \sigma^2 I)
% \end{equation}
% This unstructured perturbation promotes general robustness across all directions in the parameter space.

% \paragraph{Fisher-Guided Gaussian Noise.}
% Unlike isotropic noise, Fisher-guided perturbations introduce structure by scaling the variance of each parameter’s noise according to its estimated importance, as measured by the Fisher Information Matrix. Following simiali design \cite{liRevisitingRandomWeight2024}, we designed a 


% \begin{equation}
%     {\epsilon_{ij}} \sim \mathcal{N}\left(\mathbf{0},\  \exp\left(-\alpha \cdot \tilde{\mathcal{F}}_{ij}  \right)^2 \cdot \sigma^2 \cdot \mathbf{I} \right)
% \end{equation}

% where $\tilde{\mathcal{F}}_{ij} = \frac{\log(\mathcal{F}_{ij} + \epsilon) - \min}{\max - \min + \epsilon}$. This formulation reduces noise in highly sensitive (high-Fisher) regions and allows stronger regularization in flat regions.


% \begin{equation}
% \boldsymbol{\epsilon}_{r, t} \sim \mathbb{N}\left(\mathbf{0}, \frac{\sigma^2 \boldsymbol{I}}{\sqrt{1+\eta \sum_{i=1}^{t-1} \beta^{t-i-1} \boldsymbol{g}_i^2}}\right),
% \end{equation}
% where η and β are two hyper-parameters.



% We accumulate a diagonal Fisher estimate for the LoRA parameters via exponential smoothing:
% \begin{equation}
% F^{(i)} = \alpha F^{(i)} + \mathbb{E}_{(x, y) \sim S_t} \left[ \left( \frac{\partial \log p(y|x, W)}{\partial \Delta W} \right)^2 \right], F^{0}=0
% \end{equation}
% where $\alpha$ is a smoothing coefficient, $S_t$ is the current task dataset, and $i$ represent the training process. We impose Gaussian priors whose variance is inversely proportional to Fisher information:
% \begin{equation}
% P(\Delta W) \sim \mathcal{N}\left(0, \frac{\sigma^2}{\sqrt{1+\eta F}}\right),
% \end{equation}
% where $\sigma^2$ is a base variance and $\eta$ controls the influence of Fisher information.

% Among these methods, SAM perturbs parameters in the direction of immediate sharpness, whereas Gaussian noise provides task-agnostic stochasticity. Fisher-guided perturbations offer a middle ground: adapting perturbation strength based on task to preserve important parameters while encouraging flatness elsewhere. In our framework, these distinct perturbation strategies enable us to analyze how structured versus unstructured noise interacts with parameter scope (LoRA-only vs. global) in shaping robustness and generalization during continual learning. The five representative configurations explored in this study are summarized in Table~\ref{tab:design_flatness}.

% \begin{table}[hpbt]
% \centering

% \resizebox{\linewidth}{!}{
% \begin{tabular}{l|cc|cc}
% \toprule
% \textbf{Method} & \textbf{LoRA} & \textbf{Global} & \textbf{Gauss} & \textbf{Fisher} \\
% \midrule
% LoRA-SAM       & \checkmark &            &             &             \\
% LoRA-Gaussian  & \checkmark &            & \checkmark  &             \\
% LoRA-Fisher    & \checkmark &            &             & \checkmark  \\
% \midrule
% Full-Gaussian  &            & \checkmark & \checkmark  &             \\
% Full-Fisher    &            & \checkmark & \checkmark  & \checkmark  \\
% \bottomrule
% \end{tabular}
% }
% \caption{Summary of perturbation strategies. Scope: which parameters are perturbed. Type: how noise is injected.}
% \label{tab:design_flatness}
% \end{table}

% \begin{table}[t]
% \centering
% \resizebox{\linewidth}{!}{
% \begin{tabular}{lcc|cc}
% \toprule
% \multirow{2}{*}{\textbf{Method}} & \multicolumn{2}{c|}{\textbf{Scope}} & \multicolumn{2}{c}{\textbf{Type}} \\
% \cmidrule(lr){2-3} \cmidrule(lr){4-5}
% & \textbf{LoRA} & \textbf{Global} & \textbf{Gauss} & \textbf{Fisher} \\
% \midrule
% LoRA-SAM       & \checkmark &            &             &             \\
% LoRA-Gaussian  & \checkmark &            & \checkmark  &             \\
% LoRA-Fisher    & \checkmark &            &             & \checkmark  \\
% \midrule
% Full-Gaussian  &            & \checkmark & \checkmark  &             \\
% Full-Fisher    &            & \checkmark & \checkmark  & \checkmark  \\
% \bottomrule
% \end{tabular}
% }
% \caption{Summary of perturbation strategies. Scope indicates perturbed parameter set; Type denotes noise generation scheme.}
% \label{tab:design_flatness}
% \end{table}

\subsubsection{Datasets}
We evaluate our methods on three benchmarks covering a spectrum of task diversity and difficulty:

% \subsubsection{Datasets}
% To comprehensively evaluate the impact of flatness interventions on continual learning, we consider benchmarks covering varying levels of task complexity, similarity, and domain diversity. This setup allows us to systematically assess robustness under both conventional and more challenging continual learning scenarios.


\textbf{Short Sequence Benchmark.}  
This benchmark consists of four topic classification datasets—AG News, Amazon Reviews, DBpedia, and Yahoo Answers—from \cite{NIPS2015_250cf8b5}. These tasks are thematically related and represent within-domain transfer scenarios.

% \paragraph{Short Sequence of Tasks.}
% We first evaluate on a standard continual learning benchmark  \shortcite{NIPS2015_250cf8b5} comprising four topic classification datasets: AG News, Amazon Reviews, DBpedia, and Yahoo Answers. This setting represents typical within-domain task sequences where tasks share certain linguistic characteristics.

\textbf{Long Sequence Benchmark.}  
To simulate continual learning under high task heterogeneity, we use a 15-task benchmark following \cite{razdaibiedina2023progressivepromptscontinuallearning}. It includes datasets from CL, GLUE \cite{wang2019gluemultitaskbenchmarkanalysis}, SuperGLUE \cite{NEURIPS2019_4496bf24}, and IMDB \cite{maas2011learning}, with each task comprising 1,000 training examples. This setup introduces domain shift.


% \paragraph{Long Sequence of Tasks.}
% To further examine robustness under more demanding continual learning conditions, we adopt a continual learning benchmark of 15 datasets \shortcite{razdaibiedina2023progressivepromptscontinuallearning}, including five from the CL benchmark, four from GLUE \shortcite{wang2019gluemultitaskbenchmarkanalysis}, five from SuperGLUE \shortcite{NEURIPS2019_4496bf24}, and IMDB \shortcite{maas2011learning}. Each task uses 1000 training examples\shortcite{razdaibiedina2023progressivepromptscontinuallearning}. This setup introduces substantial domain shift and diversity in reasoning complexity.


\textbf{TRACE Benchmark.}  
The TRACE benchmark \cite{wang2023tracecomprehensivebenchmarkcontinual} features eight diverse tasks spanning multiple modalities and domains, including multilingual stance detection (C‑STANCE), text simplification (20Minuten), financial sentiment analysis (FOMC), summarization (MeetingBank), code completion (Py150), science QA (ScienceQA), and mathematical reasoning (NumGLUE-cm/ds). This benchmark offers a rigorous testbed for evaluating generalization in continual learning.

% \paragraph{TRACE Benchmark.}
% We evaluate on the TRACE benchmark\cite{wang2023tracecomprehensivebenchmarkcontinual}, comprising eight tasks: C‑STANCE (multilingual stance detection), 20Minuten (German text simplification), FOMC (financial sentiment classification), MeetingBank (meeting summarization), Py150 (code completion), ScienceQA (science question answering) and NumGLUE‑cm/ds (mathematical reasoning). The benchmark’s diversity in language, task type and reasoning complexity offers a rigorous testbed for assessing continual learning generalization.  


% We employ the TRACE benchmark, which comprises eight diverse tasks covering multilingual stance detection (C-STANCE, Chinese), German text simplification (20Minuten), financial sentiment classification (FOMC), meeting summarization (MeetingBank), code completion (Py150), science question answering (ScienceQA), and two mathematical reasoning tasks (NumGLUE-cm, NumGLUE-ds). . This benchmark features substantial diversity in language , task types, and reasoning complexity, providing a challenging and realistic testbed to systematically assess the generalization ability of large language models in continual learning scenarios.

% Each task contains 5,000 training samples in a standardized format


\subsubsection{Base Models}

We adopt two widely used pretrained language models reflecting different architectural paradigms in NLP. The first is the encoder-decoder \textbf{T5-large model} \shortcite{qin2022lfpt5unifiedframeworklifelong}, which we use for both the short- and long-sequence benchmarks \shortcite{wang2023orthogonal}. The second is \textbf{LLaMA-2-7B-Chat}, a decoder-only model \shortcite{touvron2023llama}, employed for the TRACE benchmark that emphasizes multi-task and multi-domain generalization \shortcite{wang2023tracecomprehensivebenchmarkcontinual}. This combination allows us to evaluate our methods under both autoregressive and sequence-to-sequence settings commonly used in continual learning for NLP.


% We adopt two widely used language models from prior continual learning work in NLP: the encoder-decoder T5 model \shortcite{qin2022lfpt5unifiedframeworklifelong} and the decoder-only LLaMA model \shortcite{touvron2023llama}. Specifically, we use T5-large for both short- and long-sequence tasks \shortcite{wang2023orthogonal}, and LLaMA-2-7B-Chat for multi-task continual learning \shortcite{wang2023tracecomprehensivebenchmarkcontinual}. This selection ensures coverage of both autoregressive and sequence-to-sequence modeling settings.

% \subsubsection{Base Models}
% We adopt two widely used pretrained models reflecting different architectural paradigms in NLP:


% \begin{itemize}
%     \item \textbf{T5-large} \cite{qin2022lfpt5unifiedframeworklifelong}: An encoder-decoder model used for the short and long sequence benchmarks \cite{wang2023orthogonal}.
%     \item \textbf{LLaMA-2-7B-Chat} \cite{touvron2023llama}: A decoder-only model used for TRACE, which emphasizes multi-task and multi-domain generalization \cite{wang2023tracecomprehensivebenchmarkcontinual}.
% \end{itemize}

%  This selection ensures coverage of both autoregressive and sequence-to-sequence modeling settings.

%  \textbf{T5-large} \cite{qin2022lfpt5unifiedframeworklifelong}: An encoder-decoder model used for the short and long sequence benchmarks \cite{wang2023orthogonal}. \textbf{LLaMA-2-7B-Chat} \cite{touvron2023llama}: A decoder-only model used for TRACE, which emphasizes multi-task and multi-domain generalization \cite{wang2023tracecomprehensivebenchmarkcontinual}. This selection ensures coverage of both autoregressive and sequence-to-sequence modeling settings.



% \subsubsection{BaseModel}
% In our experiments, we use two language models adopted by the previous lines of works in CL for NLP: encoder-decoder T5 model \shortcite{qin2022lfpt5unifiedframeworklifelong} and decoder-only LLaMA model\shortcite{touvron2023llama}.  We use T5-large on the short and long sequence task \shortcite{wang2023orthogonal} and use LLaMA-2-7B-Chat for multi task\shortcite{wang2023tracecomprehensivebenchmarkcontinual}.

% \subsubsection{Metrics}
% % \label{metrics: acc flatness}

% % \begin{figure}[thpb]
% %   \centering
% %   \includegraphics[width=0.8\linewidth]{images/eval_metrix.pdf}
% %   \caption{Evaluation metrics across tasks in continual learning}
% %   \label{fig2: Continual eval metric}
% % \end{figure}

% % \textbf{Continual Learning Performance}
% To evaluate continual learning performance, we report four key metrics: Average Accuracy ( AA) \shortcite{NIPS2017_f8752278}, Backward Transfer (BWT) \shortcite{NIPS2017_f8752278}, Forward Transfer (FWT) \shortcite{NIPS2017_f8752278}, and Last Average Accuracy (LA). These metrics correspond to specific regions of the accuracy matrix:  AA is defined by the diagonal entries, FWT by the upper triangular region, LA by the final row, and BWT by the difference between the last row and the diagonal. However, the standard BWT only reflects the final model state, overlooking the full trajectory of forgetting. To better capture this dynamic, we instead report the average backward transfer (ABWT) over the entire training process.

% We employ loss‐landscape visualizations \cite{liVisualizingLossLandscape2018} to directly depict the loss surface. Additionally, we adopt the E-Sh definition to quantify its flatness.


% While these metrics provide a concise summary of continual learning performance, they focus solely on the model's final state, overlooking how forgetting and transfer evolve throughout training. To better capture these dynamics, we use  forward transfer and average backward transfer from whole process, which compute the mean performance over the entire training process. These metrics reflect the overall trends of generalization and forgetting, offering a more comprehensive view of continual adaptation.


% This metrics provides a clear and intuitive assessment of adaptation, retention, and transfer throughout continual learning. 


% In addition to task-aware sharpness, we also utilize \textbf{the maximal eigenvalue of the Hessian matrix}, $\lambda_{\max}$, as a complementary quantitative measure of local curvature at the solution. A lower $\lambda_{\max}$ indicates a flatter and typically more robust minimum. To further illustrate and compare the flatness of solutions obtained under different adaptation strategies, 

% \textbf{Flatness}



\subsubsection{Metrics}
% To comprehensively assess continual learning performance, we follow standard practice \cite{NIPS2017_f8752278} and report four core metrics: average accuracy (AA), backward transfer (BWT), forward transfer (FWT), and last average accuracy (LA). 
% %  AA measures the mean final accuracy across all tasks, while LA reflects performance on the final task. BWT captures the degree of forgetting by comparing post-training performance on earlier tasks to their best historical accuracy. FWT evaluates how well the model generalizes to unseen tasks before observing their data.
% To more faithfully characterize the dynamics of learning and forgetting over time, we extend these metrics by introducing the \textbf{average backward transfer (ABWT)}, computed over the entire training trajectory rather than at the final state only. These averaged variants provide a more robust view of how knowledge is retained or propagated across the task sequence.


% To comprehensively assess continual learning behavior, we evaluate both local task-level performance and global progression over the entire training process. At each training step, we consider three core quantities: the model’s accuracy on the current task immediately after training, denoted by the diagonal entry $R_{i,i}$; the backward transfer for task $i$, defined as the average change in accuracy on task $i$ across all future time steps, $\mathrm{BWT}_i = \frac{1}{T - i - 1} \sum_{j = i+1}^{T-1} (R_{j,i} - R_{i,i})$; and the forward influence of task $i$ on all subsequent tasks, captured by all-future transfer $\mathrm{AF}_i = \frac{1}{T - i - 1} \sum_{j = i+1}^{T-1} R_{i,j}$. Unlike standard forward transfer, which considers only the immediate next task, AF$_i$ reflects the broader generalization of knowledge to future tasks. To capture learning dynamics at the sequence level, we report four global metrics. Average Accuracy (AA) is computed as the mean of the diagonal entries of the accuracy matrix, measuring task performance immediately after each training phase. Last Accuracy (LA) is defined as the average performance of the final model across all tasks, summarizing end-of-training generalization. To complement these, we introduce Average Backward Transfer (ABWT) and Average Forward Transfer (AFWT), which aggregate the BWT$_i$ and AF$_i$ values over the entire task sequence, respectively. Formally,
% $$
% \mathrm{ABWT} = \frac{1}{T - 1} \sum_{i=0}^{T-2} \mathrm{BWT}_i, \quad 
% \mathrm{AFWT} = \frac{1}{T - 1} \sum_{i=0}^{T-2} \mathrm{AF}_i.
% $$
% Together, AA, LA, ABWT, and AFWT offer a comprehensive and temporally aware view of continual learning, capturing immediate task learning, end performance, forward generalization, and accumulated forgetting across time.



To assess continual learning behavior, we distinguish between step-wise and sequence-level metrics. At each training step $t$, we evaluate three quantities: the model’s performance on the current task, termed \textbf{Current Task Accuracy (AC)}, given by the diagonal entry $R_{t,t}$; the \textbf{Current Backward Transfer (BWT)}, which quantifies forgetting on previously learned tasks and is defined as $\mathrm{BWT}_t = \frac{1}{t} \sum_{i=0}^{t-1} (R_{t,i} - R_{i,i})$ and the \textbf{Current All-Future Transfer (AF)}, which measures generalization to all future tasks $\mathrm{AF}_t = \frac{1}{T - t - 1} \sum_{j = t+1}^{T-1} R_{t,j}.$ To capture overall trends, we report four global metrics. \textbf{Average Accuracy (AA)} is the mean of all diagonal entries, reflecting per-task learning immediately after training. \textbf{Last Accuracy (LA)} is the average performance of the final model across all tasks. In addition, we define \textbf{Average Backward Transfer (ABWT)} and \textbf{Average Forward Transfer (AFWT)} by aggregating $\mathrm{BWT}_t$ and $\mathrm{AF}_t$ over the sequence:
\[
\mathrm{ABWT} = \frac{1}{T - 1} \sum_{t = 1}^{T-1} \mathrm{BWT}_t, \quad 
\mathrm{AFWT} = \frac{1}{T - 1} \sum_{t = 0}^{T-2} \mathrm{AF}_t.
\]
% Together, these metrics provide a structured view of adaptation, forgetting, and knowledge transfer throughout the continual learning process.






We visualize the loss landscape following \cite{liVisualizingLossLandscape2018} to qualitatively assess solution flatness. For quantitative evaluation, we use the E-SH metric based on stochastic perturbations (eq. \ref{def: Task-Aware Expected Sharpness}). Unlike Hessian-based measures, which are costly and unstable for large models \cite{pmlr-v70-dinh17b}, E-SH provides an efficient and scalable estimate of local flatness.

% In addition to accuracy-based evaluation, we quantify optimization geometry via expected sharpness (E-Sh) \cite{liVisualizingLossLandscape2018}, which reflects the local curvature of the loss surface. We further include loss landscape visualizations to complement this metric and provide qualitative insight into the flatness of the obtained solutions.







% \begin{table*}[thbp]
% \centering
% \begin{tabular}{ll
%     cccc  % Order1
%     cccc  % Order2
%     cccc  % Order3
% }
% \toprule
%  & \multirow{2}{*}{\textbf{Method}}
%     & \multicolumn{4}{c}{\textbf{Order1}}
%     & \multicolumn{4}{c}{\textbf{Order2}}
%     & \multicolumn{4}{c}{\textbf{Order3}} \\
% \cmidrule(lr){3-6} \cmidrule(lr){7-10} \cmidrule(lr){11-14}
%  & & \textbf{BWT} & \textbf{ AA} & \textbf{FWT} & \textbf{LA}
%    & \textbf{BWT} & \textbf{ AA} & \textbf{FWT} & \textbf{LA}
%    & \textbf{BWT} & \textbf{ AA} & \textbf{FWT} & \textbf{LA} \\
% \midrule
% \multirow{6}{*}{}
%     & SeqLoRA No Noise
%          & -9.02   & 80.66    & 45.59       & 73.41
%         & \textbf{-3.95}      & 79.95      & 44.13     & 75.74
%          &  -7.57 & 76.46  & 48.86   & 76.46 \\
% \midrule
%     & LoRA-SAM
%         & -7.39    & 80.19        & 46.14       & 73.69
%         & -5.14      & 80.10      & 43.93     & 76.07
%          &  -7.60 & 76.12  & \textbf{54.22}   & 76.12 \\
%     & LoRA-Gaussian
%        & -6.65     & 80.70      & \textbf{47.74}       & 74.75
%         & -4.09     & {80.70}      & \textbf{45.70}       & \textbf{76.67}
%        & {-6.91}  & 75.72  & 49.49    & 75.72 \\
%     & LoRA-Fisher
%         & {-6.79}  & 80.43          & 43.31       & \textbf{74.92}
%         & -4.89     & 80.37      & 39.65       & 75.39
%         & -7.50  & 75.54  & 49.38     & 75.54 \\
% \midrule
%     & Full-Gaussian
%        & \textbf{-6.31}    & 80.53         & 44.82      & 74.56
%         & -4.20      & 80.49      & 42.11      & 76
%         &  -7.23 & 76.72  & 50.19    & 76.72 \\
%     & Full-Fisher
%         & -6.50    & \textbf{80.79}        & 44.62       & 74.72
%         & -4.90      & \textbf{80.86}      & 39.71       & 75.27
%         & \textbf{-6.30} & \textbf{77.50}  & 51.22    & \textbf{77.50} \\
% \bottomrule
% \end{tabular}
% \caption{Performance comparison of SeqLoRA}
% \label{tab:seqlora_orders_benchmark}
% \end{table*}

\subsection{Flatness in SeqLoRA: Local vs. Global Effects}


\begin{table*}[thbp]
\centering
\resizebox{\textwidth}{!}{
\begin{tabular}{cc|cccc|cccc|cccc}
\toprule
& & \multicolumn{4}{c|}{\textbf{Order1}} 
  & \multicolumn{4}{c|}{\textbf{Order2}} 
  & \multicolumn{4}{c}{\textbf{Order3}} \\
\cmidrule(lr){3-6} \cmidrule(lr){7-10} \cmidrule(lr){11-14}
\textbf{Baseline} & \textbf{Method} & ABWT &  AA & AFWT & LA 
                  & ABWT &  AA & AFWT & LA 
                  & ABWT &  AA & AFWT & LA \\
\midrule
\multirow{6}{*}{\rotatebox[origin=c]{90}{SeqLoRA}} 
&No Noise         
 & -9.02   & 80.66    & 45.59       & 73.41
& \textbf{-3.95}      & 79.95      & 44.13     & 75.74
&  -7.57 & 76.46  & 48.86   & 76.46 \\

&LoRA-SAM         
  & -7.39    & 80.19        & 46.14       & 73.69
 & -5.14      & 80.10      & 43.93     & 76.07
&  -7.60 & 76.12  & \textbf{54.22}   & 76.12 \\
&LoRA-Gaussian     
 & -6.65     & 80.70      & \textbf{47.74}       & 74.75
 & -4.09     & {80.70}      & \textbf{45.70}       & \textbf{76.67}
 & {-6.91}  & 75.72  & 49.49    & 75.72 \\
&LoRA-Fisher       
& {-6.79}  & 80.43          & 43.31       & \textbf{74.92}
& -4.89     & 80.37      & 39.65       & 75.39
& -7.50  & 75.54  & 49.38     & 75.54 \\

&Full-Gaussian     
 & \textbf{-6.31}    & 80.53         & 44.82      & 74.56
        & -4.20      & 80.49      & 42.11      & 76
        &  -7.23 & 76.72  & 50.19    & 76.72 \\
&Full-Fisher       
 & -6.50    & \textbf{80.79}        & 44.62       & 74.72
        & -4.90      & \textbf{80.86}      & 39.71       & 75.27
        & \textbf{-6.30} & \textbf{77.50}  & 51.22    & \textbf{77.50} \\
\bottomrule
\end{tabular}
}
\caption{Performance of SeqLoRA under flatness-promoting perturbations applied either to the LoRA subspace or the full model. Results are reported across three task orders on the short sequence benchmark. Metrics include average backward transfer (ABWT), average accuracy ( AA),  forward transfer (FWT), and last accuracy (LA).}
\label{tab:comparison_seqlora}
\end{table*}

% We investigate the role of flatness in continual learning under the SeqLoRA framework by comparing flatness-promoting perturbations applied either to the LoRA adapters or to the full model. Table~\ref{tab:comparison_seqlora} summarizes the results in three different short sequence of tasks.



We investigate the role of flatness in continual learning under the SeqLoRA framework by comparing flatness-promoting perturbations applied either to the LoRA adapters or to the full model. Table~\ref{tab:comparison_seqlora} summarizes the results across three task orders on the short-sequence benchmark. In all settings, flatness-aware training improves the ABWT by 2–3 points relative to the baseline without noise, confirming that promoting flatter minima reduces forgetting. Among all methods, LoRA-Gaussian achieves the best trade-off across ABWT, FWT, and LA, demonstrating the value of localized perturbation within the adapter subspace. These results suggest that promoting flatness in the updated LoRA subspace is generally effective in improving stability and generalization. In comparison, global perturbation appears less beneficial and may introduce additional variability due to interference with frozen parameters.


% Across all settings, flatness-aware training consistently improves average backward transfer (ABWT) by 2–3 points over the no-noise baseline, confirming that flatter solutions mitigate forgetting. Among all methods, \textbf{LoRA-Gaussian} achieves the best overall trade-off, yielding the lowest ABWT and the highest  forward transfer (FWT), while maintaining competitive final accuracy and retention. This result supports the hypothesis that promoting flatness locally within the LoRA subspace is sufficient to improve robustness and transfer, without requiring full-model perturbation.








% \begin{table}[htbp]
% \centering
% \begin{tabular}{llcccccccc}
% \toprule
%  & \multirow{2}{*}{\textbf{Method}} 
%     & \multicolumn{4}{c}{\textbf{Short-Dataset}}  \\
% \cmidrule(lr){3-6} 
%  & & \textbf{ AA} & \textbf{FWT} & \textbf{BWT} & \textbf{LA} \\
% \midrule
% \multirow{6}{*}{} 
%     & SeqLoRA No Noise      & 80.8  & 49.0  & -4.33  & 77.50   \\
%     \midrule
%     & LoRA-SAM      & 79.4  & 49.86 & -2.26  & 77.7     \\
%     & LoRA-Gaussian & 79.75 & \textbf{50.2}  & \textbf{-0.86}  & \textbf{79.1}     \\
%     & LoRA-Fisher   & 80.7  & 49.86 & -2.26  & 78.95   \\
%     \midrule
%     & Full-Gaussian & 80.75 & 49.06 & -4.86  & 77.1      \\
%     & Full-Fisher   & \textbf{80.95} & 48.86 & -3.2   & 78.55    \\
% \bottomrule
% \end{tabular}
% \caption{Performance comparison of SeqLoRA on Short Dataset. Metrics include  AA, BWT, FWT, and LA.}
% \label{tab:seqlora_cl_benchmark}
% \end{table}


\begin{figure}[thpb]
  \centering
  \includegraphics[width=1\linewidth]{images/lossland_3x4_zoom.pdf}
  \caption{ Loss landscape visualization of the final model trained on the Short-Sequence benchmark (Order 1), evaluated across four test sets. The top row shows the loss surface under full-model perturbations, while the bottom row applies perturbations only to the LoRA adapters. The middle row presents a zoomed-in view to highlight curvature differences. The visualization reflects the geometry of the learned solution and its performance across tasks.}
  \label{fig:lossland_short}
\end{figure}


To examine the geometric implications, we visualize the loss landscapes in Figure~\ref{fig:lossland_short}. Interestingly, all configurations, including those without perturbations, exhibit broadly similar and flat surfaces. Whether perturbations are applied to the LoRA subspace or to the full model, the basins appear consistently wide and shallow, with only minor visible differences. This suggests that under SeqLoRA, where only a single adapter is trained across tasks, the final solution naturally tends toward flatter regions. As such, visual geometry alone may not capture the subtle effects of perturbation strategies on performance.



% To gain geometric insight, we visualize the loss surface under perturbations in Figure~\ref{fig:lossland_short}. However, a closer look at the loss landscapes in Figure~\ref{fig:lossland_short} reveals a surprising result. Regardless of whether perturbations are applied to the LoRA adapters, the full model, or not used at all, the resulting surfaces appear similarly flat. Across all configurations, the loss basins are wide and shallow, and the visual differences between methods are minimal, even when zoomed in. This suggests that, within the context of LoRA-based continual learning, most solutions tend to lie in relatively flat regions, and the visualized geometry alone offers limited insight into the performance differences observed in Table~\ref{tab:comparison_seqlora}.





% To better understand these effects geometrically, we visualize the loss landscape under perturbations (Figure~\ref{fig:lossland_short}). Perturbing the full model reveals sharp valleys near the optimum, while restricting perturbations to the LoRA subspace produces significantly flatter and broader basins. Among all strategies, Gaussian noise leads to the widest minima, closely aligning with its strong empirical performance in Table~\ref{tab:comparison_seqlora}.

% To understand these improvements geometrically, we examine the loss landscapes in Figure~\ref{fig:lossland_short}. Perturbing the full model reveals sharp curvature around the optimum, while perturbations restricted to LoRA adapters yield much flatter regions. Notably, Gaussian noise most effectively broadens the loss basin, consistent with its superior performance in Table~\ref{tab:comparison_seqlora}.



% The results reveal that promoting flatness locally within LoRA adapters consistently improves robustness against forgetting. Specifically, LoRA-Gaussian achieves the best performance in both backward transfer (BWT: -0.86) and last average accuracy (LA: 79.1), outperforming the baseline and other LoRA-based perturbation methods. This suggests that enhancing the flatness of the updated subspace leads to more robust solutions for continual learning, effectively mitigating catastrophic forgetting.
% In contrast, perturbing the full model parameters (Full-Gaussian and Full-Fisher) shows less stable effects. While Full-Fisher marginally improves accuracy ( AA: 80.95), it fails to enhance BWT or LA and even degrades performance in some cases compared to LoRA-only perturbations. This observation suggests that indiscriminate global perturbations, although capable of reducing the overall sharpness of the solution space, may disrupt the preserved knowledge in frozen parameters and hence offer limited benefits to continual learning stability.



\begin{figure}[hptb]
  \centering
  \includegraphics[width=0.9\linewidth]{images/Paper_SeqLoRA_Heatmap_Full_AccSharpHessian.pdf}
  \caption{Visualization of task accuracy, LoRA expected sharpness, and global expected sharpness across different perturbation strategies on the Short-Dataset benchmark Order1. The first row shows accuracy for each task, the second row shows global expected sharpness, and the third row shows LoRA-specific expected sharpness. }
  \label{fig:short_Heatmap_Global_LoRA_ESharp}
\end{figure}

% Figure~\ref{fig:short_Heatmap_Global_LoRA_ESharp} further quantifies flatness using expected sharpness. In the LoRA subspace, configurations with sharpness below 20 correspond to task accuracies exceeding 90%. In contrast, full-model sharpness often exceeds 40, particularly when accuracy declines below 80%. These patterns confirm a strong inverse correlation between flatness and performance, especially when perturbations are confined to the active adapter subspace.

% Taken together, these three perspectives—performance metrics, loss geometry, and sharpness indicators—converge on a consistent conclusion: LoRA-Gaussian perturbation induces flat, robust solutions that mitigate forgetting without disrupting previous knowledge. These findings also align with our theoretical results (cf. Eq.~8), which link flatness to generalization in PAC-Bayesian terms. Thus, encouraging flatness within the LoRA subspace is both effective and sufficient for parameter-efficient continual learning.

% To better quantify the relationship between flatness and performance, Figure~\ref{fig:short_Heatmap_Global_LoRA_ESharp} presents expected sharpness alongside task accuracy across different perturbation strategies. Both global and LoRA-specific sharpness exhibit a negative correlation with accuracy. Methods such as LoRA-Gaussian and LoRA-Fisher produce lower sharpness in the adapter space, which aligns with higher accuracy and improved forgetting mitigation. Full-Gaussian and Full-Fisher elevated sharpness at the global level, which corresponds to reduced performance. This discrepancy suggests that reducing sharpness outside the actively updated subspace may not benefit learning and can even interfere with previously acquired knowledge.  These results highlight that local flatness within the LoRA subspace more effectively reflects and supports generalization under sequential adaptation.

To quantify this effect more precisely, we compare expected sharpness and task accuracy in Figure~\ref{fig:short_Heatmap_Global_LoRA_ESharp}. A strong negative correlation emerges between flatness and accuracy. When global expected sharpness exceeds 100, task performance tends to degrade, while LoRA-specific sharpness below 20 is consistently associated with accuracy above 90. However, further comparison reveals that LoRA-level sharpness is a more reliable indicator of performance. In contrast, Full-Fisher achieves low LoRA sharpness but high global sharpness, which is not always accompanied by consistent performance gains. These patterns suggest that reducing sharpness outside the updated subspace may introduce instability, especially when the base model remains frozen.


Overall, our results show that promoting flatness improves continual learning performance by enhancing both generalization and knowledge retention. In SeqLoRA, where only a single shared adapter is updated while the backbone and past parameters remain frozen, encouraging flatness in the LoRA subspace offers a focused and effective strategy. Full-model perturbations, by contrast, may disrupt stable components without providing additional benefit. Crucially, the advantage of flatness-oriented perturbation arises not just from the geometry of the final solution but from its influence on the optimization trajectory. By smoothing the loss surface during task transitions, these methods help steer learning away from sharp regions that lead to forgetting and instability. This observation aligns with our PAC-Bayesian analysis: since the posterior in SeqLoRA is concentrated in the adapter subspace, reducing sharpness there tightens the generalization bound without increasing model complexity. Thus, localized flatness is both empirically effective and theoretically grounded.



% Overall, our results demonstrate that enhancing flatness meaningfully improves performance in continual learning. In the context of SeqLoRA, where only a single shared adapter is updated across tasks, promoting flatness within this localized subspace offers a targeted and effective way to improve both generalization and memory retention. Since the backbone and previous task parameters remain frozen, perturbing the full model may unnecessarily disrupt stable representations, especially when global sharpness is already high. Importantly, the effectiveness of flatness-oriented perturbation does not stem solely from the geometry of the final solution, but from its impact on the optimization trajectory. By smoothing the loss surface encountered during task transitions, these methods guide gradient updates away from sharp minima, reducing the risk of catastrophic forgetting and improving forward transfer. This observation aligns with our theoretical analysis based on PAC-Bayesian generalization bounds. In the case of SeqLoRA, Promoting flatness in this subspace therefore directly tightens the generalization bound by lowering the expected sharpness without increasing model complexity. This makes localized perturbation not only effective in practice, but also theoretically justified within our framework. 







% Overall, the results demonstrate that enhancing flatness can effectively improve performance in continual learning. In the context of SeqLoRA, where only a single shared adapter is updated across tasks, promoting flatness within this localized subspace is sufficient to improve both generalization and retention. In contrast, injecting perturbations into the full model—especially when most components are frozen—can introduce unnecessary noise and destabilize previously learned representations. Importantly, the benefit of flatness does not arise from the geometry of the final solution alone, but from its influence on the optimization trajectory. By smoothing the loss surface encountered during training, flatness-oriented strategies adjust the gradient dynamics and guide the model toward more stable, transferable minima that are better suited to accommodate multiple tasks sequentially.






% In the LoRA subspace, sharpness values below 20 consistently correspond to task accuracies above 90\%, whereas full-model perturbations often yield sharpness above 40 when accuracy drops below 80\%. These results reinforce a strong inverse correlation between flatness and generalization. More importantly, they demonstrate that LoRA-specific flatness serves as a more direct and reliable indicator of continual learning stability than global sharpness.


% Together, the performance metrics, geometric analysis, and flatness profiles converge on a consistent insight: perturbing only the LoRA subspace, particularly with Gaussian noise, produces flatter and more stable solutions that enhance forward transfer and mitigate forgetting. These findings align with PAC-Bayesian theory, which links flatness to generalization, and suggest that local flatness within the actively updated subspace is not only sufficient but preferable for parameter-efficient continual learning.


% To further understand the connection between flatness and performance, we visualize the expected sharpness of both the global model and the LoRA-specific parameters, as shown in Figure~\ref{short_Heatmap_Global_LoRA_ESharp}. These results indicate that global expected sharpness correlates well with overall accuracy ( AA), aligning with PAC-Bayesian interpretations that flatter solutions tend to generalize better. However, changes in LoRA-specific sharpness more directly correspond to improvements in forgetting mitigation (BWT) and retention (LA). This suggests that robustness within the actively updated subspace plays a more critical role in continual learning dynamics than the global flatness alone.




% In summary, these findings reveal a clear distinction between global and local flatness effects. While global flatness affects generalization, promoting local robustness within LoRA adapters is more effective for improving knowledge retention and stability across sequential tasks. 


% \begin{table*}[htbp]
% \centering
% \begin{tabular}{llcccccccccccc}
% \toprule
%  & \multirow{2}{*}{\textbf{Method}} 
%     & \multicolumn{4}{c}{\textbf{Short-Dataset}} 
%     & \multicolumn{4}{c}{\textbf{Long-Dataset}} 
%     & \multicolumn{4}{c}{\textbf{TRACE}} \\
% \cmidrule(lr){3-6} \cmidrule(lr){7-10} \cmidrule(lr){11-14}
%  & & \textbf{ AA} & \textbf{FWT} & \textbf{BWT} & \textbf{LA}
%    & \textbf{ AA} & \textbf{FWT} & \textbf{BWT} & \textbf{LA}
%    & \textbf{ AA} & \textbf{FWT} & \textbf{BWT} & \textbf{LA} \\
% \midrule
% \multirow{6}{*}{} 
%     & No Noise      & 80.8  & 49.0  & -4.33  & 77.50  & 81.00  & 60.78  & -45.16 & 38.85  & - & - & - & - \\
%     & LoRA-SAM      & 79.4  & 49.86 & -2.26  & 77.7   & 79.64  & 60.54  & -65.93 & 18.10  & - & - & - & - \\
%     & LoRA-Gaussian & 79.75 & 50.2  & -0.86  & 79.1   & 81.42  & 56.44  & -53.14 & 31.82  & - & - & - & - \\
%     & LoRA-Fisher   & 80.7  & 49.86 & -2.26  & 78.95  & 80.8   & 59.13  & -38.28 & 45.06  & - & - & - & - \\
%     & Full-Gaussian & 80.75 & 49.06 & -4.86  & 77.1   & 81.42  & 58.4   & -57.47 & 27.78  & - & - & - & - \\
%     & Full-Fisher   & 80.95 & 48.86 & -3.2   & 78.55  & 81.37  & 58.67  & -53.14 & 31.54  & - & - & - & - \\
% \bottomrule
% \end{tabular}
% \caption{Performance comparison of SeqLoRA on Short-, Long-Dataset and TRACE benchmarks. Metrics include  AA, BWT, FWT, and LA.}
% \label{tab:seqlora_cl_benchmark}
% \end{table*}





% \begin{table}[htbp]
% \centering
% \begin{tabular}{llcccccccc}
% \toprule
%  & \multirow{2}{*}{\textbf{Method}} 
%     & \multicolumn{4}{c}{\textbf{Short-Dataset}}  \\
% \cmidrule(lr){3-6} 
%  & & \textbf{ AA} & \textbf{FWT} & \textbf{BWT} & \textbf{LA} \\
% \midrule
% \multirow{6}{*}{} 
%     & No Noise      & 80.8  & 49.0  & -4.33  & 77.50   \\
%     & LoRA-SAM      & 79.4  & 49.86 & -2.26  & 77.7     \\
%     & LoRA-Gaussian & 79.75 & 50.2  & -0.86  & 79.1     \\
%     & LoRA-Fisher   & 80.7  & 49.86 & -2.26  & 78.95   \\
%     & Full-Gaussian & 80.75 & 49.06 & -4.86  & 77.1      \\
%     & Full-Fisher   & 80.95 & 48.86 & -3.2   & 78.55    \\
% \bottomrule
% \end{tabular}
% \caption{Performance comparison of SeqLoRA on Short Dataset. Metrics include  AA, BWT, FWT, and LA.}
% \label{tab:seqlora_cl_benchmark}
% \end{table}


















% \begin{table*}[htp]
% \centering
% \label{tab:seqlora_cl_benchmark}
% \begin{tabular}{lcccccccc}
% \toprule
% \multirow{2}{*}{\textbf{Method}} & \multicolumn{4}{c}{\textbf{Short-Dataset}} & \multicolumn{4}{c}{\textbf{Long-Dataset}} \\
% \cmidrule(lr){2-5} \cmidrule(lr){6-9}
% & \textbf{ AA} & \textbf{BWT} & \textbf{FWT} & \textbf{LA} & \textbf{ AA} & \textbf{BWT} & \textbf{FWT} & \textbf{LA} \\
% \midrule
% No Noise      & 80.8 & -4.33 & 49   & 77.50 & 81.00 & -45.16 & 60.78 & 38.85 \\
% LoRA-SAM      & 79.4 & -2.26 & 49.86 & 77.7 & 79.64 & -65.93 & 60.54 & 18.10 \\
% LoRA-Gaussian & 79.75 & -0.86 & 50.2 & 79.1 & 81.42 & -53.14 & 56.44 & 31.82 \\
% LoRA-Fisher   & 80.7 & -2.26 & 49.86 & 78.95 & 80.8 & -38.28 & 59.13 & 45.06 \\
% Full-Gaussian & 80.75 & -4.86 & 49.06 & 77.1 & 81.42 & -57.47 & 58.4 & 27.78 \\
% Full-Fisher   & 80.95 & -3.2 & 48.86 & 78.55 & 81.37 & -53.14 & 58.67 & 31.54 \\
% \bottomrule
% \end{tabular}
% \caption{Performance comparison of SeqLoRA on Short- and Long-Dataset benchmarks. Metrics include  AA, BWT, FWT, and LA.}
% \end{table*}














\subsection{Flatness in IncLoRA: Incremental Adaptation Under Flatness}


% \begin{table*}[thbp]
% \centering
% \resizebox{\textwidth}{!}{
% \begin{tabular}{cc|cccc|cccc|cccc}
% \toprule
% & & \multicolumn{4}{c|}{\textbf{Order4}} 
%   & \multicolumn{4}{c|}{\textbf{Order5}} 
%   & \multicolumn{4}{c}{\textbf{Order6}} \\
% \cmidrule(lr){3-6} \cmidrule(lr){7-10} \cmidrule(lr){11-14}
% \textbf{Baseline} & \textbf{Method} & ABWT &  AA & AFWT & LA 
%                   & ABWT &  AA & AFWT & LA 
%                   & ABWT &  AA & AFWT & LA \\
% \midrule
% SeqLoRA & No Noise  & -14.45 & 81.19 & 20.21 & 70.28   & -22.57 & 81.99 & 14.22 & 68.11   & -16.61 & 81.00 & 43.84 & 38.85 \\
% \cmidrule(lr){1-14}
% \multirow{6}{*}{\rotatebox[origin=c]{90}{IncLoRA}} 
% & No Noise        & -12.23 & 80.09 & 22.18 & 70.13   & -14.60 & \textbf{81.03} & 14.31 & 70.13   & -13.77 & 66.70 & 43.52 & 43.16 \\
% & LoRA-SAM        & -10.72 & 79.91 & \textbf{24.62} & 68.69   & -16.45 & 80.11 & 14.40 & 68.52   & \textbf{-9.45} & 69.68 & 45.39 & 70.69 \\
% & LoRA-Gaussian   & -11.32 & \textbf{80.65} & 24.42 & \textbf{70.44}   & -17.96 & 80.87 & 13.96 & 71.07  & -10.01 & 70.60 & 44.11 & 67.52 \\
% & LoRA-Fisher     & \textbf{-10.50} & 80.63 & 23.68 & 69.29   & \textbf{-12.44} & 72.57 & \textbf{16.90} & \textbf{72.52}   & -9.46 & \textbf{70.91} & 43.02 & \textbf{72.20} \\
% & Full-Gaussian   & -12.87 & 80.43 & 23.35 & 69.48   & -18.67 & 80.67 & 15.26 & 69.17   & -9.80 & 70.86 & 44.54 & 65.16 \\
% & Full-Fisher     & -10.95 & 80.05 & 24.48 & 69.24  & -16.52 & 80.73 & 14.04 & 70.02   & -9.00 & 70.54 & \textbf{47.17} & 70.75 \\
% \bottomrule
% \end{tabular}
% }
% \caption{Performance of IncLoRA with fatness-promoting perturbation applied to the LoRA subspace versus the full model across three task orders on the long sequence benchmark.}
% \label{tab:comparison_seqlora_inclora_vertical}
% \end{table*}



\begin{table*}[thbp]
\centering
\resizebox{\textwidth}{!}{
\begin{tabular}{cc|cccc|cccc|cccc}
\toprule
& & \multicolumn{4}{c|}{\textbf{Order4}} 
  & \multicolumn{4}{c|}{\textbf{Order5}} 
  & \multicolumn{4}{c}{\textbf{Order6}} \\
\cmidrule(lr){3-6} \cmidrule(lr){7-10} \cmidrule(lr){11-14}
\textbf{Baseline} & \textbf{Method} & ABWT &  AA & AFWT & LA 
                  & ABWT &  AA & AFWT & LA 
                  & ABWT &  AA & AFWT & LA \\
\midrule
SeqLoRA & No Noise  & -14.45 & 81.19 & 20.21 & 70.28   & -22.57 & 81.99 & 14.22 & 68.11   & -16.61 & 81.00 & 43.84 & 38.85 \\
\cmidrule(lr){1-14}
\multirow{6}{*}{\rotatebox[origin=c]{90}{IncLoRA}} 
& No Noise        & -12.23 & 80.09 & 22.18 & 70.13   & -14.60 & \textbf{81.03} & 14.31 & 70.13   & -13.77 & 66.70 & 43.52 & 43.16 \\
& LoRA-SAM        & -10.72 & 79.91 & \textbf{24.62} & 68.69   & -16.45 & 80.11 & 14.40 & 68.52   & \textbf{-9.45} & 69.68 & 45.39 & 70.69 \\
& LoRA-Gaussian   & -11.32 & \textbf{80.65} & 24.42 & \textbf{70.44}   & -17.96 & 80.87 & 13.96 & 71.07  & -10.01 & 70.60 & 44.11 & 67.52 \\
& LoRA-Fisher     & \textbf{-10.50} & 80.63 & 23.68 & 69.29   & \textbf{-12.44} & 72.57 & \textbf{16.90} & \textbf{72.52}   & -9.46 & \textbf{70.91} & 43.02 & \textbf{72.20} \\
& Full-Gaussian   & -12.87 & 80.43 & 23.35 & 69.48   & -18.67 & 80.67 & 15.26 & 69.17   & -9.80 & 70.86 & 44.54 & 65.16 \\
& Full-Fisher     & -10.95 & 80.05 & 24.48 & 69.24  & -16.52 & 80.73 & 14.04 & 70.02   & -9.00 & 70.54 & \textbf{47.17} & 70.75 \\
\midrule
\multirow{6}{*}{\rotatebox[origin=c]{90}{OLoRA}} 
& No Noise        & -6.27 & 75.33 & 23.42 & 66.95   & \textbf{-7.70} & 74.91 & 13.35 & 66.91   & -5.14 & 70.37  &\textbf{56.05}  & 70.84 \\
& LoRA-SAM        & -6.55 & \textbf{75.85} & 24.02 & \textbf{68.04}   & -9.22 & \textbf{76.12} & 13.40 & \textbf{66.99}   & -4.00 & 71.62 & 52.34 & 71.68 \\
& LoRA-Gaussian   & \textbf{-5.84} & 75.80 & 24.20 & 66.41   & -10.26 & 75.60 & 13.46 & 65.21   & \textbf{-3.18} & \textbf{72.45} & 54.94 & \textbf{73.11} \\
& LoRA-Fisher     & -6.15 & 75.27 & \textbf{24.32} & 66.43   & -10.22 & 75.80 & \textbf{13.65} & 66   & -3.85 & 72.29 & 54.65 & 72.17 \\
& Full-Gaussian   & -5.19 & 74.79 & 23.35 & 67.64   & -10.02 & 75.52 & {13.19} & 65.03   & -4.61 & 72.23 & 55.72 & 72.80 \\
& Full-Fisher     & -6.01 & 75.36 & 24.08 & 67.93   & -10.20 & 75.07 & 13.22 & 65.73   & -4.11 & 71.88 & 55.46 & 72.76 \\
% \bottomrule
\bottomrule
\end{tabular}
}
\caption{Performance of IncLoRA with fatness-promoting perturbation applied to the LoRA subspace versus the full model across three task orders on the long sequence benchmark.}
\label{tab:comparison_seqlora_inclora_vertical}
\end{table*}



We now shift to IncLoRA, a structurally decoupled framework where a new task-specific adapter is introduced at the start of each task and frozen after training, allowing each task to update its own parameters without interfering with previously learned knowledge.Within this setup, we examine whether flatness-promoting perturbations remain beneficial for IncLoRA continual learning.


Table~\ref{tab:comparison_seqlora_inclora_vertical} presents results across three task orders on Long Sequence Benchmark. Consistent with findings in SeqLoRA, perturbing only the current LoRA adapter outperforms full-model perturbation across all metrics. For example, in Order 5, LoRA-Fisher achieves an ABWT of 12.44, outperforming Full-Gaussian while also yielding higher last accuracy. Similar trends are observed in Orders 4 and 6. These results reinforce the conclusion that perturbations should be restricted to actively updated components, as injecting noise into frozen adapters or the backbone may compromise knowledge retention and reduce overall stability.





% To explore this, we compare performance across three task orders (Table~\ref{tab:comparison_seqlora_inclora_vertical}). As in SeqLoRA, restricting perturbations to the LoRA subspace consistently outperforms full-model perturbation. For instance, in Order 5, LoRA-Fisher achieves an ABWT of –12.44, substantially better than Full-Gaussian, while also achieving higher LA. This pattern is replicated across Orders 4 and 6, reinforcing the earlier conclusion that perturbations should remain confined to the actively updated subspace. Full-model noise, particularly when affecting frozen adapters or the backbone, can destabilize preserved knowledge and lead to performance degradation.


% Beyond perturbation scope, we examine how different strategies influence learning dynamics. Figure~\ref{fig:Order4_Inc_Process} tracks the progression of BWT, AC (current task accuracy), and FWT over all 15 tasks. Despite similar trends in overall learning curves, key differences emerge during task transitions. For example, the WiC task introduces a significant distributional shift. Under this challenge, the no-noise baseline suffers a marked drop in BWT, whereas LoRA-SAM and LoRA-Fisher maintain stable retention. This suggests that flatness plays a crucial role in absorbing task-specific shocks without erasing prior knowledge.



\begin{figure}[t]
  \centering
  \includegraphics[width=0.9\linewidth]{images/Combined_Orders_BWT_CTA_AF_2.pdf}
  \caption{Performance comparison of different methods on the Long Sequence Benchmark with three different task orders.}
  \label{fig:Order4_Inc_Process}
\end{figure}


Beyond the scope of perturbation, we further investigate the impact of different perturbation strategies on learning dynamics. Figure~\ref{fig:Order4_Inc_Process} illustrates the evolution of  BWT,  task accuracy (AC), and all future task (AF) across all 15 tasks. Although overall trajectories appear similar, notable differences emerge at key transitions. For instance, the introduction of the WiC task causes a sharp degradation in BWT for the no-noise baseline. In contrast, LoRA-SAM and LoRA-Fisher maintain stable retention, suggesting that flatness-based regularization can buffer against abrupt task shifts and mitigate catastrophic forgetting.
Across all task orders, LoRA-SAM exhibits the smoothest forgetting trajectory, with consistent BWT curves and strong forward transfer. LoRA-Fisher achieves the best overall ABWT and LA, demonstrating superior long-term retention despite slightly higher variance. LoRA-Gaussian delivers balanced performance across all metrics, offering a simple yet robust solution. In contrast, full-model perturbation strategies such as Full-Gaussian and Full-Fisher produce more erratic behavior, occasionally improving forward transfer but often at the expense of stability and memory retention.


These results further confirm that local flatness within the updated LoRA subspace remains a critical factor for continual learning, even in structurally decoupled settings like IncLoRA. Each perturbation method imposes a different inductive bias: SAM promotes temporal smoothness, Fisher prioritizes retention, and Gaussian provides generalizable robustness. Together, they extend the conclusions drawn from SeqLoRA and emphasize that adapter-level flatness is both sufficient and necessary for effective and stable adaptation across sequential tasks.



% Across all task orders, LoRA-SAM displays the most stable forgetting behavior, with smooth BWT dynamics and strong FWT. In contrast, LoRA-Fisher yields the best overall ABWT and LA, highlighting its advantage in long-term retention despite slightly more variability. LoRA-Gaussian offers consistent, moderate improvements across metrics, confirming its robustness and simplicity. In contrast, full-model strategies such as Full-Gaussian and Full-Fisher yield more erratic behavior, occasionally boosting FWT but often at the expense of memory retention.

% Taken together, these results reinforce the role of local flatness as a key driver of continual learning, even under structurally decoupled adaptation. While SAM prioritizes stability, Fisher favors long-term retention, and Gaussian offers a robust, generalizable balance. These dynamics echo earlier findings in SeqLoRA and confirm that promoting flatness within the actively updated subspace remains both necessary and effective for continual adaptation.





% % To evaluate how different flatness-promoting strategies affect continual learning under IncLoRA, we analyze performance across three task orders in both table and trajectory plots (Figure\~\ref{fig\:Order4\_Inc\_Process}). Across all settings, noise-based perturbations applied within the LoRA subspace consistently outperform the no-noise baseline in backward transfer (BWT) and forward transfer (FWT), confirming the benefits of flatness regularization.

% Notably, while all methods follow similar performance trends over the 15-task sequence, key differences emerge when facing distributional shifts. Around challenging tasks such as WiC (task 3), the no-noise baseline suffers sharp BWT degradation, while LoRA-SAM maintains stable retention, indicating improved robustness under task drift. LoRA-Fisher also shows strong retention, achieving the best ABWT in Order 5, while LoRA-Gaussian provides moderate improvements with smoother dynamics.

% Current task accuracy (CTA) remains stable across all methods, suggesting that flatness interventions do not hinder new task acquisition. In contrast, FWT highlights more substantial differences. The no-noise baseline consistently underperforms, especially in later tasks, while LoRA-Fisher achieves the highest FWT in two of three orders, indicating better generalization to unseen tasks.

% These results demonstrate that while all perturbation strategies enhance continual learning under IncLoRA, they exhibit distinct strengths: LoRA-SAM is most robust to sharp domain shifts, LoRA-Fisher balances stability and transfer, and LoRA-Gaussian offers lightweight, consistent improvements. This suggests that selecting a perturbation strategy should be guided by the nature of task transitions and the target trade-off between retention and generalization.






% We now turn to the IncLoRA framework to assess whether the benefits of flatness observed in SeqLoRA persist under a structurally decoupled learning setting. 




% Table~\ref{tab:comparison_seqlora_inclora_vertical} compares performance across three task orders on a long-sequence benchmark. SeqLoRA achieves higher average accuracy, reflecting its stronger adaptability to new tasks. However, this comes at the cost of greater forgetting. In contrast, IncLoRA shows better memory retention, with improved ABWT and LA. These results confirm the core trade-off between adaptability and stability: SeqLoRA optimizes for current-task performance, while IncLoRA prioritizes long-term knowledge preservation through structural separation.


% The previous section examined SeqLoRA, where parameter updates accumulate within a shared low-rank subspace. In contrast, IncLoRA attaches a new task-specific adapter at each stage and freezes it after training. This design minimizes interference across tasks through structural decoupling. However, it remains an open question whether local flatness continues to affect robustness and generalization when parameter sharing is absent.


% To understand the behavioral differences between SeqLoRA and IncLoRA, we compare their performance under identical long-sequence settings. As shown in Table~\ref{tab:comparison_seqlora_inclora_vertical}, SeqLoRA achieves a higher average accuracy in all tasks, reflecting its stronger adaptability to new tasks. This is because SeqLoRA uses a single shared LoRA adapter, which makes it easier for the model to adapt to new tasks. However, this flexibility comes at the cost of greater interference with prior knowledge, which is evident from the lower backward transfer. In contrast, IncLoRA exhibits stronger retention of earlier tasks, with improved backward transfer and final-task accuracy. This improvement stems from its incremental adapter-freezing mechanism, which isolates task-specific parameters after training. By preventing subsequent updates from overwriting earlier representations, IncLoRA preserves knowledge more effectively, though at the cost of reduced adaptability and parameter reuse.



% To validate the two main conclusions drawn from the SeqLoRA setting, namely that promoting flatness improves continual learning performance and that the scope of perturbation has a significant impact, we examine the IncLoRA results. Across all task orders, perturbing only the LoRA adapters leads to better outcomes than perturbing the full model. For instance, in Order 4, LoRA-Fisher reduces the ABWT from 12.23 to 10.50, while Full-Gaussian increases it to 12.87. This trend remains consistent in Orders 5 and 6. These findings confirm that enhancing flatness within the actively updated subspace is effective, whereas applying perturbations to the entire model, including frozen parameters, may destabilize prior knowledge and hinder retention.


% While the distinction between LoRA-only and full-model perturbations remains consistent with the observations in SeqLoRA, here we focus on the effect of different perturbation strategies within the LoRA subspace under IncLoRA. As shown in Table~\ref{tab:comparison_seqlora_inclora_vertical}, LoRA-Gaussian achieves the highest average accuracy ( AA = 80.65 in Order4) and a strong last-task accuracy (LA = 70.44), indicating robust generalization and stable final performance. It also delivers consistent forward transfer (FWT = 24.42), suggesting that unstructured noise helps explore flatter regions without disrupting adaptation.

% LoRA-Fisher, while slightly behind in overall accuracy, yields the best backward transfer in Order4 (ABWT = –10.50) and the highest last accuracy in Order6 (LA = 72.20), confirming its effectiveness in preserving past knowledge. This aligns with the design of Fisher-guided perturbations, which suppress noise on important parameters and promote selective flatness. LoRA-SAM achieves the highest FWT (24.62) in Order4, indicating improved immediate task adaptation, but its performance fluctuates across orders, and it underperforms in long-term stability metrics like ABWT and LA.

% These differences are further illustrated in Figure~\ref{fig:Order6_Inc_Process}. In the BWT plot (top), LoRA-Fisher consistently outperforms other strategies on earlier tasks, reflecting better retention. In contrast, SAM exhibits sharp drops, particularly in tasks 5 to 10, confirming its limited ability to maintain prior knowledge. The CTA plot (middle) shows that LoRA-Gaussian maintains the most stable task-wise accuracy throughout the sequence, while Fisher slightly improves in later tasks. In FWT (bottom), both Gaussian and SAM achieve higher early-task transfer, but Fisher gradually catches up and becomes more balanced toward the end.

% Taken together, these results show that Gaussian noise offers the most stable and general improvement across metrics, Fisher provides a more conservative trade-off favoring knowledge retention, and SAM boosts short-term adaptation but suffers from instability over longer task sequences. This confirms that perturbation strategies affect not only optimization geometry but also the stability–plasticity trade-off critical to continual learning.




% \begin{figure}[t]
%   \centering
%   \includegraphics[width=1\linewidth]{images/Paper_Combined_Orders_BWT_CTA_FWT-all.pdf}
%   \caption{Performance comparison of different methods on the Long Sequence Benchmark with three different task orders.}
%   \label{fig:Order4_Inc_Process}
% \end{figure}
















% \begin{table*}[thbp]
% \centering
% \caption{Performance comparison of two baselines: \textbf{SeqLoRA} and \textbf{IncLoRA}. We report BWT,  AA, FWT, and LA across different methods under Order1.}
% \resizebox{\textwidth}{!}{
% \begin{tabular}{c|l|cccc}
% \toprule
% \textbf{Baseline} & \textbf{Method} & \textbf{BWT} & \textbf{ AA} & \textbf{FWT} & \textbf{LA} \\
% \midrule
% \multirow{6}{*}{\rotatebox[origin=c]{90}{SeqLoRA}} 
%     & No Noise       & -9.02 & 80.66 & 45.59 & 73.41 \\
%     & LoRA-SAM       & -7.39 & 80.19 & 46.14 & 73.69 \\
%     & LoRA-Gaussian  & -6.65 & 80.70 & 47.74 & 74.75 \\
%     & LoRA-Fisher    & -6.79 & 80.43 & 43.31 & 74.92 \\
%     & Full-Gaussian  & -6.31 & 80.53 & 44.82 & 74.56 \\
%     & Full-Fisher    & -6.50 & 80.79 & 44.62 & 74.72 \\
% \midrule
% \multirow{6}{*}{\rotatebox[origin=c]{90}{IncLoRA}} 
%     & No Noise       & --    & --    & --    & --    \\
%     & LoRA-SAM       & --    & --    & --    & --    \\
%     & LoRA-Gaussian  & --    & --    & --    & --    \\
%     & LoRA-Fisher    & --    & --    & --    & --    \\
%     & Full-Gaussian  & --    & --    & --    & --    \\
%     & Full-Fisher    & --    & --    & --    & --    \\
% \bottomrule
% \end{tabular}
% }
% \label{tab:seqlora_inclora_order1}
% \end{table*}









% \begin{table*}[thbp]
% \centering
% \resizebox{\textwidth}{!}{
% \begin{tabular}{cc|cccc|cccc|cccc}
% \toprule
% & & \multicolumn{4}{c|}{\textbf{Order4}} 
%   & \multicolumn{4}{c|}{\textbf{Order5}} 
%   & \multicolumn{4}{c}{\textbf{Order6}} \\
% \cmidrule(lr){3-6} \cmidrule(lr){7-10} \cmidrule(lr){11-14}
% \textbf{Baseline} & \textbf{Method} & BWT &  AA & FWT & LA 
%                   & BWT &  AA & FWT & LA 
%                   & BWT &  AA & FWT & LA \\
% \midrule
% \multirow{6}{*}{\rotatebox[origin=c]{90}{SeqLoRA}} 
% & No Noise        & -- & -- & -- & --   & -- & -- & -- & --   & -- & -- & -- & -- \\
% & LoRA-SAM        & -- & -- & -- & --   & -- & -- & -- & --   & -- & -- & -- & -- \\
% & LoRA-Gaussian   & -- & -- & -- & --   & -- & -- & -- & --   & -- & -- & -- & -- \\
% & LoRA-Fisher     & -- & -- & -- & --   & -- & -- & -- & --   & -- & -- & -- & -- \\
% & Full-Gaussian   & -- & -- & -- & --   & -- & -- & -- & --   & -- & -- & -- & -- \\
% & Full-Fisher     & -- & -- & -- & --   & -- & -- & -- & --   & -- & -- & -- & -- \\
% \midrule
% \multirow{6}{*}{\rotatebox[origin=c]{90}{IncLoRA}} 
% & No Noise        & -12.23 & -- & -- & --   & -- & -- & -- & --   & -- & -- & -- & -- \\
% & LoRA-SAM        & -- & -- & -- & --   & -- & -- & -- & --   & -- & -- & -- & -- \\
% & LoRA-Gaussian   & -- & -- & -- & --   & -- & -- & -- & --   & -- & -- & -- & -- \\
% & LoRA-Fisher     & -- & -- & -- & --   & -- & -- & -- & --   & -- & -- & -- & -- \\
% & Full-Gaussian   & -- & -- & -- & --   & -- & -- & -- & --   & -- & -- & -- & -- \\
% & Full-Fisher     & -- & -- & -- & --   & -- & -- & -- & --   & -- & -- & -- & -- \\
% \midrule
% \multirow{6}{*}{\rotatebox[origin=c]{90}{OLoRA}} 
% & No Noise        & -- & -- & -- & --   & -- & -- & -- & --   & -- & -- & -- & -- \\
% & LoRA-SAM        & -- & -- & -- & --   & -- & -- & -- & --   & -- & -- & -- & -- \\
% & LoRA-Gaussian   & -- & -- & -- & --   & -- & -- & -- & --   & -- & -- & -- & -- \\
% & LoRA-Fisher     & -- & -- & -- & --   & -- & -- & -- & --   & -- & -- & -- & -- \\
% & Full-Gaussian   & -- & -- & -- & --   & -- & -- & -- & --   & -- & -- & -- & -- \\
% & Full-Fisher     & -- & -- & -- & --   & -- & -- & -- & --   & -- & -- & -- & -- \\
% \bottomrule
% \end{tabular}
% }
% \caption{Performance comparison between \textbf{SeqLoRA} and \textbf{IncLoRA} across three task orders. We report Backward Transfer (BWT), Accuracy ( AA), Forward Transfer (FWT), and Learning Accuracy (LA).}
% \label{tab:comparison_seqlora_inclora_vertical}
% \end{table*}






% The previous section examined SeqLoRA, where parameter updates accumulate within a shared low-rank subspace. In contrast, IncLoRA attaches a new task-specific adapter at each stage and freezes it after training. This design minimizes interference across tasks through structural decoupling. However, it remains an open question whether local flatness continues to affect robustness and generalization when parameter sharing is absent.


% To understand the behavioral differences between SeqLoRA and IncLoRA, we compare their performance under identical long-sequence settings. As shown in Table~\ref{tab:comparison_seqlora_inclora_vertical}, SeqLoRA achieves a higher average accuracy in all tasks, reflecting its stronger adaptability to new tasks. This is because SeqLoRA uses a single shared LoRA adapter, which makes it easier for the model to adapt to new tasks. However, this flexibility comes at the cost of greater interference with prior knowledge, which is evident from the lower backward transfer. In contrast, IncLoRA exhibits stronger retention of earlier tasks, with improved backward transfer and final-task accuracy. This improvement stems from its incremental adapter-freezing mechanism, which isolates task-specific parameters after training. By preventing subsequent updates from overwriting earlier representations, IncLoRA preserves knowledge more effectively, though at the cost of reduced adaptability and parameter reuse.


% While the distinction between LoRA-only and full-model perturbations remains consistent with the observations in SeqLoRA, here we focus on the effect of different perturbation strategies within the LoRA subspace under IncLoRA. As shown in Table~\ref{tab:comparison_seqlora_inclora_vertical}, LoRA-Gaussian achieves the highest average accuracy ( AA = 80.65 in Order4) and a strong last-task accuracy (LA = 70.44), indicating robust generalization and stable final performance. It also delivers consistent forward transfer (FWT = 24.42), suggesting that unstructured noise helps explore flatter regions without disrupting adaptation.

% LoRA-Fisher, while slightly behind in overall accuracy, yields the best backward transfer in Order4 (ABWT = –10.50) and the highest last accuracy in Order6 (LA = 72.20), confirming its effectiveness in preserving past knowledge. This aligns with the design of Fisher-guided perturbations, which suppress noise on important parameters and promote selective flatness. LoRA-SAM achieves the highest FWT (24.62) in Order4, indicating improved immediate task adaptation, but its performance fluctuates across orders, and it underperforms in long-term stability metrics like ABWT and LA.

% These differences are further illustrated in Figure~\ref{fig:Order6_Inc_Process}. In the BWT plot (top), LoRA-Fisher consistently outperforms other strategies on earlier tasks, reflecting better retention. In contrast, SAM exhibits sharp drops, particularly in tasks 5 to 10, confirming its limited ability to maintain prior knowledge. The CTA plot (middle) shows that LoRA-Gaussian maintains the most stable task-wise accuracy throughout the sequence, while Fisher slightly improves in later tasks. In FWT (bottom), both Gaussian and SAM achieve higher early-task transfer, but Fisher gradually catches up and becomes more balanced toward the end.

% Taken together, these results show that Gaussian noise offers the most stable and general improvement across metrics, Fisher provides a more conservative trade-off favoring knowledge retention, and SAM boosts short-term adaptation but suffers from instability over longer task sequences. This confirms that perturbation strategies affect not only optimization geometry but also the stability–plasticity trade-off critical to continual learning.





% Despite this decoupling, the effectiveness of local flatness remains evident. As reported in Table~\ref{tab:comparison_seqlora_inclora_vertical}, LoRA-based perturbations lead to consistent improvements across all orders. LoRA-Fisher achieves the lowest ABWT (–10.50 in Order4) and the highest LA (72.20 in Order6), while LoRA-Gaussian and LoRA-SAM also enhance forward transfer (FWT), with LoRA-SAM reaching a peak of 24.62 in Order4. These results indicate that encouraging flatness within each individual adapter still benefits both retention and adaptation. In contrast, perturbing the full model leads to less predictable outcomes. Full-Gaussian, for example, worsens ABWT to –18.67 in Order5, and full-model methods display more variability across tasks. This suggests that global noise may inadvertently affect frozen components, which are particularly sensitive in IncLoRA’s sequential structure.

% To further illustrate these trends, Figure~\ref{fig:Order6_Inc_Process} presents the evolution of BWT, current task accuracy (CTA), and FWT over the 15-task sequence. LoRA-Fisher and LoRA-Gaussian maintain consistently higher performance across all three metrics, confirming their ability to support stable learning dynamics. In contrast, full-model perturbations show larger fluctuations and minimal gains, especially in the latter half of the sequence.

% These findings suggest that the benefits of flatness regularization are not exclusive to shared-adapter architectures. Even in IncLoRA, where task-wise decoupling naturally reduces interference, local flatness remains a key factor in enabling stability and transfer. The results reinforce the importance of applying regularization within the actively updated subspace rather than across the entire model.



% To explore this, we begin with a direct comparison between the two architectures using identical long-sequence settings. As shown in Table~\ref{tab:comparison_seqlora_inclora_vertical}, IncLoRA demonstrates stronger resistance to forgetting in the no-perturbation baseline, yielding higher backward transfer (ABWT = –12.23 vs. –14.45) and last accuracy (LA = 43.16 vs. 38.85 in Order6). In contrast, SeqLoRA achieves higher average accuracy across tasks. This performance difference can be attributed to architectural trade-offs: by continuously updating a shared adapter, SeqLoRA adapts more effectively to the current task, leading to stronger immediate performance. However, this comes at the cost of increased interference with prior knowledge, resulting in more severe forgetting. IncLoRA, by freezing each adapter after task completion, maintains a more faithful memory of previous tasks. While this limits parameter reuse and reduces adaptability to new data, it helps preserve past knowledge, which is reflected in its higher ABWT. Interestingly, forward transfer (FWT) is slightly better in IncLoRA, indicating that the base model remains a shared foundation across tasks. This suggests that generalization to unseen tasks relies more heavily on the fixed backbone, while the adapters capture task-specific variations. The comparable final performance between SeqLoRA and IncLoRA further implies that both approaches converge to solutions within a sufficiently expressive shared space. This observation reinforces the motivation to promote flatness during training, in order to guide optimization toward regions in parameter space that generalize well across multiple downstream tasks.












% To explore this, we begin with a direct comparison between the two architectures using identical long-sequence settings. As shown in Table~\ref{tab:comparison_seqlora_inclora_vertical}, IncLoRA demonstrates stronger resistance to forgetting in the no-perturbation baseline, yielding higher backward transfer (ABWT = –12.23 vs. –14.45) and final task accuracy (LA = 43.16 vs. 38.85 in Order6). At the same time, SeqLoRA consistently achieves higher average accuracy across tasks, benefiting from the reuse of shared representations. These results confirm that IncLoRA suppresses interference but may also limit cross-task generalization.




% \begin{table*}[hpbt]
% \centering
% \begin{tabular}{ll
%     cccc  % Order4
%     cccc  % Order5
%     cccc  % Order6
% }
% \toprule
%  & \multirow{2}{*}{\textbf{Method}}
%     & \multicolumn{4}{c}{\textbf{Order4}}
%     & \multicolumn{4}{c}{\textbf{Order5}}
%     & \multicolumn{4}{c}{\textbf{Order6}} \\
% \cmidrule(lr){3-6} \cmidrule(lr){7-10} \cmidrule(lr){11-14}
%  & & \textbf{ AA} & \textbf{FWT} & \textbf{BWT} & \textbf{LA}
%    & \textbf{ AA} & \textbf{FWT} & \textbf{BWT} & \textbf{LA}
%    & \textbf{ AA} & \textbf{FWT} & \textbf{BWT} & \textbf{LA} \\
% \midrule
% % \multirow{6}{*}{}
% %     & SeqLoRA No Noise
% %         & -     & —        & —        & — 
% %         &  —      &  —      & —        &— 
% %         &  —  & —    & —   & —  \\
% % \midrule
%     & IncLoRA No Noise
%         & 80.09      & 43.81       & -\textbf{10.67}       & 70.13
%         & 81.03      &  44.21     & -11.67       & 70.13
%         & 79.28  & 60.14   & -39.21  & 42.68 \\
% \midrule
%     & LoRA-SAM
%         & 79.91      & 45.57       & -12.01       & 68.69
%         & 80.11      & 41.94       & -12.41       & 68.52
%         & 78.83  & 58.85   & -9.62   & 69.85 \\
%     & LoRA-Gaussian
%         & \textbf{80.65}      & 45.36       & -10.94       & \textbf{70.44}
%         &  80.87     & 40.76       & -10.50       & 71.07
%         & \textbf{79.87}  & 59.31   & -13.23  & 67.52 \\
%     & LoRA-Fisher
%         & 80.63     & \textbf{45.76}       & -12.14       & 69.29
%         & 80.87      & \textbf{46.04}       & \textbf{-8.94}       & \textbf{72.52}
%         & 79.71  & 58.37   & \textbf{-8.04}  & \textbf{72.20} \\
% \midrule
%     & Full-Gaussian
%         & 80.43      & 44.90       & -11.73       & 69.48
%         & 80.67      & 42.30       & -12.31       & 69.17
%         & 79.50  & 58.61   & -16.15  & 64.43 \\
%     & Full-Fisher
%         & 80.05      & 46.33       & -11.59       & 69.24
%         & 80.73      & 40.61       & -11.29       & 70.20
%         & 79.05  & \textbf{61.44}   & -8.90   & 70.75 \\
% \bottomrule
% \end{tabular}
% \caption{Performance comparison of IncLoRA }
% \label{tab:inclora_orders_benchmark}
% \end{table*}



% \begin{table}[htbp]
% \centering
% \begin{tabular}{llcccccccc}
% \toprule
%  & \multirow{2}{*}{\textbf{Method}} 
%     & \multicolumn{4}{c}{\textbf{Long-Dataset}}  \\
% \cmidrule(lr){3-6} 
%  & & \textbf{ AA} & \textbf{FWT} & \textbf{BWT} & \textbf{LA} \\
% \midrule
% \multirow{6}{*}{} 
%     & IncLoRA No Noise      & 79.28  & 60.14  & -39.21  & 42.68   \\
%     \midrule
%     & LoRA-SAM      & 78.83  & 58.85 & -9.62  & 69.85     \\
%     & LoRA-Gaussian &\textbf{79.87} & 59.31  &{-13.23}  & 67.52     \\
%     & LoRA-Fisher   & 79.71  & 58.37 & \textbf{-8.04}  & \textbf{72.20}   \\
%     \midrule
%     & Full-Gaussian & 79.50  & 58.61 & -16.15  & 64.43      \\
%     & Full-Fisher   & {79.05} & \textbf{61.44} & -8.90   & 70.75    \\
% \bottomrule
% \end{tabular}
% \caption{Performance comparison of IncLoRA on Short Dataset. Metrics include  AA, BWT, FWT, and LA.}
% \label{tab:seqlora_cl_benchmark}
% \end{table}









% \begin{table*}[htbp]
% \centering
% \begin{tabular}{llcccccccccccc}
% \toprule
%  & \multirow{2}{*}{\textbf{Method}} 
%     & \multicolumn{4}{c}{\textbf{Short-Dataset}} 
%     & \multicolumn{4}{c}{\textbf{Long-Dataset}} 
%     & \multicolumn{4}{c}{\textbf{TRACE}} \\
% \cmidrule(lr){3-6} \cmidrule(lr){7-10} \cmidrule(lr){11-14}
%  & & \textbf{ AA} & \textbf{FWT} & \textbf{BWT} & \textbf{LA}
%    & \textbf{ AA} & \textbf{FWT} & \textbf{BWT} & \textbf{LA}
%    & \textbf{ AA} & \textbf{FWT} & \textbf{BWT} & \textbf{LA} \\
% \midrule
% \multirow{6}{*}{} 
%     & No Noise      & 80.8  & 49.0  & -4.33  & 77.50  & 81.00  & 60.78  & -45.16 & 38.85  & - & - & - & - \\
%     & LoRA-SAM      & 79.4  & 49.86 & -2.26  & 77.7   & 79.64  & 60.54  & -65.93 & 18.10  & - & - & - & - \\
%     & LoRA-Gaussian & 79.75 & 50.2  & -0.86  & 79.1   & 81.42  & 56.44  & -53.14 & 31.82  & - & - & - & - \\
%     & LoRA-Fisher   & 80.7  & 49.86 & -2.26  & 78.95  & 80.8   & 59.13  & -38.28 & 45.06  & - & - & - & - \\
%     & Full-Gaussian & 80.75 & 49.06 & -4.86  & 77.1   & 81.42  & 58.4   & -57.47 & 27.78  & - & - & - & - \\
%     & Full-Fisher   & 80.95 & 48.86 & -3.2   & 78.55  & 81.37  & 58.67  & -53.14 & 31.54  & - & - & - & - \\
% \bottomrule
% \end{tabular}
% \caption{Performance comparison of SeqLoRA on Short-, Long-Dataset and TRACE benchmarks. Metrics include  AA, BWT, FWT, and LA.}
% \label{tab:seqlora_cl_benchmark}
% \end{table*}






\subsection{Flatness in OLoRA: Orthogonal Adaptation with Stability}



\begin{table}[htbp]
\centering
\resizebox{.95\linewidth}{!}{
\begin{tabular}{cc|cccc}
\toprule
& & \multicolumn{4}{c}{\textbf{TRACE Benchmark}} \\
\cmidrule(lr){3-6}
\textbf{Base} & \textbf{Method} & ABWT &  AA & AFWT & LA \\
\midrule
\multirow{6}{*}{\rotatebox[origin=c]{90}{OLoRA}} 
& No Noise        & -13.04     & 44.76     & 14.03      & 27.78     \\
% \cmidrule(lr){2-6}
& LoRA-SAM        & \textbf{-3.84}     & 42.28     & 13.53      & 39.60     \\
& LoRA-Gaussian   &-12.13  &45.08  &12.70  & 32.94 \\
& LoRA-Fisher     &-15.85  &46.00  &12.75  & 28.83 \\
% \cmidrule(lr){2-6}
& Full-Gaussian   &-8.07  &45.61  &13.87  &\textbf{41.50}  \\
& Full-Fisher     &-10.25  &\textbf{46.11}  &\textbf{14.08}  &38.29 \\
\bottomrule
\end{tabular}
}
\caption{Performance comparison of {OLoRA} under the TRACE benchmark.}
\label{tab:seqlora_trace_dualcol}
\end{table}


We now turn to OLoRA, an extension of IncLoRA that incorporates orthogonality constraints between task-specific LoRA adapters. This structure enforces representational disentanglement by ensuring that each new adapter lies in a subspace orthogonal to all previous ones, thereby reducing interference. While this constraint improves knowledge retention, it also limits the expressiveness of each subspace, potentially affecting learning on new tasks.


As shown in Table~\ref{tab:comparison_seqlora_inclora_vertical}, OLoRA without perturbation achieves significantly better ABWT and AFWT than both SeqLoRA and IncLoRA, confirming the benefit of orthogonal subspace partitioning in mitigating interference. However, this improvement is accompanied by a noticeable drop in average accuracy (AA), indicating a trade-off between forgetting reduction and task-specific learning capacity. This degradation suggests that strict orthogonality, while effective in preventing interference, may constrain the expressiveness of newly allocated subspaces.

To address this limitation, flatness-promoting perturbations provide a complementary regularization mechanism. For example, LoRA-Gaussian and LoRA-Fisher yield further gains in ABWT and LA relative to the no-noise baseline. These improvements imply that encouraging flatter minima within each orthogonal subspace can compensate for reduced flexibility by improving robustness and generalization without violating structural constraints.


This interaction becomes more pronounced under the TRACE benchmark (Table~\ref{tab:seqlora_trace_dualcol}), which introduces greater task diversity and domain shift. Despite the increased complexity, methods such as LoRA-SAM and LoRA-Gaussian still yield consistent improvements in forgetting and retention. These results confirm that orthogonality and flatness address distinct challenges: the former prevents cross-task interference by structural separation, while the latter promotes stability and transfer within constrained spaces. Their synergy enables more robust continual learning, even when model capacity is explicitly partitioned.



% To further investigate the role of flatness in structurally constrained continual learning, we analyze OLoRA, which enforces orthogonality between task-specific LoRA subspaces. This design explicitly eliminates parameter overlap across tasks, aiming to reduce interference. However, while orthogonality constrains the representational space geometrically, it does not inherently control the optimization trajectory or local loss curvature. We therefore examine whether promoting flatness can complement orthogonality to enhance stability and generalization.


% The OLoRA framework introduces orthogonality constraints across task-specific LoRA subspaces, which effectively mitigate interference by enforcing parameter disentanglement. As shown in Table~\ref{tab:comparison_seqlora_inclora_vertical}, this leads to substantial improvements in ABWT and AFWT compared to SeqLoRA and IncLoRA. However, the average accuracy (AA) declines, suggesting that orthogonalization, while reducing forgetting, limits subspace flexibility and thus impairs current-task learning. 

% Flatness-oriented perturbations address this limitation by promoting broader, low-curvature solutions within each constrained subspace. This helps stabilize optimization under reduced capacity, guiding updates toward more robust minima without violating the orthogonality constraint. Empirically, applying flatness (e.g., LoRA-Gaussian or LoRA-Fisher) further improves ABWT and LA while preserving forward transfer, confirming its compensatory effect.

% Notably, this conclusion is further validated on the TRACE benchmark (Table~\ref{tab:seqlora_trace_dualcol}), where task diversity is high and structural constraints are more likely to conflict with generalization demands. Even under these conditions, LoRA-SAM and LoRA-Gaussian yield substantial gains in ABWT and LA relative to the no-noise baseline, confirming that flatness serves as a critical regularizer when subspace separation alone is insufficient.



% Our results suggest that orthogonality and flatness target distinct yet complementary sources of forgetting. As shown in Table~\ref{tab:comparison_seqlora_inclora_vertical}, the unperturbed OLoRA baseline already achieves significantly improved ABWT compared to both SeqLoRA and IncLoRA. This confirms the effectiveness of orthogonal subspace partitioning in mitigating interference. However, further applying flatness-promoting noise yields consistent additional gains.  However, further applying flatness-promoting noise yields consistent additional gains. These improvements are meaningful given the reduced interference already achieved through orthogonalization.


% This pattern highlights a critical interaction: while orthogonality prevents parameter overlap, it may restrict the expressive capacity of each new subspace, especially as the number of tasks increases. In such a setting, promoting flatness within each subspace helps compensate by guiding optimization toward more stable, transferable minima within the constrained manifold. Flatness reduces the sensitivity of each adapter to noisy gradients from new tasks, thereby improving forward transfer without undermining the non-overlapping structure imposed by orthogonality.





% We extend our investigation to OLoRA, a variant of IncLoRA that incorporates orthogonality constraints across task-specific LoRA subspaces. In this framework, each newly introduced adapter is optimized under the constraint that its parameter subspace remains orthogonal to those of all previously learned tasks. This mechanism aims to reduce representational interference and promote explicit parameter disentanglement. The central question in this section is whether flatness-promoting perturbations retain their efficacy when applied within such structurally regularized, orthogonal subspaces.


% To evaluate the impact of flatness under orthogonality constraints, we report performance across three task orders on the long-sequence benchmark (Table~\ref{tab:comparison_seqlora_inclora_vertical}) and the more heterogeneous TRACE benchmark (Table~\ref{tab:seqlora_trace_dualcol}). In all settings, OLoRA without perturbation achieves notably stronger retention than prior baselines, confirming that structural orthogonality alone substantially mitigates forgetting. However, our results also demonstrate that flatness-aware perturbations remain beneficial even under this constrained setting.


% \begin{table}[htbp]
% \centering
% \begin{tabular}{llcccccccc}
% \toprule
%  & \multirow{2}{*}{\textbf{Method}} 
%     & \multicolumn{4}{c}{\textbf{TRACE Benchmark}}  \\
% \cmidrule(lr){3-6} 
%  & & \textbf{ AA} & \textbf{FWT} & \textbf{BWT} & \textbf{LA} \\
% \midrule
% \multirow{6}{*}{} 
%     & No Noise      & -  & -  & -  & -   \\
%     \midrule
%     & LoRA-SAM      & -  & - & -  & -     \\
%     & LoRA-Gaussian & 45.08 & {12.00}  & {-13.87}  & {32.94}     \\
%     & LoRA-Fisher   & 46.00  & 12.75 & -19.62  & 28.83   \\
%     \midrule
%     & Full-Gaussian & 45.61 & 16.96 & -4.69  & 41.50      \\
%     & Full-Fisher   & {46.11} & 17.35 & -8.94   & 38.29    \\
% \bottomrule
% \end{tabular}
% \caption{Performance comparison of SeqLoRA on Short Dataset. Metrics include  AA, BWT, FWT, and LA.}
% \label{tab:olora_cl_benchmark}
% \end{table}















% \begin{table*}[htbp]
% \centering
% \begin{tabular}{llcccccccccccc}
% \toprule
%  & \multirow{2}{*}{\textbf{Method}} 
%     & \multicolumn{4}{c}{\textbf{Short-Dataset}} 
%     & \multicolumn{4}{c}{\textbf{Long-Dataset}} 
%     & \multicolumn{4}{c}{\textbf{TRACE}} \\
% \cmidrule(lr){3-6} \cmidrule(lr){7-10} \cmidrule(lr){11-14}
%  & & \textbf{ AA} & \textbf{FWT} & \textbf{BWT} & \textbf{LA}
%    & \textbf{ AA} & \textbf{FWT} & \textbf{BWT} & \textbf{LA}
%    & \textbf{ AA} & \textbf{FWT} & \textbf{BWT} & \textbf{LA} \\
% \midrule
% \multirow{6}{*}{} 
%     & No Noise      & 80.8  & 49.0  & -4.33  & 77.50  & 81.00  & 60.78  & -45.16 & 38.85  & - & - & - & - \\
%     & LoRA-SAM      & 79.4  & 49.86 & -2.26  & 77.7   & 79.64  & 60.54  & -65.93 & 18.10  & - & - & - & - \\
%     & LoRA-Gaussian & 79.75 & 50.2  & -0.86  & 79.1   & 81.42  & 56.44  & -53.14 & 31.82  & - & - & - & - \\
%     & LoRA-Fisher   & 80.7  & 49.86 & -2.26  & 78.95  & 80.8   & 59.13  & -38.28 & 45.06  & - & - & - & - \\
%     & Full-Gaussian & 80.75 & 49.06 & -4.86  & 77.1   & 81.42  & 58.4   & -57.47 & 27.78  & - & - & - & - \\
%     & Full-Fisher   & 80.95 & 48.86 & -3.2   & 78.55  & 81.37  & 58.67  & -53.14 & 31.54  & - & - & - & - \\
% \bottomrule
% \end{tabular}
% \caption{Performance comparison of SeqLoRA on Short-, Long-Dataset and TRACE benchmarks. Metrics include  AA, BWT, FWT, and LA.}
% \label{tab:seqlora_cl_benchmark}
% \end{table*}





\subsection{Ablation Study}
\subsubsection{Hyperparameter Sensitivity }
\begin{figure}[ht]
  \centering
  \includegraphics[width=0.8\linewidth]{images/Paper_Order3_accuracy_lr_epoch_compare-3methods.pdf}
  \caption{Performance comparison of SeqFT, IncLoRA, and FT-LoRA, where FT-LoRA denotes training a LoRA on the finetuned model. Results are reported on Order3 with varying learning rates (1e-3, 1e-4, 1e-5) and training epochs (1, 3, 5, and 10), while the sample size is fixed at 5000.}
  \label{fig result: epoch VS Learning rate}
\end{figure}
\begin{figure}[h]
  \centering
  \includegraphics[width=0.8\linewidth]{images/Paper_Order3_accuracy_samples_epoch_compare-3methods.pdf}
  \caption{Performance comparison of SeqFT, IncLoRA, and
FT-LoRA on Order3, varying sample size (500, 1000, 3000, 5000) and training epochs (1, 3, 5, 10), with learning rate fixed at 1e-4 for IncLoRA and 1e-5 for SeqFT.}
  \label{fig result: epoch VS sample sizes}
\end{figure}
\paragraph{Trade-offs among Learning Rate, Sample Size, and Training Epochs} We investigate how key hyperparameters, including learning rate, training sample size, and number of epochs, affect the trade-off between task adaptation and knowledge retention in continual learning.
As shown in Fig.\ref{fig result: epoch VS Learning rate} and Fig.\ref{fig result: epoch VS sample sizes}, we observe that higher learning rates and excessive training epochs accelerate forgetting and reduce generalization to previous and future tasks, despite yielding higher accuracy on the current task.Increasing the training set size improves immediate task performance, but also exacerbates catastrophic forgetting. LoRA-based methods (IncLoRA, FT-LoRA) achieve performance comparable to full finetuning(SeqFT) while retaining parameter efficiency.
Based on these findings, we select a baseline hyperparameter configuration that carefully balances current-task optimization and knowledge retention, and use it consistently for all subsequent experiments to ensure fairness and reproducibility.




\begin{figure}[h]
  \centering
  \includegraphics[width=1\linewidth]{images/Paper_SAM_RWP_Rho_Paper3.pdf}
  \caption{Performance of IncLoRA under varying perturbation strengths and perturbation scopes. The figure reports accuracy on past (Yahoo), current (Amazon), and future tasks (AGNews and DBpedia), comparing three flatness-promoting strategies: LoRA-SAM, LoRA-Gaussian, and Full-Gaussian.}
  \label{fig:sam_rho_tradeoff}
\end{figure}
\paragraph{Trade-offs Perturbation Strength in Flatness-Promoting Methods} We further examine the effect of perturbation strength and the scope of flatness regularization on continual learning performance.
As illustrated in Fig.~\ref{fig:sam_rho_tradeoff}, different perturbation strength lead to  trade-offs. On the one hand, a small perturbation may be insufficient to effectively promote generalization. On the other hand, an excessively large perturbation can negatively affect the model's current task performance. These observations inform the choice of perturbation strength in the method comparisons.


\section{Conclusion}


% As shown in Fig.~\ref{fig:sam_rho_tradeoff}, varying $\rho$ produces trade-offs: while small values may insufficiently promote flatness, excessive perturbation—especially when applied only to the LoRA branch—can harm stability and generalization, particularly for previous and future tasks. These results  guide the selection of $\rho$ for the subsequent method comparison.

\bibliography{aaai2026}



\appendix




% \input{ReproducibilityChecklist_clean}
\input{ReproducibilityChecklist}





\onecolumn
\section{Appendix}

\subsection{A: Proof of the theorem 2}

\begin{theorem}[PAC--Bayes Bound with flatness for  Continual Learning]
\label{thm:general-pacbayes-eng}
Under the hierarchical Bayesian continual‐learning framework for $n$ sequential tasks, for any hyperposterior $\mathcal Q$ and any family of task‐wise posteriors $\{Q_t\}_{t=1}^n$, with probability at least $1-\delta$ over all training samples,
\[
L(\mathcal Q_{1:n})
\leq
\frac{1}{n}\sum_{t=1}^n\mathbb{E}_{P_t\sim\mathcal Q_t}\Bigl[
   \widehat L_{S_t}(W_t)+\mathrm{E\text{-}Sh}_t(W_t,\Sigma_Q)
\Bigr]
+\frac{1}{n(\bar m_n-1)}
\sqrt{\frac{KL^{\mathrm{task}}_n + KL^{\mathrm{hyper}}_n + \ln\bigl(\bar m_n/\delta\bigr)}{2(\bar m_n-1)}}\,,
\]
where
\[
KL^{\mathrm{task}}_n = KL\bigl(Q_{1:n}\|P_{1:n}\bigr),
\quad
KL^{\mathrm{hyper}}_n = KL(\mathcal Q\|\mathcal P),
\quad
\bar m_n = \frac{n}{\sum_{t=1}^n 1/m_t}.
\]
\end{theorem}

\begin{proof}
We prove the bound by induction on the number of tasks~$n$.  Let $[a,b]$ be the range of the bounded loss.

\medskip\noindent\textbf{Base case ($n=1$).}  
For a single task with dataset $S_1$ of size $m_1$, classical PAC--Bayes (McAllester, 2003) states: for any prior $P_1$ and any posterior $Q_1$ and with probability at least $1-\delta_1$,

\begin{equation}
    \mathbb{E}_{f\sim Q_1}[L(f)]
\le
\mathbb{E}_{f\sim Q_1}[\widehat L_{S_1}(f)]
+ \sqrt{\frac{KL(Q_1\|P_1)+\ln(m_1/\delta_1)}{2(m_1-1)}}.
\end{equation}


\medskip\noindent\textbf{Hierarchical extension with hyperdistributions} 
Introduce a hyperprior $\mathcal{P}_1$ and hyperposterior $\mathcal{Q}_1$  over task priors $P_1$. Taking expectation over $P_1 \sim \mathcal{Q}_1$ on both sides, we use Jensen's inequality to get:

\begin{equation}
\mathbb{E}_{P_1 \sim \mathcal{Q}_1}\mathbb{E}_{f \sim Q_1(P_1)}[L(f)]
\le
\mathbb{E}_{P_1 \sim \mathcal{Q}_1}\mathbb{E}_{f \sim Q_1(P_1)}[\widehat{L}_{S_1}(f)]
+\sqrt{\frac{\mathbb{E}_{P_1 \sim \mathcal{Q}_1}[KL(Q_1\|P_1)] + KL(\mathcal{Q}_1\|\mathcal{P}_1) + \ln(m_1/\delta_1)}{2(m_1-1)}}.
\end{equation}

\medskip\noindent\textbf{Bridging distributional expectations and point estimates} While the classical PAC–Bayes bound features the expected disribution of the model $\mathbb{E}_{f\sim Q_1}[\widehat L_{S_1}(f)]$, in practice one uses a single parameter $\hat f$ (e.g.\ the posterior‐mean or MAP) for prediction.  We therefore decompose
\[
\mathbb{E}_{f\sim Q_1}[\widehat L_{S_1}(f)]
=
\underbrace{\widehat L_{S_1}(\hat f)}_{\text{empirical loss at the reference}}
\;+\;
\underbrace{
  \mathbb{E}_{f\sim Q_1}[\widehat L_{S_1}(f)] - \widehat L_{S_1}(\hat f)
}_{\displaystyle \mathrm{E\text{-}Sh}_1},
\]
where we define the {expected sharpness}:
\[
\mathrm{E\text{-}Sh}_1(f,Q_1)
:=\mathbb{E}_{f\sim Q_1}[\widehat L_{S_1}(f)] - \widehat L_{S_1}(\hat f).
\]

Concretely, choose $\hat f$ to minimize $\mathbb{E}_{f\sim Q_1}[\widehat L_{S_1}(f)]$.  Then
\[
\mathbb{E}_{f\sim Q_1}[\widehat L_{S_1}(f)]
= \widehat L_{S_1}(\hat f) + \mathrm{E\text{-}Sh}_1,
\]



and substituting into the hierarchical bound gives the final base‐case inequality
\[
\mathbb{E}_{P_1\sim\mathcal Q_1}\mathbb{E}_{f\sim Q_1(P_1)}[L(f)]
\;\le\;
\mathbb{E}_{P_1}\Bigl[\widehat L_{S_1}(\hat f) + \mathrm{E\text{-}Sh}_1\Bigr]
\;+\;
\sqrt{\frac{KL^{\mathrm{task}}_1 + KL^{\mathrm{hyper}}_1 + \ln(m_1/\delta_1)}{2(m_1-1)}}.
\]
 where $KL^{\text{task}}_1 = \mathbb{E}_{P_1 \sim \mathcal{Q}_1}[KL(Q_1\|P_1)]$ and $KL^{\text{hyper}}_1 = KL(\mathcal{Q}_1\|\mathcal{P}_1)$. With $\delta_1 = \delta/2$, the confidence is $1-\delta/2$, satisfying the base case.
Setting $\delta_1=\delta/2$ yields the bound at confidence $1-\delta/2$.  This introduction of $\mathrm{E\text{-}Sh}_1$ is both theoretically natural measuring posterior‐induced “flatness” and practically aligned with using a single point estimate.  













% Define the expected sharpness for task $t$ as:
% \begin{equation}\label{eq:esh}
% \mathrm{E\text{-}Sh}_t(W_t, \mathcal{Q}) 
% := \mathbb{E}_{W_t \sim Q_t}[\widehat{L}_{S_t}(W_t)] - \widehat{L}_{S_t}(\hat{W}_t),
% \end{equation}
% where $\hat{W}_t$ is the mode of $Q_t$ (i.e., $\arg\min_{W} \mathbb{E}_{W_t \sim Q_t}[\widehat{L}_{S_t}(W_t)]$).

% By definition, $\mathbb{E}[\widehat{L}] = \widehat{L}(\hat{W}_t) + \mathrm{E\text{-}Sh}_t$, so substituting gives:
% \[
% \mathbb{E}_{P_1 \sim \mathcal{Q}_1}\mathbb{E}_{f \sim Q_1(P_1)}[L(f)]
% \le
% \mathbb{E}_{P_1 \sim \mathcal{Q}_1}\left[\widehat{L}_{S_1}(\hat{W}_1) + \mathrm{E\text{-}Sh}_1\right]
% + (b-a)\sqrt{\frac{KL^{\text{task}}_1 + KL^{\text{hyper}}_1 + \ln(m_1/\delta_1)}{2(m_1-1)}},
% \]
% where $KL^{\text{task}}_1 = \mathbb{E}_{P_1 \sim \mathcal{Q}_1}[KL(Q_1\|P_1)]$ and $KL^{\text{hyper}}_1 = KL(\mathcal{Q}_1\|\mathcal{P}_1)$. With $\delta_1 = \delta/2$, the confidence is $1-\delta/2$, satisfying the base case.



% Introduce a hyperprior $\mathcal P_1$ over $P_1$ and hyperposterior $\mathcal Q_1$, and take expectation over $P_1\sim\mathcal Q_1$.  Then with probability at least $1-\delta_1$,
% \[
% \mathbb{E}_{P_1\sim\mathcal Q_1}\mathbb{E}_{f\sim Q_1(P_1)}[L(f)]
% \le
% \mathbb{E}_{P_1}\mathbb{E}_{f}[\widehat L_{S_1}(f)]
% + \sqrt{\frac{\mathbb{E}_{P_1}[KL(Q_1\|P_1)] + KL(\mathcal Q_1\|\mathcal P_1) + \ln(m_1/\delta_1)}{2(m_1-1)}}.
% \]
% Define
% \[
% KL^{\mathrm{task}}_1 = \mathbb{E}_{P_1}[KL(Q_1\|P_1)],
% \quad
% KL^{\mathrm{hyper}}_1 = KL(\mathcal Q_1\|\mathcal P_1),
% \quad
% \bar m_1=m_1,
% \quad
% \delta_1=\delta/2.
% \]
% Absorb $(b-a)$ into the definition of expected sharpness:
% \[
% \mathrm{E\text{-}Sh}_1 = \mathbb{E}_{f\sim Q_1}[\widehat L_{S_1}(f)] - \widehat L_{S_1}(\hat f_1),
% \]
% and note
% \[
% \mathbb{E}_{f\sim Q_1}[\widehat L_{S_1}(f)]
% = \widehat L_{S_1}(\hat f_1) + \mathrm{E\text{-}Sh}_1.
% \]
% Thus the stated bound holds for $n=1$ with confidence $1-\delta/2$.


\medskip\noindent\textbf{Inductive step (\(n\to n+1\) tasks).}

Assume the bound holds for $n$ tasks with confidence $1-\delta/2$:
\begin{equation}
\label{eq: combine0}
    L(\mathcal{Q}_{1:n})
\le
\frac{1}{n}\sum_{t=1}^n \mathbb{E}_{P_t \sim \mathcal{Q}_t}\left[\widehat{L}_{S_t}(\hat{W}_t) + \mathrm{E\text{-}Sh}_t\right]
+\sqrt{\frac{KL^{\text{task}}_n + KL^{\text{hyper}}_n + \ln(\bar{m}_n/(\delta/2))}{2n^2(\bar{m}_n - 1)}}.
\end{equation}
\label{eq: combine1}
For task $n+1$, apply the single-task bound (confidence $1-\delta/2$):
\begin{equation}
\mathbb{E}_{P_{n+1} \sim \mathcal{Q}_{n+1}}\mathbb{E}_{f \sim Q_{n+1}}[L(f)]
\le
\mathbb{E}_{P_{n+1} \sim \mathcal{Q}_{n+1}}\left[\widehat{L}_{S_{n+1}}(\hat{W}_{n+1}) + \mathrm{E\text{-}Sh}_{n+1}\right]
+ \sqrt{\frac{KL^{\text{task}}_{n+1,\text{single}} + KL^{\text{hyper}}_{n+1,\text{single}} + \ln(m_{n+1}/(\delta/2))}{2(m_{n+1} - 1)}}.
\end{equation}
By the union bound, both  eq(\ref{eq: combine0}) and eq(\ref{eq: combine1}) hold with probability $1-\delta$. We combine them via $\alpha = \frac{n}{n+1}$ and $1-\alpha = \frac{1}{n+1}$.


% ---
% Assume that for \(n\) tasks the bound holds with confidence \(1-\tfrac\delta2\):
% \[
% L(\mathcal Q_{1:n})
% \;\le\;
% \frac{1}{n}\sum_{t=1}^n \mathbb{E}_{P_t\sim\mathcal Q_t}\bigl[\widehat L_{S_t}(W_t)+\mathrm{E\text{-}Sh}_t\bigr]
% \;+\;
% \frac{1}{n(\bar m_n-1)}
% \sqrt{\frac{KL_n^{\mathrm{task}} + KL_n^{\mathrm{hyper}} + \ln\bigl(\bar m_n/(\delta/2)\bigr)}{2(\bar m_n-1)}}.
% \]
% Similarly, applying the single‐task hierarchical PAC–Bayes bound to task \(n+1\) (also at confidence \(1-\tfrac\delta2\)), we have for any \(P_{n+1}\sim\mathcal Q_{n+1}\):
% \[
% \mathbb{E}_{f\sim Q_{n+1}}[L(f)]
% \;\le\;
% \mathbb{E}_{P_{n+1}}\bigl[\widehat L_{S_{n+1}}(W_{n+1})+\mathrm{E\text{-}Sh}_{n+1}\bigr]
% \;+\;
% \frac{1}{\bar m_{n+1}-1}
% \sqrt{\frac{KL(Q_{n+1}\|P_{n+1}) + \ln\bigl(m_{n+1}/(\delta/2)\bigr)}{2(\bar m_{n+1}-1)}},
% \]
% where
% \(\bar m_{n+1}=(n+1)/\sum_{t=1}^{n+1}(1/m_t)\).
% By the union bound, both high‐probability events hold simultaneously with probability at least
% \[
% 1-\frac\delta2-\frac\delta2
% \;=\;
% 1-\delta.
% \]

% We now blend the two inequalities via a convex combination.  Set
% \[
% \alpha=\frac{n}{n+1},
% \quad
% 1-\alpha=\frac{1}{n+1}.
% \]
% Multiply the \(n\)-task inequality by \(\alpha\) and the \((n+1)\)th‐task inequality by \(1-\alpha\).

1. Combination of empirical‐risk and sharpness terms 

The average generalization loss for $n+1$ tasks is:
\[
L(\mathcal{Q}_{1:n+1}) = \alpha \cdot L(\mathcal{Q}_{1:n}) + (1-\alpha) \cdot \mathbb{E}\left[L_{n+1}\right],
\]
and the empirical terms combine as:
\begin{align*}
&\alpha \cdot \frac{1}{n}\sum_{t=1}^n \mathbb{E}[\cdot] + (1-\alpha) \cdot \mathbb{E}[\cdot] \\
&\quad=\alpha\;\frac{1}{n}\sum_{t=1}^n\bigl(\widehat L_{S_t}+\mathrm{E\!-\!Sh}_t\bigr)
\;+\;(1-\alpha)\,\bigl(\widehat L_{S_{n+1}}+\mathrm{E\!-\!Sh}_{n+1}\bigr)\\
&\quad=\;\frac{n}{n+1}\,\frac{1}{n}\sum_{t=1}^n\!\bigl(\widehat L_{S_t}+\mathrm{E\!-\!Sh}_t\bigr)
\;+\;\frac{1}{n+1}\bigl(\widehat L_{S_{n+1}}+\mathrm{E\!-\!Sh}_{n+1}\bigr)\\
&\quad=\;\frac{1}{n+1}\Bigl(\sum_{t=1}^n(\widehat L_{S_t}+\mathrm{E\!-\!Sh}_t)
+(\widehat L_{S_{n+1}}+\mathrm{E\!-\!Sh}_{n+1})\Bigr)\\
&\quad=\;\frac{1}{n+1}\sum_{t=1}^{n+1}\bigl(\widehat L_{S_t}+\mathrm{E\!-\!Sh}_t\bigr).
\end{align*}

2. Chain-rule decomposition of KL divergences (task and hyper levels)  


Let \(Q^{\mathrm{joint}}_{1:n+1}\) and \(P^{\mathrm{joint}}_{1:n+1}\) denote the full joint posterior and prior (over both parameters and hyperparameters), which factorize in a chain‐structured manner.  Then by the general chain rule for KL ,
\[
KL\bigl(Q^{\mathrm{joint}}_{1:n+1}\big\|P^{\mathrm{joint}}_{1:n+1}\bigr)
=KL(Q_1\|P_1)
+\sum_{t=2}^{n+1}
\mathbb{E}_{Q^{\mathrm{joint}}_{1:t-1}}
\Bigl[
  KL\bigl(Q_t\|\;P_t\bigl\|\mid P_{1:t-1}\bigr)
\Bigr].
\]
Marginalizing out hyperparameters then defines the parameter‐level divergence
\[
KL^{\mathrm{task}}_{n+1}
=KL\bigl(Q_{1:n+1}(\Theta)\big\|P_{1:n+1}(\Theta)\bigr),
\]
which in the special case of independent factors reduces to \(\sum_{t=1}^{n+1}KL(Q_t\|P_t)\).

An analogous expansion holds for the hyperparameter KL, yielding
\[
KL^{\mathrm{hyper}}_{n+1}
=KL(\mathcal Q_1\|\mathcal P_1)
+\sum_{t=2}^{n+1}
\mathbb{E}_{\mathcal Q_{1:t-1}}
\bigl[\,KL(\mathcal Q_t\|\mathcal P_t(\cdot\mid\mathcal Q_{1:t-1}))\bigr].
\]

3. Combination of complexity terms  
We must verify that
\[
\alpha\,\frac{1}{n(\bar m_n-1)}
\;+\;
(1-\alpha)\,\frac{1}{\bar m_{n+1}-1}
\;=\;\frac{1}{(n+1)(\bar m_{n+1}-1)}.
\]
From the definition
\(\bar m_{n+1}=(n+1)/\sum_{t=1}^{n+1}(1/m_t)\)
it follows that
\[
\bar m_{n+1}-1
=\frac{n(\bar m_n-1)+(m_{n+1}-1)}{\sum_{t=1}^{n+1}(1/m_t)}
\;\Longrightarrow\;
\frac{1}{\bar m_{n+1}-1}
=\frac{\sum_{t=1}^{n+1}(1/m_t)}{n(\bar m_n-1)+(m_{n+1}-1)}.
\]
Substituting this identity into the left‐hand side immediately yields the claimed equality.



Putting these three pieces together recovers exactly the stated bound for \(n+1\) tasks:
\[
L(\mathcal Q_{1:n+1})
\;\le\;
\frac{1}{n+1}
\sum_{t=1}^{n+1}\mathbb{E}_{P_t\sim\mathcal Q_t}\bigl[\widehat L_{S_t}+\mathrm{E\text{-}Sh}_t\bigr]
\;+\;
\frac{1}{(n+1)(\bar m_{n+1}-1)}
\sqrt{\frac{KL^{\mathrm{task}}_{n+1}+KL^{\mathrm{hyper}}_{n+1}+\ln(\bar m_{n+1}/\delta)}{2(\bar m_{n+1}-1)}}.
\]
This completes the inductive argument.  \qed



\end{proof}




\subsection{B: Proof of the theorem 3}
\textbf{Theorem: Multi-Task PAC-Bayes Generalization Bound for Incremental LoRA with Chain-Dependent Hyperdistributions}


Consider a continual learning scenario with $n$ sequential tasks under the Incremental LoRA framework, subject to the following conditions:

1. Orthogonality of Parameter Spaces: For each task $t$, the LoRA parameter block $\Delta\mathbf{W}_t$ is constrained to be orthogonal to all previous parameter blocks, i.e., for any $i < t$, $A_i^\top A_t = 0$, where $A_i, A_t$ denote the adapter matrices of tasks $i$ and $t$, respectively.

2. Gaussian Posterior Assumption: The posterior parameter distribution for task $t$ is Gaussian, $Q_t = \mathcal{N}(\hat{\boldsymbol{\theta}}_t, \rho_t^2 \mathbf{I})$, where $\hat{\boldsymbol{\theta}}_t$ is the optimal parameter for task $t$, and $\rho_t^2$ reflects the posterior uncertainty.

3. Chain-Structured Hyperdistribution: The hyperprior for task $t$, $\mathcal{P}_t$, is conditioned on the cumulative hyperposterior of all previous tasks, i.e., the joint hyperprior is
   $$
   \mathcal{P}_{1:n} = \mathcal{P}_1 \otimes \mathcal{P}_2(\cdot|\mathcal{Q}_1) \otimes \cdots \otimes \mathcal{P}_n(\cdot|\mathcal{Q}_{1:n-1}),
   $$ where $\mathcal{Q}_{1:t-1}$ denotes the joint hyperposterior for tasks $1$ through $t-1$.

4. Sharpness Definition: The expected sharpness for task $t$ is defined as
% \begin{align}
% &\mathrm{Sharpness}_t(\mathbf{W}_{t-1} + \hat{\Delta\mathbf{W}}_t,\, \rho_t^2) \notag\\
% & = \mathbb{E}_{\epsilon \sim \mathcal{N}(0,\, \rho_t^2 \mathbf{I})}
% \left[ \widehat{L}_{S_t}(\mathbf{W}_{t-1} + \hat{\Delta\mathbf{W}}_t + \epsilon) \right] \notag\\
% &\quad - \widehat{L}_{S_t}(\mathbf{W}_{t-1} + \hat{\Delta\mathbf{W}}_t) \notag
% \end{align}
\begin{align}
\mathrm{Sharpness}_t(\mathbf{W}_{t-1} + \hat{\Delta\mathbf{W}}_t,\, \rho_t^2)  = \mathbb{E}_{\epsilon \sim \mathcal{N}(0,\, \rho_t^2 \mathbf{I})}
\left[ \widehat{L}_{S_t}(\mathbf{W}_{t-1} + \hat{\Delta\mathbf{W}}_t + \epsilon) \right] \notag\quad - \widehat{L}_{S_t}(\mathbf{W}_{t-1} + \hat{\Delta\mathbf{W}}_t) \notag
\end{align}
quantifying the increase in empirical loss due to parameter perturbation.

\begin{theorem}[PAC-Bayes Generalization Bound for Incremental Multi-Task LoRA]
With probability at least $1-\delta$, the generalization risk $L(\mathcal{Q}_{1:n})$ for $n$ sequential tasks under the incremental LoRA framework satisfies
% \begin{align}
% L(\mathcal{Q}_{1:n}) &\leq \frac{1}{n} \sum_{t=1}^n 
%     \mathbb{E}_{P_{1:n} \sim \mathcal{Q}_{1:n}}
%     \Big[
%         \widehat{L}_{S_t}(\mathbf{W}_{t-1} + \hat{\Delta\mathbf{W}}_t)
%     \notag \\
% &\quad +\; \mathrm{Sharpness}_t(\mathbf{W}_{t-1} + \hat{\Delta\mathbf{W}}_t,\, \rho_t^2)
%     \Big] 
%     \notag \\
% & + (b-a) \left\{ 
%     \frac{KL_n^{\text{task}} + KL_n^{\text{hyper}} + \log(m/\delta)}
%          {2(\bar{m}_n - 1)}
% \right\}^{1/2}
% \label{eq:lora-pac-bayes-bound}
% \end{align}

\begin{align}
L(\mathcal{Q}_{1:n}) &\leq \frac{1}{n} \sum_{t=1}^n 
    \mathbb{E}_{P_{1:n} \sim \mathcal{Q}_{1:n}}
    \Big[
        \widehat{L}_{S_t}(\mathbf{W}_{t-1} + \hat{\Delta\mathbf{W}}_t)
     +\; \mathrm{Sharpness}_t(\mathbf{W}_{t-1} + \hat{\Delta\mathbf{W}}_t,\, \rho_t^2)
    \Big] 
    \notag \\
& + (b-a) \left\{ 
    \frac{KL_n^{\text{task}} + KL_n^{\text{hyper}} + \log(m/\delta)}
         {2(\bar{m}_n - 1)}
\right\}^{1/2}
\label{eq:lora-pac-bayes-bound}
\end{align}
where:
% \begin{align}
% KL_n^{\mathrm{task}} &= \sum_{t=1}^n KL(Q_t\|P_t), \notag  \\
% KL_n^{\mathrm{hyper}} &= KL(\mathcal{Q}_1 \| \mathcal{P}_1) +  \\ &+\sum_{t=2}^n \mathbb{E}_{\mathcal{Q}_{1:t-1}}\big[KL(\mathcal{Q}_t \| \mathcal{P}_t(\cdot \| \mathcal{Q}_{1:t-1}))\big], \notag\\
% \bar{m}_n &= n \big/ \left(\sum_{t=1}^n 1/m_t\right), \notag
% \end{align}
\begin{align}
KL_n^{\mathrm{task}} &= \sum_{t=1}^n KL(Q_t\|P_t), \notag  \\
KL_n^{\mathrm{hyper}} &= KL(\mathcal{Q}_1 \| \mathcal{P}_1) +\sum_{t=2}^n \mathbb{E}_{\mathcal{Q}_{1:t-1}}\big[KL(\mathcal{Q}_t \| \mathcal{P}_t(\cdot \| \mathcal{Q}_{1:t-1}))\big], \notag\\
\bar{m}_n &= n \big/ \left(\sum_{t=1}^n 1/m_t\right), \notag
\end{align}
with $m_t$ denoting the sample size for task $t$, and $b-a$ the range of the bounded loss function.
\end{theorem}

\begin{proof}
\textbf{Step 1: Single-task PAC-Bayes Bound (0→1)}

The PAC-Bayesian theory provides a theoretical foundation for analyzing the generalization properties of stochastic model selection, where algorithms probabilistically select models based on a probability distribution over the hypothesis space \(\mathcal{F}\). 


\begin{theorem} \label{the: PAC1999}
\citep{mcallesterPACBayesianStochasticModel2003}
For loss functions $\ell^{'}$ mapping to the range $[0,1]$, the following generalization bound holds for any prior $P$ over parameters with probability $1 - \delta$ over the choice of the training set $S$, for any posterior $Q$ over parameters.
% \begin{equation}\label{eq:PAC2003}
% \begin{aligned}
%     &\forall^{\delta} S,\;\; \forall Q, \\
%     &\quad
%     \mathbb{E}_{f \in Q}\big[L_{\mathcal{D}}^{\ell'}(f)\big] \;\leq\;
%     \mathbb{E}_{f \in Q}\big[\hat{L}^{\ell'}_{S}(f)\big] \\
%     &\qquad
%     +\, \sqrt{
%         \frac{
%             KL(Q\|P) + \log\frac{m}{\delta}
%         }{
%             2(m-1)
%         }
%     }
% \end{aligned}
% \end{equation}

\begin{equation}\label{eq:PAC2003}
\begin{aligned}
    &\forall^{\delta} S,\;\; \forall Q, 
    \mathbb{E}_{f \in Q}\big[L_{\mathcal{D}}^{\ell'}(f)\big] \;\leq\;
    \mathbb{E}_{f \in Q}\big[\hat{L}^{\ell'}_{S}(f)\big] 
    +\, \sqrt{
        \frac{
            KL(Q\|P) + \log\frac{m}{\delta}
        }{
            2(m-1)
        }
    }
\end{aligned}
\end{equation}
where $m=|S|$ is the number of sample in the training set $S$ and $f$ denotes a prediction function sampled from the posterior distribution $Q$ over the hypothesis space $\mathcal{F}$.
\end{theorem}

Theorem \ref{the: PAC1999} is limited to loss functions $\ell^{'}$ mapping to the range $[0,1]$. According to \citep{germainPACBayesianTheoryMeets2016}, by straightforward rescaling, we can extend it to any bounded loss, i.e., $\ell: \mathcal{F} \times \mathcal{X} \times \mathcal{Y} \rightarrow[a, b]$, where $[a, b] \subset \mathbb{R}$. This is done using $\beta:=b-a$ and the rescaled loss function $\ell^{\prime}(f, x, y):=(\ell(f, x, y)-a) /(b-a) \in[0,1]$. After few arithmetic manipulations, we can rewrite Equation (\ref{eq:PAC2003}) as
\begin{equation}
\begin{aligned}
        \mathbb{E}_{f \sim Q}[{L}_\mathcal{D}^{\ell}(f)] \leq \mathbb{E}_{f\sim Q}[\widehat{{L}}_S^{\ell}(f)] + (b - a) \cdot \sqrt{\frac{\mathrm{KL}(Q \| P) + \log \frac{m}{\delta}}{2(m - 1)}} .
\end{aligned}
\end{equation}
% \begin{equation}
% \begin{aligned}
%         &\mathbb{E}_{f \sim Q}[{L}_\mathcal{D}^{\ell}(f)] \leq \mathbb{E}_{f\sim Q}[\widehat{{L}}_S^{\ell}(f)] \\
%         &+ (b - a) \cdot \sqrt{\frac{\mathrm{KL}(Q \| P) + \log \frac{m}{\delta}}{2(m - 1)}} .
% \end{aligned}
% \end{equation}
If the right-hand side of the bound exceeds $b-a$, the bound reduces to
$\mathbb{E}_{f \sim Q}[L_D^{\ell}(f)] \leq b,$
which is always true by definition. So we focus on the nontrivial regime where the complexity term is less than $b-a$, ,which equals to
$\sqrt{\frac{KL(Q\|P)+\log \frac{m}{\delta}}{2(m-1)}} <1. $


Following \citep{pentinaPACBayesianBoundLifelong2014}, we introduce the notion of hyperdistributions. Let $\mathcal{P}$ denote the hyperprior (a distribution over parameter priors $P$) and $\mathcal{Q}$ denote the hyperposterior, the updated distribution over priors after observing data. For each choice of prior $P$, the parameter posterior is denoted $Q(f|P)$, while $P(f)$ is the corresponding parameter prior under $P$.

In PAC-Bayes theory with hyperdistributions, the learned joint distribution over parameters and hyperparameters is given by $Q_\text{joint}(f, P) = Q(f|P)\mathcal{Q}(P)$, and the reference joint distribution before seeing data is $P_\text{joint}(f, P) = P(f)\mathcal{P}(P)$. The KL divergence between these two joint distributions quantifies the adaptation to data in both parameter space and hyperparameter space, and can be decomposed as:
% \begin{align}
% & KL(Q_\text{joint} \| P_\text{joint}) \notag\\
% &\quad =\; 
% \mathbb{E}_{P \sim \mathcal{Q}} \big[\, KL(Q(f|P)\;\|\;P(f))\, \big] 
% + KL(\mathcal{Q} \| \mathcal{P}).
% \end{align}

\begin{align}
& KL(Q_\text{joint} \| P_\text{joint})  =\; 
\mathbb{E}_{P \sim \mathcal{Q}} \big[\, KL(Q(f|P)\;\|\;P(f))\, \big] 
+ KL(\mathcal{Q} \| \mathcal{P}).
\end{align}



Taking the expectation over the hyperposterior $P \sim \mathcal{Q}$, the PAC-Bayes bound for hyperdistributions is:



% \begin{equation}
% \begin{aligned}
%     & \mathbb{E}_{P \sim \mathcal{Q}} \Big[\, \mathbb{E}_{f \sim Q(P)} [L_{\mathcal{D}}(f)] \,\Big] \\
%     &\leq\; \mathbb{E}_{P \sim \mathcal{Q}} \Big[\, \mathbb{E}_{f \sim Q(P)} [\hat{L}_{S}(f)] \,\Big] \\
%     &\quad +\, (b-a)\Bigg\{
%         \frac{1}{2(m-1)} \bigg[
%             \mathbb{E}_{P \sim \mathcal{Q}}\left[ KL(Q(P)\|P) \right] \\
%     &\qquad\qquad
%             +\; KL(\mathcal{Q}\|\mathcal{P}) \\
%     &\qquad\qquad
%             +\; \log \frac{m}{\delta}
%         \bigg]
%     \Bigg\}^{1/2}
% \end{aligned}
% \end{equation}


\begin{equation}
\begin{aligned}
    \mathbb{E}_{P \sim \mathcal{Q}} \Big[\, \mathbb{E}_{f \sim Q(P)} [L_{\mathcal{D}}(f)] \,\Big] 
    & \leq\; \mathbb{E}_{P \sim \mathcal{Q}} \Big[\, \mathbb{E}_{f \sim Q(P)} [\hat{L}_{S}(f)] \,\Big] \\
    &\quad +\, (b-a)\Bigg\{
        \frac{1}{2(m-1)} \bigg[
            \mathbb{E}_{P \sim \mathcal{Q}}\left[ KL(Q(P)\|P) \right] 
            +\; KL(\mathcal{Q}\|\mathcal{P}) 
            +\; \log \frac{m}{\delta}
        \bigg]
    \Bigg\}^{1/2}
\end{aligned}
\end{equation}


Now, let $Q = \mathcal{N}(\hat{\boldsymbol{\theta}}, \rho^2 I)$, where $\hat{\boldsymbol{\theta}}$ is the trained parameter vector. Then,
$$\mathbb{E}_{\boldsymbol{\theta} \sim Q}[\hat{L}_S(\boldsymbol{\theta})] = \mathbb{E}_{\boldsymbol{\epsilon} \sim \mathcal{N}(0, \rho^2 I)}[\hat{L}_S(\hat{\boldsymbol{\theta}} + \boldsymbol{\epsilon})].$$
Refer the \textbf{expected sharpness} definition (\ref{def: Task-Aware Expected Sharpness}):
$$\mathrm{Sh}_S(\hat{\boldsymbol{\theta}}, \rho^2) := \mathbb{E}_{\boldsymbol{\epsilon} \sim \mathcal{N}(0, \rho^2 I)}\left[ \hat{L}_S(\hat{\boldsymbol{\theta}} + \boldsymbol{\epsilon}) \right] - \hat{L}_S(\hat{\boldsymbol{\theta}})$$
Thus,
$$\mathbb{E}_{\boldsymbol{\theta} \sim Q}[\hat{L}_S(\boldsymbol{\theta})] = \hat{L}_S(\hat{\boldsymbol{\theta}}) + \mathrm{Sh}_S(\hat{\boldsymbol{\theta}}, \rho^2)$$

In the LoRA framework, the full-model parameters are $\boldsymbol{W}_1 = \boldsymbol{W}_0 + \Delta\boldsymbol{W}_1$, where $\boldsymbol{W}_0$ denotes the frozen pre-trained backbone and $\Delta\boldsymbol{W}_1$ is the task-specific low-rank adaptation. The full-model distributions factorize as $Q_1^{\text{full}} = \delta_{\boldsymbol{W}_0} \otimes Q_1$ and $P_1^{\text{full}} = \delta_{\boldsymbol{W}_0} \otimes P_1$, leading to:
\begin{equation}
    \text{KL}^{\text{task}}(Q_1^{\text{full}} \| P_1^{\text{full}}) = \text{KL}^{\text{task}}(Q_1 \| P_1),
\end{equation}
where $Q_1$ and $P_1$ are distributions over $\Delta\boldsymbol{W}_1$. 


Although the posterior $Q$ is defined only over $\Delta\boldsymbol{W}_1$, the model prediction and corresponding loss are determined by the entire effective parameter $\boldsymbol{W}_0 + \Delta\boldsymbol{W}_1$. Therefore, when evaluating sharpness, the perturbation is applied to the full parameter vector $\boldsymbol{W}_0 + \Delta\boldsymbol{W}_1$, not just to the low-rank increment $\Delta\boldsymbol{W}_1$. Formally, the expected sharpness in the learned model $\boldsymbol{W}_0 + \hat{\Delta\boldsymbol{W}}_1$ is defined as:
% \begin{equation}
% \begin{aligned}
% &\mathrm{Sh}_S\big(\boldsymbol{W}_0 + \hat{\Delta\mathbf{W}}_1, \rho^2\big) := \\
% & \qquad \mathbb{E}_{\boldsymbol{\epsilon} \sim \mathcal{N}(0, \rho^2 I)} \Big[
%     \widehat{L}_S\big(\boldsymbol{W}_0 + \hat{\Delta\mathbf{W}}_1 + \boldsymbol{\epsilon}\big)
% \Big]  \\
% & \qquad- \widehat{L}_S\big(\boldsymbol{W}_0 + \hat{\Delta\mathbf{W}}_1\big) 
% \notag
% \end{aligned}
% \end{equation}

\begin{equation}
\begin{aligned}
&\mathrm{Sh}_S\big(\boldsymbol{W}_0 + \hat{\Delta\mathbf{W}}_1, \rho^2\big) :=  \mathbb{E}_{\boldsymbol{\epsilon} \sim \mathcal{N}(0, \rho^2 I)} \Big[
    \widehat{L}_S\big(\boldsymbol{W}_0 + \hat{\Delta\mathbf{W}}_1 + \boldsymbol{\epsilon}\big)
\Big]  
- \widehat{L}_S\big(\boldsymbol{W}_0 + \hat{\Delta\mathbf{W}}_1\big) 
\notag
\end{aligned}
\end{equation}
where $\boldsymbol{W}_0 +\hat{\Delta\boldsymbol{W}}_1$ is the optimal LoRA with the base model for task $1$, the perturbation $\boldsymbol{\epsilon}$ is added to the full parameter vector. Thus,
% \begin{equation}
%     \begin{aligned}
%     &\mathbb{E}_{\Delta\mathbf{W} \sim Q}[\widehat{L}_S(\mathbf{W}_0 + \Delta\mathbf{W}_1)] \\
%     &= \mathbb{E}_{\boldsymbol{\epsilon} \sim \mathcal{N}(0, \rho^2 I)}[\widehat{L}_S(\mathbf{W}_0 + \hat{\Delta\mathbf{W}}_1 + \boldsymbol{\epsilon})]  \\
%     &=\widehat{L}_S(\mathbf{W}_0 + \hat{\Delta\mathbf{W}}_1) + \mathrm{Sh}_S(\mathbf{W}_0+\hat{\Delta\mathbf{W}}_1, \rho^2)
% \end{aligned}
% \end{equation}

\begin{equation}
    \begin{aligned}
    &\mathbb{E}_{\Delta\mathbf{W} \sim Q}[\widehat{L}_S(\mathbf{W}_0 + \Delta\mathbf{W}_1)] \\
    &= \mathbb{E}_{\boldsymbol{\epsilon} \sim \mathcal{N}(0, \rho^2 I)}[\widehat{L}_S(\mathbf{W}_0 + \hat{\Delta\mathbf{W}}_1 + \boldsymbol{\epsilon})]  \\
    &=\widehat{L}_S(\mathbf{W}_0 + \hat{\Delta\mathbf{W}}_1) + \mathrm{Sh}_S(\mathbf{W}_0+\hat{\Delta\mathbf{W}}_1, \rho^2)
\end{aligned}
\end{equation}


Finally, the corresponding PAC-Bayes bound, incorporating sharpness under the LoRA setting, is as follows:



% \begin{align}
% & \mathbb{E}_{P \sim \mathcal{Q}}\, \mathbb{E}_{\Delta\mathbf{W} \sim Q(P)} 
%     \big[ L_{\mathcal{D}}(\boldsymbol{W}_0 + \Delta\mathbf{W}_1) \big] \notag \\
% &\leq\; \mathbb{E}_{P \sim \mathcal{Q}} \Big[
%         \widehat{L}_S(\boldsymbol{W}_0 + \hat{\Delta\mathbf{W}}_1)
%         + \mathrm{Sh}_S(\boldsymbol{W}_0 + \hat{\Delta\mathbf{W}}_1,\, \rho^2)
%     \Big] \notag \\
% &\quad +\, (b-a) \Bigg\{
%         \frac{1}{2(m-1)} \bigg[
%             KL^{\text{task}}_1 \notag \\
% &\qquad\qquad
%             +\; KL^{\text{hyper}}_1 \notag \\
% &\qquad\qquad
%             +\; \log\frac{m}{\delta}
%         \bigg]
%     \Bigg\}^{1/2}
% \label{eq:single-task-pac-bayes}
% \end{align}

\begin{align}
 \mathbb{E}_{P \sim \mathcal{Q}}\, \mathbb{E}_{\Delta\mathbf{W} \sim Q(P)} 
    \big[ L_{\mathcal{D}}(\boldsymbol{W}_0 + \Delta\mathbf{W}_1) \big] \notag 
&\leq \mathbb{E}_{P \sim \mathcal{Q}} \Big[
        \widehat{L}_S(\boldsymbol{W}_0 + \hat{\Delta\mathbf{W}}_1)
        + \mathrm{Sh}_S(\boldsymbol{W}_0 + \hat{\Delta\mathbf{W}}_1,\, \rho^2)
    \Big] \notag \\
&\quad +\, (b-a) \Bigg\{
        \frac{1}{2(m-1)} \bigg[
            KL^{\text{task}}_1
            +\; KL^{\text{hyper}}_1
            +\; \log\frac{m}{\delta}
        \bigg]
    \Bigg\}^{1/2}
\label{eq:single-task-pac-bayes}
\end{align}



where $KL^{\text{task}}_1 = \mathbb{E}_{P \sim \mathcal{Q}}\left[ KL(Q(P)|P) \right]$,
$KL^{\text{hyper}}_1 = KL(\mathcal{Q}|\mathcal{P})$,
and $\log(m/\delta)$ is the parameter term.


\textbf{Step 2: Extending the PAC-Bayes Bound from One Task to Two Tasks  (1→2)}
To extend the result to two sequential tasks ($n=2$), we consider the incremetnal LoRA scenario in which the LoRA parameter block for the second task, $\Delta\mathbf{W}_2$, is trained independently, with the backbone fixed at $\mathbf{W}_1 = \mathbf{W}_0 + \hat{\Delta\mathbf{W}}_1$. 


In incremental LoRA, we consider the parameter increments for each task to be the formation of mutually orthogonal subspaces in the parameter space. Formally, let $U_t$ denote the subspace spanned by the LoRA parameter block for task $t$. The orthogonality condition requires that for any pair of tasks $i \neq t$,
$$ \forall\, u \in U_i,\, w \in U_t, \quad \langle u, w \rangle = 0$$
where $\langle \cdot, \cdot \rangle$ denotes the inner product in the parameter space. Therefore, achieving orthogonality between the LoRA subspaces of task $i$ and task $t$ can be expressed as
$${O}_{i,t} = A_i^{\top} A_t = 0,$$
where $A_i$ and $A_t$ are the adaptation matrices for tasks $i$ and $t$, respectively. Under this assumption, the joint KL divergence of product measures satisfies:   \begin{equation}
      KL\left(\bigotimes_{t=1}^n Q_t \, \bigg\| \, \bigotimes_{t=1}^n P_t\right) = \sum_{t=1}^n KL(Q_t \| P_t).
  \end{equation}
where $Q_t$ and $P_t$ denote the posterior and prior distributions over the LoRA parameters for task $t$.



The corresponding joint hyperprior and hyperposterior for the two-task setting are denoted as $\mathcal{P}_{1:2}$ and $\mathcal{Q}_{1:2}$. Here, $\mathcal{P}_{1:2}$ represents the joint hyperprior over both tasks, where the hyperprior for the second task is conditionally dependent on the hyperposterior of the first task, i.e.,
$\mathcal{P}_{1:2} = \mathcal{P}_1 \otimes \mathcal{P}_2(\cdot\,|\,\mathcal{Q}_1).$
Similarly, the joint hyperposterior is given by
$\mathcal{Q}_{1:2} = \mathcal{Q}_1 \otimes \mathcal{Q}_2,$
assuming conditional independence between the task-wise hyperposteriors. Because $\mathcal{P}_2$ is conditionally dependent on $\mathcal{Q}_1$, the joint KL divergence between the hyperdistributions admits a chain-structured decomposition:
\begin{equation}
\begin{aligned}
    KL^{\text{hyper}}_2&=KL(\mathcal{Q}_{1:2}\,\|\,\mathcal{P}_{1:2})  \\
    &= KL(\mathcal{Q}_1\,\|\,\mathcal{P}_1) + \mathbb{E}_{\mathcal{Q}_1}\left[ KL(\mathcal{Q}_2\,\|\,\mathcal{P}_2(\cdot\,|\,\mathcal{Q}_1)) \right].
\end{aligned}
\end{equation}
This structure reflects that the uncertainty or inductive bias for the second task is informed by the knowledge gained from the first task. In practical terms, the conditional hyperprior $\mathcal{P}_2(\cdot,|,\mathcal{Q}_1)$ can take the form of a Gaussian distribution whose parameters (e.g., covariance) are derived from the posterior statistics of the previous task. For example,
$P_2 = \mathcal{N}\big(\mathbf{0},\, \mathbf{F}_1^{-1}\big),$
where $\mathbf{F}_1$ is the Fisher information matrix estimated from the first task, thus encoding the structure and uncertainty learned in the previous stage.


Analogously, for the parameter-level distributions, the KL divergence decomposes as
\begin{align*}
 &KL^{\text{task}}_2 =KL(Q_{1:2}\,\|\,P_{1:2})  \\
 &   = KL(Q_1\,\|\,P_1) + KL(Q_2\,\|\,P_2)
\end{align*}
where $P_2$ is sampled from the conditional hyperprior $\mathcal{P}_2(\cdot|\mathcal{Q}_1)$.



Therefore, the generalization bound for two sequential tasks with recursive, dependent hyperdistributions takes the following form:



% \begin{align}
% & L(\mathcal{Q}_{1:2}) \leq\;  \frac{1}{2} \sum_{t=1}^2 
%     \mathbb{E}_{P_{1:2} \sim \mathcal{Q}_{1:2}} 
%     \Big[\, 
%         \widehat{L}_{S_t}(\mathbf{W}_{t-1} + \hat{\Delta\mathbf{W}}_t) \notag \\
% &\qquad
%     +\; \mathrm{Sh}_t\big(\mathbf{W}_{t-1} + \hat{\Delta\mathbf{W}}_t,\, \rho_t^2\big)
%     \Big] \notag \\
% &\quad
%     +\, (b-a)\Bigg\{ \frac{1}{2(\bar{m}_2 - 1)} \bigg[
%             KL^{\text{task}}_2 \notag \\
% &\qquad\qquad
%             +\; KL^{\text{hyper}}_2 \notag \\
% &\qquad\qquad
%             +\; \log\frac{\bar{m}_2}{\delta}
%         \bigg] \Bigg\}^{1/2}
% \end{align}

\begin{align}
 L(\mathcal{Q}_{1:2}) &\leq \frac{1}{2} \sum_{t=1}^2 
    \mathbb{E}_{P_{1:2} \sim \mathcal{Q}_{1:2}} 
    \Big[\, 
        \widehat{L}_{S_t}(\mathbf{W}_{t-1} + \hat{\Delta\mathbf{W}}_t) 
    + \mathrm{Sh}_t\big(\mathbf{W}_{t-1} + \hat{\Delta\mathbf{W}}_t,\, \rho_t^2\big)
    \Big] \notag \\
&+\, (b-a)\Bigg\{ \frac{1}{2(\bar{m}_2 - 1)} \bigg[
            KL^{\text{task}}_2 
            +\; KL^{\text{hyper}}_2 
            +\; \log\frac{\bar{m}_2}{\delta}
        \bigg] \Bigg\}^{1/2}
\end{align}


where 


% \begin{align*}
% L(\mathcal{Q}_{1:2}) &= \\
% &\mathbb{E}_{P_{1:2} \sim \mathcal{Q}_{1:2}} 
% \left[
%     \frac{1}{2} \mathbb{E}_{\Delta\mathbf{W}_1 \sim Q_1(P_1)} 
%     \left[
%         L_{\mathcal{D}_1}(\mathbf{W}_0 + \Delta\mathbf{W}_1)
%     \right] \right.\\
% &\quad \left.
%     +\, \frac{1}{2} \mathbb{E}_{\Delta\mathbf{W}_2 \sim Q_2(P_2)} 
%     \left[
%         L_{\mathcal{D}_2}(\mathbf{W}_1 + \Delta\mathbf{W}_2)
%     \right]
% \right] \\
% KL^{\text{task}}_2 &=\; 
%     \mathbb{E}_{P_1 \sim \mathcal{Q}_1}\left[ KL(Q_1(P_1)\|P_1) \right] \\
% & \quad +\; \mathbb{E}_{P_2 \sim \mathcal{Q}_2}\left[ KL(Q_2(P_2)\|P_2) \right] \\[1.2ex]
% KL^{\text{hyper}}_2 &=\; 
%     KL(\mathcal{Q}_1\|\mathcal{P}_1) \\
% & \quad +\; \mathbb{E}_{\mathcal{Q}_1}\left[ KL(\mathcal{Q}_2 \| \mathcal{P}_2(\cdot\,|\,\mathcal{Q}_1)) \right] \\
% \bar{m}_2 &=\;  2\, /\, \left(\frac{1}{m_1} + \frac{1}{m_2}\right)
% \end{align*}


\begin{align*}
L(\mathcal{Q}_{1:2}) &=\mathbb{E}_{P_{1:2} \sim \mathcal{Q}_{1:2}} 
\left[
    \frac{1}{2} \mathbb{E}_{\Delta\mathbf{W}_1 \sim Q_1(P_1)} 
    \left[
        L_{\mathcal{D}_1}(\mathbf{W}_0 + \Delta\mathbf{W}_1)
    \right] \right. \left.
    +\, \frac{1}{2} \mathbb{E}_{\Delta\mathbf{W}_2 \sim Q_2(P_2)} 
    \left[
        L_{\mathcal{D}_2}(\mathbf{W}_1 + \Delta\mathbf{W}_2)
    \right]
\right] \\
KL^{\text{task}}_2 &=
    \mathbb{E}_{P_1 \sim \mathcal{Q}_1}\left[ KL(Q_1(P_1)\|P_1) \right] + \mathbb{E}_{P_2 \sim \mathcal{Q}_2}\left[ KL(Q_2(P_2)\|P_2) \right] \\[1.2ex]
KL^{\text{hyper}}_2 &=
    KL(\mathcal{Q}_1\|\mathcal{P}_1)  +\; \mathbb{E}_{\mathcal{Q}_1}\left[ KL(\mathcal{Q}_2 \| \mathcal{P}_2(\cdot\,|\,\mathcal{Q}_1)) \right] \\
\bar{m}_2 &=\;  2\, /\, \left(\frac{1}{m_1} + \frac{1}{m_2}\right)
\end{align*}


 All terms involving empirical risk, sharpness, and complexity are averaged over the joint hyperposterior $\mathcal{Q}_{1:2}$. The chain-structured KL terms reflect that each new task's hyperprior may depend on the hyperposterior of previous tasks, capturing the transfer and adaptation of uncertainty and prior structure in incremental learning.


\textbf{Step 3: Multi-task PAC-Bayes Bound (n→n+1)}

Assume the following holds for $n$ sequential tasks:


\textbf{Parameter Distribution KL:}
    \begin{align}
     KL^{\text{task}}_n=KL(Q_{1:n} \| P_{1:n}) 
    &= \sum_{t=1}^n KL(Q_t \| P_t)
    \end{align}

 
\textbf{Hyperdistribution KL:}
    \begin{align}
    &KL^{\text{hyper}}_n =KL(\mathcal{Q}_{1:n} \| \mathcal{P}_{1:n}) \notag \\
    & \quad  \qquad= KL(\mathcal{Q}_1 \| \mathcal{P}_1)  + \sum_{t=2}^n 
        \mathbb{E}_{\mathcal{Q}_{1:t-1}} \left[ 
            KL\left(\mathcal{Q}_t \| \mathcal{P}_t\big(\cdot\,|\,\mathcal{Q}_{1:t-1}\big)\right) 
        \right]
    \end{align}


    %  \begin{align}
    % &KL^{\text{hyper}}_n =KL(\mathcal{Q}_{1:n} \| \mathcal{P}_{1:n}) \notag \\
    % & \quad  \qquad= KL(\mathcal{Q}_1 \| \mathcal{P}_1) \notag  \\
    %   & \qquad \qquad + \sum_{t=2}^n 
    %     \mathbb{E}_{\mathcal{Q}_{1:t-1}} \left[ 
    %         KL\left(\mathcal{Q}_t \| \mathcal{P}_t\big(\cdot\,|\,\mathcal{Q}_{1:t-1}\big)\right) 
    %     \right]
    % \end{align}

\textbf{Generalization Bound:}

% \begin{align}
% & L(\mathcal{Q}_{1:n}) \leq\;  \frac{1}{n} \sum_{t=1}^n 
%     \mathbb{E}_{P_{1:n} \sim \mathcal{Q}_{1:2}} 
%     \Big[\, 
%         \widehat{L}_{S_t}(\mathbf{W}_{t-1} + \hat{\Delta\mathbf{W}}_t) \notag \\
% &\qquad
%     +\; \mathrm{Sh}_t\big(\mathbf{W}_{t-1} + \hat{\Delta\mathbf{W}}_t,\, \rho_t^2\big)
%     \Big] \notag \\
% &\quad
%     +\, (b-a)\Bigg\{ \frac{1}{n(\bar{m}_n - 1)} \bigg[
%             KL^{\text{task}}_n \notag \\
% &\qquad\qquad
%             +\; KL^{\text{hyper}}_n \notag \\
% &\qquad\qquad
%             +\; \log\frac{\bar{m}_n}{\delta}
%         \bigg] \Bigg\}^{1/2}
% \end{align}


\begin{align}
 L(\mathcal{Q}_{1:n}) &\leq  \frac{1}{n} \sum_{t=1}^n 
    \mathbb{E}_{P_{1:n} \sim \mathcal{Q}_{1:2}} 
    \Big[\, 
        \widehat{L}_{S_t}(\mathbf{W}_{t-1} + \hat{\Delta\mathbf{W}}_t) 
    + \mathrm{Sh}_t\big(\mathbf{W}_{t-1} + \hat{\Delta\mathbf{W}}_t,\, \rho_t^2\big)
    \Big] \notag \\
&+ (b-a)\Bigg\{ \frac{1}{n(\bar{m}_n - 1)} \bigg[
            KL^{\text{task}}_n 
            + KL^{\text{hyper}}_n 
            +\log\frac{\bar{m}_n}{\delta}
        \bigg] \Bigg\}^{1/2}
\end{align}

    
    where


% \begin{align}
% &L(\mathcal{Q}_{1:n}) = \notag\\
% & \frac{1}{n} \sum_{t=1}^n\;
%     \mathbb{E}_{P_{1:n} \sim \mathcal{Q}_{1:n}}\;
%     \mathbb{E}_{\Delta\mathbf{W}_t \sim Q_t(P_t)}
%     \big[
%         L_{\mathcal{D}_t}(\mathbf{W}_{t-1} + \Delta\mathbf{W}_t)
%     \big] \notag \\
% &\bar{m}_n =
%     n \Big/ \bigg( \sum_{t=1}^n \frac{1}{m_t} \bigg) \notag
% \end{align}

\begin{align}
L(\mathcal{Q}_{1:n}) &=  \frac{1}{n} \sum_{t=1}^n\;
    \mathbb{E}_{P_{1:n} \sim \mathcal{Q}_{1:n}}\;
    \mathbb{E}_{\Delta\mathbf{W}_t \sim Q_t(P_t)}
    \big[
        L_{\mathcal{D}_t}(\mathbf{W}_{t-1} + \Delta\mathbf{W}_t)
    \big] \notag \\
\bar{m}_n &=
    n \Big/ \bigg( \sum_{t=1}^n \frac{1}{m_t} \bigg) \notag
\end{align}


\textbf{Inductive Step: From $n$ to $n+1$ Tasks}


\textbf{(a) Extension of Orthogonality.}
The LoRA parameter block $\Delta\mathbf{W}_{n+1}$ is constrained to be orthogonal to all previous blocks:
    \begin{align}
        \forall\, i \leq n, \quad A_i^\top A_{n+1} = 0 \notag
    \end{align}
    Thus,
    \begin{equation}
        KL^{\text{task}}_{n+1}=KL(Q_{1:n+1} \| P_{1:n+1}) = \sum_{t=1}^{n+1} KL(Q_t \| P_t)
    \end{equation}

\textbf{(b) Chain-Structured Hyperprior.}
    The hyperprior for task $n+1$ is conditioned on the cumulative hyperposterior:
    \begin{align}
        \mathcal{P}_{1:n+1} = \mathcal{P}_{1:n} \otimes \mathcal{P}_{n+1}(\cdot\,|\,\mathcal{Q}_{1:n}) \notag
    \end{align}
    So,
    \begin{equation}
        \begin{aligned}
        &KL^{\text{hyper}}_{n+1}
        =KL(\mathcal{Q}_{1:n+1} \| \mathcal{P}_{1:n+1}) \\
        &= KL(\mathcal{Q}_{1:n} \| \mathcal{P}_{1:n}) +\mathbb{E}_{\mathcal{Q}_{1:n}}\left[ KL\big(\mathcal{Q}_{n+1} \| \mathcal{P}_{n+1}(\cdot\,|\,\mathcal{Q}_{1:n})\big) \right] 
        \\
        &= KL(\mathcal{Q}_1 \| \mathcal{P}_1) +
         \sum_{t=2}^{n+1} \mathbb{E}_{\mathcal{Q}_{1:t-1}}\left[ KL\left(\mathcal{Q}_t \| \mathcal{P}_t(\cdot\,|\,\mathcal{Q}_{1:t-1})\right) \right]
    \end{aligned}
    \end{equation}


\textbf{(c) Generalization Bound.}
    The empirical risk and sharpness terms for the new task are averaged in:

% \begin{align}
% & L(\mathcal{Q}_{1:n+1})  \\
% &\leq\; \frac{1}{n+1} \sum_{t=1}^{n+1}
%     \mathbb{E}_{P_{1:n+1} \sim \mathcal{Q}_{1:n+1}}\big[
%         \widehat{L}_{S_t}(\mathbf{W}_{t-1} + \hat{\Delta\mathbf{W}}_t) \notag \\
% &\qquad
%     +\; \mathrm{Sh}_t(\mathbf{W}_{t-1} + \hat{\Delta\mathbf{W}}_t,\, \rho_t^2)
%     \big] \notag \\
% &\quad
%     +\, (b-a) \Bigg\{
%         \frac{1}{2(\bar{m}_{n+1} - 1)} \bigg[
%             KL_{n+1}^{\text{task}} \notag \\
% &\qquad\qquad
%             +\; KL_{n+1}^{\text{hyper}} \notag \\
% &\qquad\qquad
%             +\; \log\frac{\bar{m}_{n+1}}{\delta}
%         \bigg]
%     \Bigg\}^{1/2}
% \end{align}

\begin{align}
 L(\mathcal{Q}_{1:n+1}) &\leq \frac{1}{n+1} \sum_{t=1}^{n+1}
    \mathbb{E}_{P_{1:n+1} \sim \mathcal{Q}_{1:n+1}}\big[
        \widehat{L}_{S_t}(\mathbf{W}_{t-1} + \hat{\Delta\mathbf{W}}_t) 
    + \mathrm{Sh}_t(\mathbf{W}_{t-1} + \hat{\Delta\mathbf{W}}_t,\, \rho_t^2)
    \big] \notag \\
&+ (b-a) \Bigg\{
        \frac{1}{2(\bar{m}_{n+1} - 1)} \bigg[
            KL_{n+1}^{\text{task}} \notag 
            +\ KL_{n+1}^{\text{hyper}} 
            + \log\frac{\bar{m}_{n+1}}{\delta}
        \bigg]
    \Bigg\}^{1/2}
\end{align}
where

 % \begin{equation}
 %        \begin{aligned}
 %        KL_{n+1}^{\text{task}} &= \sum_{t=1}^{n+1} KL(Q_t\|P_t)\\
 %        KL_{n+1}^{\text{hyper}} = & KL(\mathcal{Q}_1\|\mathcal{P}_1)  \\
 %        &+\sum_{t=2}^{n+1} \mathbb{E}_{\mathcal{Q}_{1:t-1}}
 %        \left[ KL(\mathcal{Q}_t \| \mathcal{P}_t(\cdot\,|\,\mathcal{Q}_{1:t-1})) \right] \\
 %        \bar{m}_{n+1} &= (n+1)/\big(\sum_{t=1}^{n+1} 1/m_t\big) \notag
 %    \end{aligned}
 %    \end{equation}
    
    \begin{equation}
        \begin{aligned}
        KL_{n+1}^{\text{task}} &= \sum_{t=1}^{n+1} KL(Q_t\|P_t)\\
        KL_{n+1}^{\text{hyper}} = & KL(\mathcal{Q}_1\|\mathcal{P}_1)  +\sum_{t=2}^{n+1} \mathbb{E}_{\mathcal{Q}_{1:t-1}}
        \left[ KL(\mathcal{Q}_t \| \mathcal{P}_t(\cdot\,|\,\mathcal{Q}_{1:t-1})) \right] \\
        \bar{m}_{n+1} &= (n+1)/\big(\sum_{t=1}^{n+1} 1/m_t\big) \notag
    \end{aligned}
    \end{equation}

\textbf{Step4: Conclusion General Bound for $n$ Tasks}

By induction, the PAC-Bayesian generalization bound for incremental multi-task LoRA with orthogonal blocks and recursively dependent hyperpriors is:
% \begin{align}
%         L(\mathcal{Q}_{1:n}) \leq & \frac{1}{n} \sum_{t=1}^{n}
%         \mathbb{E}_{P_{1:n} \sim \mathcal{Q}_{1:n}}\big[
%             \widehat{L}_{S_t}(\mathbf{W}_{t-1} + \hat{\Delta\mathbf{W}}_t) \notag\\
%         & + \mathrm{Sh}_t(\mathbf{W}_{t-1} + \hat{\Delta\mathbf{W}}_t,\, \rho_t^2)
%         \big] \notag\\
%         & + \frac{1 }{n(\bar{m}_n - 1)}\sqrt{
%             \frac{KL_{n}^{\text{task}} + KL_{n}^{\text{hyper}} +\log(\bar{m}_{n}/\delta)}{2(\bar{m}_{n} - 1)}
%         }
%     \end{align}

    \begin{align}
        L(\mathcal{Q}_{1:n}) \leq & \frac{1}{n} \sum_{t=1}^{n}
        \mathbb{E}_{P_{1:n} \sim \mathcal{Q}_{1:n}}\big[
            \widehat{L}_{S_t}(\mathbf{W}_{t-1} + \hat{\Delta\mathbf{W}}_t)  + \mathrm{Sh}_t(\mathbf{W}_{t-1} + \hat{\Delta\mathbf{W}}_t,\, \rho_t^2)
        \big] \notag\\
        & + \frac{1 }{n(\bar{m}_n - 1)}\sqrt{
            \frac{KL_{n}^{\text{task}} + KL_{n}^{\text{hyper}} +\log(\bar{m}_{n}/\delta)}{2(\bar{m}_{n} - 1)}
        }
    \end{align}


 This result demonstrates that, under the incremental LoRA framework, orthogonalization of parameter spaces and sequential modeling of hyperpriors allow the single-task PAC-Bayes theory to be strictly extended to multi-task continual learning. This structure ensures that model complexity, empirical risk, and sharpness are balanced while enabling transfer and retention of knowledge across tasks.




\end{proof}



\subsection{C: Implementation Details}


All experiments are implemented in PyTorch with two NVIDIA L40 GPUs, each with 46GB of memory  and bfloat16 precision. We adopt the T5-large model for short sequence tasks and long sequence tasks. For TRACE, we use LLaMA-2-7B-chat. We use a LoRA rank of 8, constant learning rate schedules without warm-up, and disable intermediate checkpoint saving.

\paragraph{Short Sequence Benchmark.} We use a per-device batch size of 16 (on 2 GPUs), maximum input length 512, output length 50, and train for 5 epochs. Each task uses up to 3000 training samples. The learning rate is set to $1\times10^{-4}$ with gradient accumulation steps of 1.

\paragraph{Long Sequence Benchmark.} Settings are the same as short sequence, except each task uses 1000 samples.

\paragraph{TRACE Benchmark.} We use LLaMA-2-7B-chat and evaluate on eight diverse tasks from the TRACE benchmark: {C-STANCE}, {FOMC}, {MeetingBank}, {Py150}, {ScienceQA}, {NumGLUE-cm}, {NumGLUE-ds}, and {20Minuten}. For each task, we use a per-device training batch size of 1 and evaluation batch size of 16.  he learning rate is set to $1\times10^{-4}$ with no weight decay, and the number of training epochs varies by task: 5 (C-STANCE), 3 (FOMC), 7 (MeetingBank), 5 (Py150), 3 (ScienceQA), 5 (NumGLUE-cm), 5 (NumGLUE-ds), and 7 (20Minuten). Both the learning rate and task-specific epoch settings are adopted from the original TRACE benchmark\cite{wang2023tracecomprehensivebenchmarkcontinual}.



% We follow the original TRACE Benchmark setup . Each task uses a training batch size of 1 and evaluation batch size of 16. 

% Prompt and answer lengths are 1024 and 512, respectively. Learning rates and epoch counts are adopted from the original TRACE configuration.

\paragraph{Flatness-Aware Methods.}  
We evaluate several perturbation-based flatness strategies:

\begin{itemize}
    \item \textbf{SAM}: Sharpness-Aware Minimization is applied with AdamW and a perturbation radius of 0.002.
    \item \textbf{LoRA-Gaussian}: Gaussian noise with zero mean and standard deviation 0.002 is injected into the LoRA adapter of the current task. Learning rate is set to $1\times10^{-4}$.
    \item \textbf{LoRA-Fisher}: Perturbations are scaled by the Fisher information of parameters. The Fisher matrix is estimated online with exponential decay (decay factor 0.99). Standard deviation is 0.002.
    \item \textbf{Full-Gaussian}: Isotropic Gaussian noise with zero mean and a standard deviation of 0.002 is independently applied to all trainable parameters in the model, including the newly introduced LoRA adapter, the backbone, and previously learned adapters. The noise is injected uniformly across all components regardless of parameter type.

    \item \textbf{Full-Fisher}: Fisher-weighted Gaussian noise is applied selectively: the newly added LoRA adapter receives adaptive noise scaled by the parameter-wise Fisher information (as in LoRA-Fisher), while all remaining components—including the frozen backbone and older adapters—receive standard Gaussian noise. This hybrid strategy retains the flatness-aware benefits in active components while introducing stochastic regularization to the full model space.

    \item \textbf{OLoRA}: Applies orthogonal regularization with $\lambda_1 = 0.5$ and $\lambda_2 = 0$. All flatness-promoting strategies are compatible and used similarly as in IncLoRA.
\end{itemize}






\subsection{D: Datasets}
Table \ref{table:data} shows details of the 15 datasets we used for our short and long experiments, along with their evaluation metrics. Overall, we used datasets from CL benchmark (Zhang et al., 2015), GLUE (Wang et al., 2018) and SuperGLUE (Wang et al., 2019) benchmarks, and added IMDB movie reviews dataset, following (Razdaibiedina et al., 2023).

\begin{table}[htbp]
\centering
\begin{tabular}{lllll}
\toprule
\textbf{Dataset name} & \textbf{Category} & \textbf{Task} & \textbf{Domain} & \textbf{Metric} \\
\midrule
1. Yelp    & CL Benchmark & sentiment analysis      & Yelp reviews     & accuracy \\
2. Amazon  & CL Benchmark & sentiment analysis      & Amazon reviews   & accuracy \\
3. DBpedia & CL Benchmark & topic classification    & Wikipedia        & accuracy \\
4. Yahoo   & CL Benchmark & topic classification    & Yahoo QA       & accuracy \\
5. AG News & CL Benchmark & topic classification    & news             & accuracy \\
6. MNLI    & GLUE         & NLI                    & various          & accuracy \\
7. QQP     & GLUE         & paragraph detection     & Quora            & accuracy \\
8. RTE     & GLUE         & NLI                    & news, Wikipedia  & accuracy \\
9. SST-2   & GLUE         & sentiment analysis      & movie reviews    & accuracy \\
10. WiC    & SuperGLUE    & word sense disambiguation & lexical databases & accuracy \\
11. CB     & SuperGLUE    & NLI                    & various          & accuracy \\
12. COPA   & SuperGLUE    & QA                     & blogs, encyclopedia & accuracy \\
13. BoolQA & SuperGLUE    & boolean QA             & Wikipedia        & accuracy \\
14. MultiRC& SuperGLUE    & QA                     & various          & accuracy \\
15. IMDB   & SuperGLUE    & sentiment analysis      & movie reviews    & accuracy \\
\bottomrule
\end{tabular}
\label{table:data}
\caption{The details of 15 datasets used in our CL experiments. NLI denotes natural language inference, QA denotes questions and answers task. First five tasks correspond to the standard CL benchmark, all other tasks are used in long-sequence experiments.}
\end{table}



% \newpage
% \input{ReproducibilityChecklist_clean}

\end{document}
